\documentclass[11pt,a4paper]{article}
\usepackage{fontspec}
\setmainfont{Noto Serif}
\setsansfont{Noto Sans}
\setmonofont{Noto Sans Mono}[Scale=0.86]
\usepackage[english]{babel}
\usepackage{amsmath,amssymb,amsthm,mathtools}
\usepackage{geometry,booktabs,longtable,tabularx,array}
\usepackage{xcolor,hyperref,xurl,fancyhdr,enumitem}
\geometry{margin=2.2cm,headheight=15pt}
\definecolor{ink}{HTML}{183449}
\definecolor{teal}{HTML}{187C80}
\hypersetup{colorlinks=true,linkcolor=ink,citecolor=teal,urlcolor=teal,
pdftitle={ECDR Mathematics: Concepts, Foundations, and Proofs},pdfauthor={ECDR Research}}
\urlstyle{same}
\makeatletter
\renewcommand{\@pnumwidth}{2.5em}
\renewcommand{\@tocrmarg}{3.5em}
\renewcommand{\l@section}[2]{\@dottedtocline{1}{0em}{3em}{\bfseries#1}{#2}}
\makeatother
\setlength{\parindent}{0pt}
\setlength{\parskip}{5pt}
\setlength{\emergencystretch}{3em}
\setlength{\tabcolsep}{4pt}
\renewcommand{\arraystretch}{1.12}
\setlist{nosep,leftmargin=1.5em}
\newcolumntype{Y}{>{\raggedright\arraybackslash}X}
\newcolumntype{P}[1]{>{\raggedright\arraybackslash}p{#1}}
\newcommand{\R}{\mathbb R}
\newcommand{\norm}[1]{\left\lVert#1\right\rVert}
\newcommand{\abs}[1]{\left\lvert#1\right\rvert}
\newcommand{\code}[1]{{\small\ttfamily\nolinkurl{#1}}}
\newcommand{\relu}{\operatorname{ReLU}}
\newcommand{\diag}{\operatorname{diag}}
\newcommand{\E}{\mathbb E}
\newcommand{\conceptid}{Front matter}
\newcommand{\concept}[3]{\clearpage\renewcommand{\conceptid}{#1}%
\section{#2}\label{concept:#1}%
\textcolor{teal}{\textbf{Concept #1.} Prerequisites and connections: #3.}\par
\typeout{ECDR-CONCEPT-START #1 PAGE=\thepage}}
\pagestyle{fancy}\fancyhf{}
\fancyhead[L]{\small ECDR Mathematics | English Concept Atlas}
\fancyhead[R]{\small\conceptid}
\fancyfoot[L]{\small 28 October 2026}
\fancyfoot[R]{\thepage}
\begin{document}
\hypersetup{pageanchor=false}
\begin{titlepage}
\vspace*{1cm}
{\Large\color{teal} FOUNDATIONS, DEFINITIONS, AND CONDITIONAL PROOFS\par}
\vspace{0.7cm}
{\huge\bfseries\color{ink} ECDR Mathematics:\\Concepts, Foundations,\\and Proofs\par}
\vspace{0.8cm}
{\Large One concept per page, from foundational axioms\\to certified decision reuse\par}
\vspace{1cm}
\begin{tabularx}{\linewidth}{@{}P{3.2cm}Y@{}}\toprule
Language & English throughout; previous Vietnamese documents remain available separately.\\
Reading unit & 104 concepts: 46 foundations, 30 data/model bridges, 28 certification and systems concepts.\\
Explanatory depth & Definition, every local symbol, derivation or example, motivation, assumptions and connections.\\
Meaning of span & Linear span, annotation temporal span, and receptive-field span are explained separately.\\
Reference model & Frozen SmallCNN1D K=12,283,692 parameters, padded Conv/BN/Pool/Flatten--MLP. Valid/ReLU/GAP concepts retain their stated historical assumptions.\\
Evidence boundary & Mathematical derivations and declared examples; no new HAPT inference, formal executable verification or runtime benchmark.\\\bottomrule
\end{tabularx}
\vfill
This atlas expands the previous mathematical explanation. It traces the
assumptions actually used by ECDR back to the real-number, vector-space and
probability frameworks. It does not claim to prove axioms from nothing,
formalize all set theory, or establish a new foundation for mathematics.
\end{titlepage}
\hypersetup{pageanchor=true}\pagenumbering{arabic}
\section*{How to read a concept page}
\addcontentsline{toc}{section}{How to read a concept page}
Start with an unfamiliar term and follow its prerequisite IDs. Each concept
starts on a separate page. Front matter, coverage tables and references are
not concept pages. Symbols are explained locally so a page can be read
without reconstructing definitions from a distant chapter.

Four kinds of statements have different logical roles. An \textbf{axiom}
is accepted in the chosen mathematical framework. A \textbf{definition}
assigns meaning to notation. A \textbf{theorem} follows from stated
hypotheses. A \textbf{modeling assumption} chooses how a physical or
computational system is represented. A 50 Hz sampling model is not a
field axiom; a positive classifier margin is not an axiom of probability.

\textbf{The foundation chain.} Real-number algebra and order support scalar
inequalities. Vector-space structure supports linear combinations and
span. Inner products and norms support distance, projection and duality.
Induction, convexity, spectral algebra, calculus, optimization and
complexity are unpacked before their applications.
Continuity and compactness support existence statements. Probability
axioms support coverage statements. Sound operator bounds and a fixed
reference connect these tools to a universal decision property.

\textbf{Canonical pipeline notation.} A raw vector is $x\in\R^d$;
normalized input is $z=\mathcal P(x)$. The neural network is $f_\theta$,
the full logit map is $F=f_\theta\circ\mathcal P$, and the class is
$h(x)=\arg\max_jF_j(x)$ with a fixed tie rule. Raw pairwise margins are
$g_{cj}(x)=F_c(x)-F_j(x)$. In normalized coordinates,
$\gamma_{cj}(z)=f_{\theta,c}(z)-f_{\theta,j}(z)$, so
$g_{cj}=\gamma_{cj}\circ\mathcal P$. Tensor versions $X,Z$ retain
channel/time axes. Dimensions, index ranges and units are restated on the
relevant concept pages.

\textbf{Three uses of span.} F12 concerns the set of all linear combinations
of vectors. B02 concerns an annotated interval of time indices. B21
concerns the extreme input indices reachable by a feature. Shared English
vocabulary does not make these objects mathematically interchangeable.

\textbf{Limits.} Real-arithmetic proofs are conditional on their stated
reference and domain. Existing empirical artifacts are not new experiments.
No value of coverage, accuracy or speedup is inferred merely from a theorem.
See C20--C28 for executable soundness, lifecycle, cost and evidence boundaries.
\clearpage
\begingroup\small\setlength{\parskip}{0pt}\tableofcontents\endgroup

\concept{F01}{Propositions and quantifiers}{None}
\paragraph{Foundation and definition.}
A proposition is a statement with a truth value after its terms have been
specified. A predicate is a statement that still depends on a variable.
The universal quantifier $\forall$ means ``for every''; the existential
quantifier $\exists$ means ``there is at least one.'' These are logical
definitions, rather than empirical claims about sensor data.
\paragraph{Symbols and meaning.}
Let $\mathcal S$ be a set of permitted input windows, let $x$ denote one
window in that set, let $h:\mathcal S\to\{1,\ldots,K\}$ be a deterministic
classifier, let the integer $K\ge2$ count classes, and let $c$ be one fixed
class. The predicate $P(x)$ is the statement $h(x)=c$. A window can be
represented by $768$ real coordinates, but the logical statement does not
depend on that particular dimension.
\paragraph{Derivation or example.}
The statements
\[
 \forall x\in\mathcal S\;P(x),\qquad
 \exists x\in\mathcal S\;P(x)
\]
have very different strengths. For $\mathcal S=\{u,v\}$, with
$h(u)=c$ and $h(v)\ne c$, the existential statement is true and the
universal statement is false. Negation changes the quantifier:
\[
 \neg(\forall x\in\mathcal S\;P(x))
 \ \Longleftrightarrow\ 
 \exists x\in\mathcal S\;\neg P(x).
\]
Thus a single counterexample defeats a universal claim. Testing a finite
selection from an infinite region does not establish that claim.
\paragraph{Why ECDR needs it.}
A regional certificate asserts that every input admitted by its domain
has the certified decision. Observing that several neighboring windows
share a prediction supplies examples, not the missing universal argument.
\paragraph{Conditions and connections.}
Quantifier order matters: ``for every input there exists a safe radius''
does not specify one radius that works for every input. Universal
statements over the empty set are vacuously true; a useful reuse domain
must also be nonempty. \hyperref[concept:F02]{F02} explains the logical role of a certificate.
\par\label{endconcept:F01}
\typeout{ECDR-CONCEPT-END F01 PAGE=\thepage}

\concept{F02}{Implication, necessity, sufficiency, and equivalence}{\hyperref[concept:F01]{F01}}
\paragraph{Foundation and definition.}
An implication $A\Rightarrow B$ says that whenever proposition $A$ holds,
proposition $B$ holds. Then $A$ is sufficient for $B$, and $B$ is necessary
for $A$. Equivalence $A\Leftrightarrow B$ requires both implications.
These relationships describe logical strength; they do not assert that
either proposition is actually true in a particular run.
\paragraph{Symbols and meaning.}
Let $A$ mean ``a sound verifier returns a strictly positive lower bound,''
and let $B$ mean ``the chosen class wins throughout the specified input
region.'' The word sound means that the verifier's claimed inequality is
mathematically valid under its stated model and arithmetic assumptions.
Neither letter is a numerical variable or a dataset label.
\paragraph{Derivation or example.}
For a scalar variable $x\in\R$, take the margin $m(x)=2-x$ on the
interval $\mathcal S=[0,1]$. Its exact minimum is $1$. A valid but loose
lower bound can be $-3$, because $-3\le m(x)$ for every permitted $x$.
Therefore a positive returned bound proves safety, but a nonpositive
returned bound need not prove unsafety:
\[
 \text{positive sound lower bound}\Rightarrow\text{safety},
 \qquad
 \text{no positive bound}\not\Rightarrow\text{unsafety}.
\]
The contrapositive $\neg B\Rightarrow\neg A$ follows from
$A\Rightarrow B$. The converse $B\Rightarrow A$ does not follow.
\paragraph{Why ECDR needs it.}
An unsuccessful certificate attempt should trigger the defined fallback.
It must not be reported as a discovered decision error. Conversely, a
successful attempt is useful only if its soundness conditions hold.
\paragraph{Conditions and connections.}
A conditional theorem proves its conclusion from hypotheses. An
unconditional claim has no remaining hypotheses within its declared
framework; it still uses the framework's axioms and definitions. Saying
``unconditional guarantee'' cannot remove model identity, input-domain,
or numerical assumptions. \hyperref[concept:F28]{F28} gives a concrete conditional theorem.
\par\label{endconcept:F02}
\typeout{ECDR-CONCEPT-END F02 PAGE=\thepage}

\concept{F03}{Sets, membership, and inclusion}{\hyperref[concept:F01]{F01}}
\paragraph{Foundation and definition.}
A set specifies which objects belong to a collection. Membership
$x\in\mathcal S$ is a yes-or-no relation. Inclusion
$\mathcal A\subseteq\mathcal B$ means every member of $\mathcal A$ is
also a member of $\mathcal B$. A set is different from its size, its
bounding box, and a list of sampled members.
\paragraph{Symbols and meaning.}
Here $x$ is an input vector; $\mathcal S$, $\mathcal A$, and
$\mathcal B$ are sets of input vectors in the same coordinate system.
The empty set $\varnothing$ has no members. The union
$\mathcal A\cup\mathcal B$ admits a member of either set; the
intersection $\mathcal A\cap\mathcal B$ admits only shared members.
The difference $\mathcal B\setminus\mathcal A$ contains members of
$\mathcal B$ outside $\mathcal A$.
\paragraph{Derivation or example.}
Let $\mathcal A=[-1,1]$ and $\mathcal B=[-2,2]$, subsets of $\R$.
Then $0.5\in\mathcal A\subseteq\mathcal B$, whereas
$1.5\in\mathcal B\setminus\mathcal A$. If a classifier $h$ satisfies
$h(x)=c$ for all $x\in\mathcal B$, it also satisfies this identity on
$\mathcal A$. Proof: choose any member $x$ of $\mathcal A$; inclusion
puts the same $x$ in $\mathcal B$, where the identity is known.
The reverse transfer is invalid without more information.
\paragraph{Why ECDR needs it.}
Membership is the online gate that connects a particular query to a
previously proved regional statement. A conservative inner gate can
reduce reuse coverage while preserving safety. Enlarging the gate beyond
the certified set can invalidate safety.
\paragraph{Conditions and connections.}
An inequality can define a set without proving that observed data usually
falls inside it. All comparisons require identical normalization and
indexing. \hyperref[concept:F24]{F24} constructs boxes; \hyperref[concept:F34]{F34} explains why individually chosen
coordinate quantiles do not automatically determine joint coverage.
\par\label{endconcept:F03}
\typeout{ECDR-CONCEPT-END F03 PAGE=\thepage}

\concept{F04}{Functions, domains, images, and preimages}{\hyperref[concept:F03]{F03}}
\paragraph{Foundation and definition.}
A function assigns exactly one output to every input in its domain.
Writing $f:\mathcal X\to\mathcal Y$ names the domain $\mathcal X$
and codomain $\mathcal Y$. The image is the set of outputs actually
attained. A preimage collects inputs whose outputs belong to a specified
output set; it does not require an inverse function.
\paragraph{Symbols and meaning.}
The variable $x\in\mathcal X$ is one input and $f(x)\in\mathcal Y$
is its output. For subsets $\mathcal A\subseteq\mathcal X$ and
$\mathcal B\subseteq\mathcal Y$, define
\[
 f(\mathcal A)=\{f(x):x\in\mathcal A\},\qquad
 f^{-1}(\mathcal B)=\{x\in\mathcal X:f(x)\in\mathcal B\}.
\]
The superscript $-1$ here denotes a preimage operation, not division.
\paragraph{Derivation or example.}
For $f:\R\to\R$ given by $f(x)=x^2$, the codomain is $\R$ but
the image is $[0,\infty)$. The preimage of $[1,4]$ is
$[-2,-1]\cup[1,2]$. No single-valued inverse on the full image can
recover whether the original input was positive or negative.
For a classifier $h:\R^{d}\to\{1,\ldots,K\}$, the decision region
of class $c$ is $h^{-1}(\{c\})$. Here integers $d\ge1$ and $K\ge2$
count input coordinates and classes.
\paragraph{Why ECDR needs it.}
Certification seeks a practical input set contained in one decision
region. Feature compression can map different inputs to the same feature
vector, so invertibility is unnecessary for classification or reuse.
\paragraph{Conditions and connections.}
The actual reference must be deterministic under its inference state;
random layers require a separately stated treatment. \hyperref[concept:F17]{F17} explains how
preprocessing, a CNN, and a decision rule form a composed function.
\par\label{endconcept:F04}
\typeout{ECDR-CONCEPT-END F04 PAGE=\thepage}

\concept{F05}{Real-number field axioms and algebra}{\hyperref[concept:F02]{F02}}
\paragraph{Foundation and definition.}
The field axioms are assumptions about addition and multiplication on
$\R$. Each operation is closed. Both are associative and commutative.
There are identities $0$ and $1$ with $0\ne1$. Every element has an
additive inverse, every nonzero element has a multiplicative inverse, and
multiplication distributes over addition. These axioms are not proved
by collecting sensor measurements.
\paragraph{Symbols and meaning.}
For arbitrary real numbers $a,b,c$, associativity includes
$(a+b)+c=a+(b+c)$; distributivity is $a(b+c)=ab+ac$. The symbols
$-a$ and $a^{-1}$ mean the respective inverses, with $a^{-1}$ defined
only for $a\ne0$. Define $a-b=a+(-b)$ and $a/b=a b^{-1}$ for $b\ne0$.
\paragraph{Derivation or example.}
Cancellation follows from the axioms rather than being another empirical
observation. If $a+c=b+c$, adding $-c$ to both sides and reassociating
gives $a=b$. Also,
\[
 a\cdot0=a(0+0)=a\cdot0+a\cdot0
 \quad\Rightarrow\quad a\cdot0=0.
\]
This proves why division by zero cannot be defined as a multiplicative
inverse: an inverse would imply $1=0$. For a logit pair $u,v\in\R$,
the margin $u-v$ is simply addition with the additive inverse of $v$.
\paragraph{Why ECDR needs it.}
Rearranging margin expressions, canceling shared bias terms, and
distributing sums over a linear head all rely on these algebraic laws.
Stating the foundation makes clear which manipulations are legitimate.
\paragraph{Conditions and connections.}
Floating-point arithmetic only approximates this field: addition can
fail associativity after rounding. Real-number proofs therefore need an
additional error argument before becoming machine-level certificates.
Order and completeness require further axioms, developed in \hyperref[concept:F06]{F06}--\hyperref[concept:F07]{F07}.
\par\label{endconcept:F05}
\typeout{ECDR-CONCEPT-END F05 PAGE=\thepage}

\concept{F06}{Ordered-field axioms and inequalities}{\hyperref[concept:F05]{F05}}
\paragraph{Foundation and definition.}
An ordered field adds a total order compatible with field operations.
For any two real numbers exactly one of $a<b$, $a=b$, or $b<a$ holds.
Order is transitive; adding the same number preserves order; and the
product of positive numbers is positive. These are axioms of the real
number system, not assumptions about the classifier.
\paragraph{Symbols and meaning.}
The symbols $a,b,c\in\R$ denote arbitrary scalars. Write $a\le b$
when $a<b$ or $a=b$. A positive scalar satisfies $c>0$; a negative scalar
satisfies $c<0$. The notation $[a,b]$ later denotes all scalars between
the endpoints when $a\le b$.
\paragraph{Derivation or example.}
If $a\le b$ and $c>0$, then $ac\le bc$: from $b-a\ge0$, product
compatibility gives $c(b-a)\ge0$, and distributivity finishes the proof.
For $c<0$, multiplying reverses the direction. For example,
\[
 -2\le3\quad\Rightarrow\quad 4\le-6
 \quad\text{is false},\qquad -2\le3\Rightarrow4\ge-6.
\]
If a margin satisfies $m\ge\ell$ and a proved bound gives $\ell>0$,
transitivity yields $m>0$. If only $\ell\le0$ is known, neither the
sign of $m$ nor a decision error follows.
\paragraph{Why ECDR needs it.}
The direction of a bound determines its usefulness. A lower margin bound
supports a safety claim; an upper bound with the same numerical value
generally does not. Negative coefficients also determine which interval
endpoint minimizes an affine expression.
\paragraph{Conditions and connections.}
Order is not a floating-point tolerance policy. Replacing $m>0$ by
$m\ge-10^{-6}$ changes the mathematical acceptance condition and needs
its own justification. \hyperref[concept:F25]{F25} distinguishes bounding values from attained
minima; \hyperref[concept:F28]{F28} uses strict positivity.
\par\label{endconcept:F06}
\typeout{ECDR-CONCEPT-END F06 PAGE=\thepage}

\concept{F07}{Completeness and the least-upper-bound axiom}{\hyperref[concept:F03]{F03}, \hyperref[concept:F06]{F06}}
\paragraph{Foundation and definition.}
Completeness distinguishes the real numbers from an arbitrary ordered
field. Its least-upper-bound formulation states: every nonempty subset
of $\R$ that is bounded above has a least upper bound in $\R$. This is
an axiom characterizing the real-number setting. It does not say that
every set contains its upper bound.
\paragraph{Symbols and meaning.}
Let $\mathcal A\subseteq\R$ be nonempty. A number $u\in\R$ is an
upper bound when $a\le u$ for every $a\in\mathcal A$. Its supremum
$s=\sup\mathcal A$ is an upper bound satisfying $s\le u$ for every
upper bound $u$. Bounded above means that at least one such finite $u$
exists.
\paragraph{Derivation or example.}
For $\mathcal A=(0,1)$, the supremum is $1$, although
$1\notin\mathcal A$. Every number below $1$ fails to be an upper bound,
because there is a larger member still below $1$.
Completeness also gives an infimum for a nonempty set bounded below:
\[
 \inf\mathcal A=-\sup(-\mathcal A),\qquad
 -\mathcal A=\{-a:a\in\mathcal A\}.
\]
The ordered rational field lacks this completeness: the positive
rational numbers with square below $2$ have no rational least upper
bound. A real square root supplies the missing bound.
\paragraph{Why ECDR needs it.}
Worst-case optimization uses infima over infinitely many possible inputs.
Completeness supports the real-valued limits and extrema theory behind
certification, even when a practical verifier returns only an enclosure.
\paragraph{Conditions and connections.}
Completeness does not supply an algorithm or guarantee finite computation.
An unbounded margin set may have extended-real infimum $-\infty$ instead
of a finite real bound. \hyperref[concept:F25]{F25} explains infimum versus minimum, and
\hyperref[concept:F27]{F27}--\hyperref[concept:F28]{F28} add hypotheses that make an extremum attainable.
\par\label{endconcept:F07}
\typeout{ECDR-CONCEPT-END F07 PAGE=\thepage}

\concept{F08}{Integer indexing, floor, and ceiling}{\hyperref[concept:F05]{F05}, \hyperref[concept:F06]{F06}}
\paragraph{Foundation and definition.}
An index identifies a discrete position. For a real scalar $u$,
$\lfloor u\rfloor$ is the greatest integer not exceeding $u$, and
$\lceil u\rceil$ is the least integer not smaller than $u$. These
definitions convert continuous counts or thresholds into valid discrete
positions; choosing zero-based indexing is a convention, not a theorem.
\paragraph{Symbols and meaning.}
Let the integers $T\ge1$, $W\ge1$, and $s\ge1$ count recording samples,
window samples, and stride samples. Let $i$ be a zero-based window index.
Window $i$ starts at sample $is$ and contains indices
$is,is+1,\ldots,is+W-1$. A sampling frequency $f_s>0$, measured in
samples per second, converts index differences to seconds.
\paragraph{Derivation or example.}
A full window exists exactly when $is+W\le T$. Hence its count is
\[
 N=\max\!\left(0,1+\left\lfloor\frac{T-W}{s}\right\rfloor\right).
\]
For $T=1000$, $W=128$, and $s=64$, there are $N=14$ windows, with
final start $13\cdot64=832$ and final included sample $959$.
The remaining $40$ samples are outside this grid's windows. If a class
must occupy at least fraction $\tau=0.6$ of a $128$-sample window,
the required integer count is $\lceil\tau W\rceil=77$.
\paragraph{Why ECDR needs it.}
The overlap structure, valid CNN output length, and available transition
examples all depend on exact indexing. A one-position convention change
can change labels or invalidate an equality used to reuse activations.
\paragraph{Conditions and connections.}
An inclusive annotation interval $[a,b]$ contains $b-a+1$ integer
positions; a half-open slice $[a,b)$ contains $b-a$. Neither convention
should be inferred from notation alone. The index extent $W/f_s$ and
first-to-last timestamp gap $(W-1)/f_s$ are distinct time quantities.
\par\label{endconcept:F08}
\typeout{ECDR-CONCEPT-END F08 PAGE=\thepage}

\concept{F09}{Finite sums and reindexing}{\hyperref[concept:F05]{F05}, \hyperref[concept:F08]{F08}}
\paragraph{Foundation and definition.}
A finite sum repeatedly adds terms over a declared finite index set.
Associativity and commutativity of real addition justify regrouping and
reordering. Reindexing is valid when the new indices describe exactly the
same terms once each, or when explicitly stated terms are added or removed.
\paragraph{Symbols and meaning.}
Let $n\ge1$ be an integer, let $a_0,\ldots,a_{n-1}\in\R$ be scalars,
and let $i$ be a summation index. Then $\sum_{i=0}^{n-1}a_i$ adds those
$n$ terms. For vector terms the same notation means coordinatewise
addition. A sum over an empty index set is defined as the additive
identity, either zero or the zero vector as appropriate.
\paragraph{Derivation or example.}
Let $H_0,\ldots,H_{n+s-1}\in\R^d$ be feature vectors, with integers
$1\le s\le n$ and $d\ge1$. Reindexing gives
\[
 \sum_{i=0}^{n-1}H_{i+s}-\sum_{i=0}^{n-1}H_i
 =\sum_{i=n}^{n+s-1}H_i-\sum_{i=0}^{s-1}H_i.
\]
Both left-hand sums contain the positions $s,\ldots,n-1$, whose terms
cancel. For $n=4,s=2$, this is
$(H_2+H_3+H_4+H_5)-(H_0+H_1+H_2+H_3)
=H_4+H_5-H_0-H_1$. Dividing by $n$ yields a difference of averages.
\paragraph{Why ECDR needs it.}
This cancellation explains why a global-average-pooled feature can change
through entering and leaving feature positions when alignment is exact.
It is an algebraic consequence, not an empirical statement that all
nearby windows have small differences.
\paragraph{Conditions and connections.}
The shared symbols must represent equal features, established by a
separate receptive-field and shift argument. Equal input-window lengths
alone do not establish that equality. In floating-point arithmetic,
different summation orders can produce different rounded outputs.
\par\label{endconcept:F09}
\typeout{ECDR-CONCEPT-END F09 PAGE=\thepage}

\concept{F10}{Vector spaces and their axioms}{\hyperref[concept:F05]{F05}}
\paragraph{Foundation and definition.}
A real vector space is a set $V$ equipped with vector addition and real
scalar multiplication. Addition is associative and commutative, has a
zero vector, and supplies an additive inverse for every vector. Scalar
multiplication satisfies identity, compatibility, and both distributive
laws. Closure requires the resulting objects to remain in $V$.
\paragraph{Symbols and meaning.}
Let $u,v,w\in V$ be vectors and $a,b\in\R$ be scalars. The remaining
laws are $1v=v$, $(ab)v=a(bv)$,
$a(u+v)=au+av$, and $(a+b)v=av+bv$. In the concrete space $\R^d$,
the positive integer $d$ counts coordinates, and these operations act
coordinatewise. The zero vector is written $0_V$ when confusion with
the real scalar zero is possible.
\paragraph{Derivation or example.}
For $u=(1,2)$ and $v=(3,-1)$ in $\R^2$,
$2u-v=(-1,5)$. Addition is legitimate because corresponding entries use
the same coordinate meanings. A set of physically realizable sensor
windows need not itself be a vector space: adding two realizable windows
or multiplying one by a negative number may leave the physical set.
The ambient coordinate space can still be a vector space containing it.
\paragraph{Why ECDR needs it.}
Differences between windows, linear heads, and perturbation directions
are vectors in an ambient space. This makes linear algebra available
without assuming that every resulting vector is a plausible activity.
\paragraph{Conditions and connections.}
Coordinate scaling must be fixed when physical channels have different
units. An acceleration coordinate in $g$ and an angular-velocity coordinate
in rad/s are distinct types of measurements. Normalized coordinates
support a declared dimensionless geometry. \hyperref[concept:F11]{F11}--\hyperref[concept:F14]{F14} develop combinations,
span, independence, and coordinates from these axioms.
\par\label{endconcept:F10}
\typeout{ECDR-CONCEPT-END F10 PAGE=\thepage}

\concept{F11}{Linear combinations}{\hyperref[concept:F09]{F09}, \hyperref[concept:F10]{F10}}
\paragraph{Foundation and definition.}
A linear combination forms a new vector by scaling finitely many vectors
and adding the results. It is a definition built from the vector-space
operations. Negative coefficients and coefficients whose sum differs
from one are allowed; convex and affine combinations impose additional
constraints and are different concepts.
\paragraph{Symbols and meaning.}
Let the integers $d\ge1$ and $r\ge1$ count coordinates and supplied
directions. Let $v_1,\ldots,v_r\in\R^d$ be direction vectors and let
$\alpha_1,\ldots,\alpha_r\in\R$ be scalar coefficients. Their linear
combination is
\[
 v=\sum_{j=1}^{r}\alpha_j v_j.
\]
Here $j$ is an index, and $v\in\R^d$ is the resulting vector.
If the directions are dimensionless and the coefficients are
dimensionless, so is the combination.
\paragraph{Derivation or example.}
Take $v_1=(1,1)$ and $v_2=(1,-1)$. The coefficients
$\alpha_1=2$ and $\alpha_2=1$ give $v=(3,1)$.
Collecting the directions as columns of
$B=\begin{pmatrix}1&1\\1&-1\end{pmatrix}$ and coefficients as
$\alpha=(2,1)^\top$ expresses the same operation as $v=B\alpha$.
Restricting coefficients to $|\alpha_1|\le2$ and $|\alpha_2|\le1$
selects a bounded parallelogram from all possible combinations.
\paragraph{Why ECDR needs it.}
A structured perturbation model can represent correlated changes using
several directions instead of an independent allowance at every input
coordinate. The coefficient bounds specify how much motion along each
direction is admitted.
\paragraph{Conditions and connections.}
Representability does not establish that a direction is physically
meaningful, statistically likely, or safe for the classifier. Those are
separate modeling, empirical, and verification questions. \hyperref[concept:F12]{F12} allows all
coefficients to define a span; a bounded coefficient box gives only a
subset of that span, which must not be confused with the whole subspace.
\par\label{endconcept:F11}
\typeout{ECDR-CONCEPT-END F11 PAGE=\thepage}

\concept{F12}{Linear span: all combinations and the smallest subspace}{\hyperref[concept:F10]{F10}, \hyperref[concept:F11]{F11}}
\paragraph{Foundation and definition.}
The linear span of vectors is the set of every vector obtainable as a
linear combination. For directions $v_1,\ldots,v_r\in\R^d$, define
\[
 \operatorname{span}\{v_1,\ldots,v_r\}
 =\left\{\sum_{j=1}^{r}\alpha_jv_j:
                   \alpha_1,\ldots,\alpha_r\in\R\right\}.
\]
The positive integers $r$ and $d$ count supplied directions and ambient
coordinates. All coefficients $\alpha_j$ are unrestricted real numbers.
This linear-algebra term differs from a temporal annotation span.
\paragraph{Symbols and meaning.}
A subspace is a subset containing the zero vector and closed under
addition and real scaling. Here $j$ indexes directions; a member of the
span is a possible direction combination. The word ``linear'' implies
that the resulting set passes through the origin.
\paragraph{Derivation or example.}
For $v_1=(1,2)$, its span is the line
$\{(a,2a):a\in\R\}$. Adding $v_2=(0,1)$ produces the entire plane:
an arbitrary $(x,y)\in\R^2$ equals $xv_1+(y-2x)v_2$.
The span is the \emph{smallest} subspace containing the supplied vectors.
Proof: zero coefficients give zero; adding combinations adds their
coefficients; scaling scales the coefficients, so the span is a subspace.
Any subspace containing the supplied vectors must contain every scaled
vector and every finite sum of them, hence contains the span.
\paragraph{Why ECDR needs it.}
If a matrix $B$ contains learned perturbation directions as columns,
$B\alpha$ stays in their span. A residual term is needed to represent
changes outside that subspace; discarding it without evidence can exclude
real queries from the modeled region.
\paragraph{Conditions and connections.}
Dependent directions do not increase span dimension. Bounding
coefficients changes an unbounded subspace into a restricted region.
Translating by an anchor creates an affine set. \hyperref[concept:F13]{F13}--\hyperref[concept:F16]{F16} explain these
distinctions; \hyperref[concept:F19]{F19} explains orthogonal residuals.
\par\label{endconcept:F12}
\typeout{ECDR-CONCEPT-END F12 PAGE=\thepage}

\concept{F13}{Linear independence and redundant directions}{\hyperref[concept:F11]{F11}, \hyperref[concept:F12]{F12}}
\paragraph{Foundation and definition.}
Vectors are linearly independent when the only linear combination equal
to zero uses zero for every coefficient. Otherwise they are dependent.
Independence is an algebraic property. It does not mean probabilistic
independence, uncorrelated measurements, or distinct class labels.
\paragraph{Symbols and meaning.}
For vectors $v_1,\ldots,v_r\in\R^d$, with integers $r,d\ge1$, and
coefficients $\alpha_1,\ldots,\alpha_r\in\R$, independence means
\[
 \sum_{j=1}^{r}\alpha_jv_j=0
 \quad\Rightarrow\quad
 \alpha_1=\cdots=\alpha_r=0.
\]
Here $0$ on the left is the $d$-dimensional zero vector, and $j$ indexes
directions. A nonzero coefficient vector producing zero demonstrates
redundancy among the directions.
\paragraph{Derivation or example.}
The vectors $v_1=(1,2)$ and $v_2=(2,4)$ are dependent because
$2v_1-v_2=0$. Their span remains one line despite two supplied columns.
In contrast, $u_1=(1,0)$ and $u_2=(1,1)$ are independent: if
$a u_1+b u_2=(a+b,b)=(0,0)$, then $b=0$ and $a=0$.
Independent directions also give unique coefficients for any vector in
their span. If two coefficient lists represent the same vector,
subtracting them gives a zero combination; independence forces each
coefficient difference to vanish.
\paragraph{Why ECDR needs it.}
An independent basis prevents ambiguity when coefficients describe a
structured perturbation. With redundant columns, the same perturbation
can have several coefficient vectors, so testing one arbitrary encoding
against coefficient bounds may not settle set membership.
\paragraph{Conditions and connections.}
Near-dependence is a numerical conditioning issue even when exact
independence holds. A tolerance-based rank estimate should therefore be
reported as a numerical decision. No more than $d$ vectors can be
independent in $\R^d$. \hyperref[concept:F14]{F14} uses independence together with spanning to
define a basis.
\par\label{endconcept:F13}
\typeout{ECDR-CONCEPT-END F13 PAGE=\thepage}

\concept{F14}{Bases, coordinates, and dimension}{\hyperref[concept:F12]{F12}, \hyperref[concept:F13]{F13}}
\paragraph{Foundation and definition.}
A basis of a vector space is a linearly independent set that spans the
space. Every vector then has a unique linear-combination representation.
For a finite-dimensional space, its dimension is the number of elements
in any basis; equality of basis sizes is a theorem, not the definition
of independence.
\paragraph{Symbols and meaning.}
Let $V\subseteq\R^d$ be a subspace of dimension $r$, where integers
$1\le r\le d$. Let $b_1,\ldots,b_r\in\R^d$ be a basis and
$B\in\R^{d\times r}$ the matrix with these vectors as columns.
The coordinate vector $\alpha\in\R^r$ of $v\in V$ satisfies
$v=B\alpha$. The ambient dimension $d$ and subspace dimension $r$ have
different meanings.
\paragraph{Derivation or example.}
Take $b_1=(1,1)$ and $b_2=(1,-1)$ as a basis of $\R^2$.
For $v=(3,1)$, solving
$\alpha_1+\alpha_2=3$ and $\alpha_1-\alpha_2=1$ gives the unique
coordinates $\alpha=(2,1)^\top$. Changing the basis changes these
coordinates without changing $v$. For example, standard-coordinate
axes represent the same vector by $(3,1)^\top$.
\paragraph{Why ECDR needs it.}
A low-dimensional model describes perturbations using $r$ coefficients
instead of $d$ ambient coordinates. This can express measured structure,
but dimensionality alone does not establish classifier invariance or
an efficient online membership algorithm.
\paragraph{Conditions and connections.}
A truncated direction matrix spans a subspace rather than the whole
input space. A residual must account for omitted directions when they
are allowed. For orthonormal columns, coordinates are $B^\top v$;
for a general independent basis they are
$(B^\top B)^{-1}B^\top v$. Using the first formula without
orthonormality generally gives incorrect coordinates.
\par\label{endconcept:F14}
\typeout{ECDR-CONCEPT-END F14 PAGE=\thepage}

\concept{F15}{Matrices as linear maps}{\hyperref[concept:F09]{F09}, \hyperref[concept:F10]{F10}, \hyperref[concept:F14]{F14}}
\paragraph{Foundation and definition.}
A map is linear if it preserves vector addition and scalar
multiplication. A real matrix represents a linear map after input and
output bases are chosen. Conversely, each linear map between
finite-dimensional coordinate spaces has a matrix representation.
\paragraph{Symbols and meaning.}
Let $A\in\R^{m\times d}$ be a matrix, $x\in\R^d$ an input, and
$y=Ax\in\R^m$ an output. Positive integers $d$ and $m$ count input
and output coordinates. Entry $A_{ji}$ multiplies input coordinate
$x_i$ when forming output coordinate $y_j$:
\[
 y_j=\sum_{i=1}^{d}A_{ji}x_i,
 \qquad j=1,\ldots,m.
\]
The transpose $A^\top\in\R^{d\times m}$ exchanges row and column
indices; it is not generally an inverse.
\paragraph{Derivation or example.}
For
$A=\begin{pmatrix}2&-1\\0&3\end{pmatrix}$ and $x=(1,2)^\top$,
the output is $(0,6)^\top$. Distributivity proves
$A(x+z)=Ax+Az$ for any $z\in\R^d$: each row distributes over the
corresponding scalar sum. Similarly, $A(a x)=a(Ax)$ for any scalar
$a\in\R$. A map containing a fixed nonzero bias generally fails
linearity because its output at zero is not zero.
\paragraph{Why ECDR needs it.}
A convolution with fixed weights and no bias can be represented as a
sparse matrix. A final linear head is a small dense matrix. This shared
language allows coordinate differences and norm bounds to propagate
through both kinds of layers.
\paragraph{Conditions and connections.}
Shape compatibility is necessary for composition and multiplication.
Different flattening orders require corresponding matrix-column orders.
Physical units must also match: a coefficient transforms input units
into output units. \hyperref[concept:F16]{F16} adds biases; \hyperref[concept:F21]{F21} measures worst-case amplification.
\par\label{endconcept:F15}
\typeout{ECDR-CONCEPT-END F15 PAGE=\thepage}

\concept{F16}{Affine maps and translated structure}{\hyperref[concept:F11]{F11}, \hyperref[concept:F15]{F15}}
\paragraph{Foundation and definition.}
An affine map has the form of a linear map plus a constant vector.
It preserves affine combinations, whose scalar coefficients sum to one.
It need not preserve the origin. A translated subspace is an affine set,
which should not be called a linear subspace unless it contains zero.
\paragraph{Symbols and meaning.}
Let $A\in\R^{m\times d}$ be fixed weights,
$b\in\R^m$ a fixed bias, and $x\in\R^d$ an input. Integers
$m,d\ge1$ count output and input coordinates. Define
$f(x)=Ax+b$. For an anchor $x_0\in\R^d$ and perturbation
$\delta\in\R^d$, the new input is $x_0+\delta$.
\paragraph{Derivation or example.}
The difference identity is exact over real numbers:
\[
 f(x_0+\delta)-f(x_0)
 =A(x_0+\delta)+b-(Ax_0+b)=A\delta.
\]
Thus a fixed bias affects absolute outputs but cancels from their
differences. In one dimension, $f(x)=2x+3$ gives
$f(4)-f(1)=11-5=6=2(4-1)$.
If $u,v\in\R^d$ and $\lambda\in\R$, then
$f(\lambda u+(1-\lambda)v)=\lambda f(u)+(1-\lambda)f(v)$,
because the two coefficients of $b$ sum to one.
\paragraph{Why ECDR needs it.}
Class logits and fixed normalization are affine operations. Their
differences can often be analyzed more tightly than separate absolute
output intervals. Translating a perturbation domain by an anchor also
explains precisely where a local certificate applies.
\paragraph{Conditions and connections.}
The same $A$ and $b$ must be used at both inputs. Updating running
normalization statistics or model weights breaks that cancellation.
ReLU layers make a network piecewise affine; one affine formula applies
only when its activation conditions are established.
\par\label{endconcept:F16}
\typeout{ECDR-CONCEPT-END F16 PAGE=\thepage}

\concept{F17}{Composition and the complete reference function}{\hyperref[concept:F04]{F04}, \hyperref[concept:F15]{F15}, \hyperref[concept:F16]{F16}}
\paragraph{Foundation and definition.}
Composition applies one function after another. If the first function's
outputs lie within the second function's domain, the composition is
well defined. Writing a pipeline as a composition exposes which changes
alter the mathematical reference being certified.
\paragraph{Symbols and meaning.}
Let $P:\R^{d}\to\R^{d}$ be preprocessing,
$G:\R^{d}\to\R^{r}$ the feature extractor,
$L:\R^{r}\to\R^{K}$ the logit head, and
$D:\R^{K}\to\{1,\ldots,K\}$ the decision rule. The positive
integers $d,r,K$ count input coordinates, features, and classes.
Define $h=D\circ L\circ G\circ P$, so
$h(x)=D(L(G(P(x))))$ for input $x\in\R^d$.
\paragraph{Derivation or example.}
In a scalar example, let $P(x)=(x-2)/4$ and $G(u)=\max(0,u)$.
Then $(G\circ P)(6)=1$, while $(G\circ P)(0)=0$.
Reversing the order gives $(P\circ G)(0)=-0.5$, which differs from
zero. Composition is therefore generally noncommutative. It is
associative when all involved domains fit: grouping stages does not
change the order of evaluation.
\paragraph{Why ECDR needs it.}
A certificate for logits under one channel ordering or normalization is
not automatically a certificate for a changed pipeline. The reference
includes preprocessing, inference state, weights, and deterministic tie
handling, even if only the CNN weights are stored in a checkpoint.
\paragraph{Conditions and connections.}
Every stage must receive the shape, units, and coordinate conventions it
expects. A certificate built in normalized coordinates needs the query
mapped by the same fixed $P$. \hyperref[concept:F21]{F21} shows how amplification bounds
multiply through composition; nonlinear stages require their own bounds.
\par\label{endconcept:F17}
\typeout{ECDR-CONCEPT-END F17 PAGE=\thepage}

\concept{F18}{Inner products and dot products}{\hyperref[concept:F09]{F09}, \hyperref[concept:F10]{F10}}
\paragraph{Foundation and definition.}
A real inner product is a symmetric bilinear operation producing a
scalar and satisfying positive definiteness: a vector's product with
itself is nonnegative and vanishes only at zero. The standard dot
product on $\R^d$ is one particular inner product.
\paragraph{Symbols and meaning.}
For vectors $u,v\in\R^d$, where integer $d\ge1$ counts coordinates,
write
\[
 \langle u,v\rangle=u^\top v=\sum_{i=1}^{d}u_i v_i.
\]
The index $i$ selects a coordinate. Symmetry means
$\langle u,v\rangle=\langle v,u\rangle$; bilinearity means linearity
in either argument when the other is fixed. A weighted product
$u^\top Mv$ is an inner product when the matrix
$M\in\R^{d\times d}$ is symmetric positive definite.
\paragraph{Derivation or example.}
With $u=(1,2)$ and $v=(3,-1)$, their dot product is $3-2=1$.
Also $u^\top u=5$. For a head row difference $q=(2,-1)$ and feature
change $\delta=(0.1,0.3)$, the scalar change is
$q^\top\delta=0.2-0.3=-0.1$. Opposing coordinate contributions can
partly cancel; their separate absolute magnitudes would sum to $0.5$.
\paragraph{Why ECDR needs it.}
A class-pair margin is a directional measurement of a feature vector.
Keeping that direction can preserve cancellations that separate logit
bounds lose. Inner products also define orthogonality, projections, and
the usual geometry of learned directions.
\paragraph{Conditions and connections.}
The standard dot product implicitly chooses a coordinate scale. Adding
products of physically different units needs a declared normalization
or weighting interpretation. \hyperref[concept:F19]{F19} develops orthogonality; \hyperref[concept:F22]{F22} derives the
dual-norm inequality that bounds a directional change by its magnitude.
\par\label{endconcept:F18}
\typeout{ECDR-CONCEPT-END F18 PAGE=\thepage}

\concept{F19}{Orthogonality and residual decomposition}{\hyperref[concept:F18]{F18}}
\paragraph{Foundation and definition.}
Two vectors are orthogonal when their inner product is zero. This
algebraic relation means a right angle for nonzero vectors under the
chosen geometry. Orthonormal vectors are pairwise orthogonal and each
has unit Euclidean length. Neither condition is equivalent to
probabilistic independence.
\paragraph{Symbols and meaning.}
Let $U\in\R^{d\times r}$ have orthonormal columns, with integers
$1\le r\le d$. Thus $U^\top U=I_r$, where $I_r$ is the
$r\times r$ identity matrix. For a vector $\delta\in\R^d$, define
its coefficient vector $\alpha=U^\top\delta\in\R^r$ and residual
$e=\delta-U\alpha\in\R^d$.
\paragraph{Derivation or example.}
Multiplying the residual by $U^\top$ gives
\[
 U^\top e=U^\top\delta-U^\top U\alpha
 =\alpha-\alpha=0.
\]
The residual is therefore orthogonal to every supplied direction, and
$\delta=U\alpha+e$. For $U=(1,0)^\top$ and $\delta=(3,4)^\top$,
the projected part is $(3,0)^\top$ and the residual is $(0,4)^\top$.
Expanding the dot product gives Pythagoras:
$\norm{\delta}_2^2=\norm{U\alpha}_2^2+\norm{e}_2^2=9+16=25$,
because the cross term vanishes.
\paragraph{Why ECDR needs it.}
A direction-based domain can distinguish explained motion from
unexplained residual motion. Bounding the residual prevents a
low-dimensional approximation from silently asserting that all real
changes lie exactly in the learned subspace.
\paragraph{Conditions and connections.}
Column scaling generally destroys orthonormality, so $U^\top\delta$
cannot then be reused as a coefficient formula. Orthogonality is relative
to an inner product, which must match normalization and membership
calculations. \hyperref[concept:C11]{C11} explains projection; \hyperref[concept:C12]{C12} explains how PCA estimates
directions from finite data.
\par\label{endconcept:F19}
\typeout{ECDR-CONCEPT-END F19 PAGE=\thepage}

\concept{F20}{Norms and the meanings of vector size}{\hyperref[concept:F10]{F10}, \hyperref[concept:F18]{F18}}
\paragraph{Foundation and definition.}
A norm assigns a nonnegative size to each vector. Its axioms are
definiteness, absolute homogeneity, and the triangle inequality:
$\norm{v}=0$ exactly at $v=0$,
$\norm{a v}=|a|\norm{v}$, and
$\norm{u+v}\le\norm{u}+\norm{v}$.
Here $u,v$ are vectors in one real vector space and $a\in\R$ is a scalar.
\paragraph{Symbols and meaning.}
For $v\in\R^d$, with integer $d\ge1$ and coordinate index
$i\in\{1,\ldots,d\}$, common choices are
\[
 \norm{v}_1=\sum_{i=1}^d|v_i|,\qquad
 \norm{v}_2=\sqrt{\sum_{i=1}^d v_i^2},\qquad
 \norm{v}_\infty=\max_{1\le i\le d}|v_i|.
\]
The subscripts specify different notions of size; $\infty$ names the
maximum-coordinate norm, not an infinite vector size.
\paragraph{Derivation or example.}
For $v=(3,-4)$, these sizes are $7$, $5$, and $4$. More generally,
\[
 \norm{v}_\infty\le\norm{v}_2\le\norm{v}_1,
 \qquad \norm{v}_2\le\sqrt d\,\norm{v}_\infty.
\]
For example, the last inequality follows by bounding each of the $d$
squared coordinates by $\norm{v}_\infty^2$ and taking square roots.
Thus changing the norm while keeping the same numerical radius changes
the admitted perturbations.
\paragraph{Why ECDR needs it.}
Norms define perturbation budgets, Lipschitz changes, and residual size.
A maximum-coordinate budget constrains every coordinate, whereas an
Euclidean budget constrains total squared change.
\paragraph{Conditions and connections.}
Norm values depend on coordinate scaling. Raw acceleration and angular
velocity should not be treated as physically commensurate just because
their numerical entries share an array. \hyperref[concept:F22]{F22} pairs norms with their duals;
\hyperref[concept:F23]{F23} turns a norm into a distance and a ball.
\par\label{endconcept:F20}
\typeout{ECDR-CONCEPT-END F20 PAGE=\thepage}

\concept{F21}{Induced operator norms and composition bounds}{\hyperref[concept:F15]{F15}, \hyperref[concept:F20]{F20}}
\paragraph{Foundation and definition.}
An induced operator norm measures the largest amplification of vector
size by a linear map. For compatible input and output norms, define
\[
 \norm{A}_{p\to q}
 =\sup_{v\ne0}\frac{\norm{Av}_q}{\norm{v}_p}.
\]
Here $A\in\R^{m\times d}$, $v\in\R^d$, integers $m,d\ge1$
count coordinates, and subscripts $p,q$ name the chosen norms.
\paragraph{Symbols and meaning.}
The supremum is the least upper bound over all nonzero input directions.
The notation $\norm{A}_2$ usually means $\norm{A}_{2\to2}$, not
the square root of the sum of every squared matrix entry. The latter is
the Frobenius norm and is a different quantity.
\paragraph{Derivation or example.}
The definition gives $\norm{Av}_q\le\norm{A}_{p\to q}\norm{v}_p$.
For another matrix $B\in\R^{d\times r}$, with integer $r\ge1$ and
input norm indexed by $t$, applying this inequality twice yields
\[
 \norm{AB}_{t\to q}
 \le\norm{A}_{p\to q}\norm{B}_{t\to p}.
\]
This is submultiplicativity. For
$A=\operatorname{diag}(2,-3)$, the Euclidean induced norm is $3$:
each squared output is bounded by $9$ times the squared input, and
input $(0,1)$ attains equality. Under maximum-coordinate norms,
$\norm{A}_{\infty\to\infty}=\max_j\sum_i|A_{ji}|$.
\paragraph{Why ECDR needs it.}
Layerwise amplification bounds can be multiplied to obtain a network
bound. This is simple and sound when every layer constant is valid,
although successive worst-case directions may not be jointly achievable.
\paragraph{Conditions and connections.}
The intermediate norms and dimensions must match. A computed spectral
estimate can underestimate the true norm unless certified with an error
bound. Biases cancel from affine differences; nonlinear stages need
separate Lipschitz arguments.
\par\label{endconcept:F21}
\typeout{ECDR-CONCEPT-END F21 PAGE=\thepage}

\concept{F22}{H\"older's inequality and dual norms}{\hyperref[concept:F18]{F18}, \hyperref[concept:F20]{F20}}
\paragraph{Foundation and definition.}
The dual norm measures the largest dot product a direction can produce
against a unit-norm vector. For a norm $\norm{\cdot}$ on $\R^d$,
define $\norm{q}_*=\sup_{\norm{v}\le1}|q^\top v|$.
Here integer $d\ge1$ counts coordinates, while $q,v\in\R^d$ are
a direction and a test vector. The star names the dual norm.
\paragraph{Symbols and meaning.}
For $L_p$ norms with $1\le p\le\infty$, the dual is $L_{p^*}$,
where $1/p+1/p^*=1$ and $1/\infty=0$. Thus $L_1$ pairs with
$L_\infty$, $L_2$ pairs with itself, and $L_\infty$ pairs with $L_1$.
The vector $\delta\in\R^d$ below denotes a perturbation.
\paragraph{Derivation or example.}
If $\delta\ne0$, then $v=\delta/\norm{\delta}$ is a unit vector.
Substitution into the dual-norm definition gives
\[
 |q^\top\delta|\le\norm{q}_*\norm{\delta}.
\]
At $\delta=0$ the result follows directly. This is H\"older's
inequality in dual-norm form. If $q=(2,-1)$ and
$\norm{\delta}_\infty\le0.1$, the largest absolute directional
change is $(|2|+|-1|)0.1=0.3$, attained at
$\delta=(0.1,-0.1)$ or its negative. Using $\norm{q}_\infty$
instead would incorrectly claim the upper bound $0.2$.
\paragraph{Why ECDR needs it.}
A margin direction and a perturbation budget become a scalar bound on
possible margin loss. The dual pairing determines the correct constant;
choosing two matching-looking subscripts can be unsound.
\paragraph{Conditions and connections.}
Equality is possible for the complete norm ball. Additional structural
constraints can reduce the attainable loss, so the same inequality may
be conservative. Nonuniform coordinate budgets lead to a weighted
absolute-value sum rather than one common radius.
\par\label{endconcept:F22}
\typeout{ECDR-CONCEPT-END F22 PAGE=\thepage}

\concept{F23}{Metrics, balls, and neighborhoods}{\hyperref[concept:F03]{F03}, \hyperref[concept:F20]{F20}}
\paragraph{Foundation and definition.}
A metric is a distance function that is nonnegative, vanishes exactly
when its two arguments agree, is symmetric, and satisfies the triangle
inequality. Every norm supplies a metric by taking the norm of a
difference. A ball is the set of points within a specified distance
from one center.
\paragraph{Symbols and meaning.}
Let $x,y,z\in\R^d$, with integer $d\ge1$ counting coordinates.
Write $d_p(x,y)=\norm{x-y}_p$ for distance under norm $L_p$.
For center $x_0\in\R^d$ and radius $\varepsilon>0$, define
\[
 B_p(x_0,\varepsilon)=\{x:d_p(x,x_0)<\varepsilon\},\qquad
 \overline B_p(x_0,\varepsilon)=\{x:d_p(x,x_0)\le\varepsilon\}.
\]
The first is open and the second is closed. A neighborhood of $x_0$
contains some open ball around $x_0$.
\paragraph{Derivation or example.}
For $x_0=(0,0)$ and $\varepsilon=1$, the Euclidean ball is a disk;
the maximum-coordinate ball is a square. The point $(1,0)$ is in both
closed balls but neither open ball. The triangle inequality gives
$d_p(x,z)\le d_p(x,y)+d_p(y,z)$, so every point within radius
$r$ of $y$ lies within radius $r+d_p(y,x_0)$ of anchor $x_0$.
Here $r\ge0$ is another radius, measured in the same coordinate units.
\paragraph{Why ECDR needs it.}
Local certification needs an exact meaning of ``near the anchor.''
Triangle inequalities can transfer distance allowances, but the enlarged
set must still fit inside a region with a valid decision proof.
\paragraph{Conditions and connections.}
Proximity alone does not imply equal class labels. A decision boundary
can pass arbitrarily close to an input. The norm, normalization, and
boundary convention belong to the certificate's contract. \hyperref[concept:F26]{F26} formalizes
local change; \hyperref[concept:F28]{F28} connects continuity with positive margin.
\par\label{endconcept:F23}
\typeout{ECDR-CONCEPT-END F23 PAGE=\thepage}

\concept{F24}{Intervals, Cartesian products, and boxes}{\hyperref[concept:F03]{F03}, \hyperref[concept:F06]{F06}}
\paragraph{Foundation and definition.}
A closed real interval includes every scalar between its two endpoints,
including the endpoints. A Cartesian product collects tuples whose
coordinates belong to specified sets. A box is the Cartesian product
of coordinate intervals. It explicitly permits every listed coordinate
combination, whether or not those combinations appeared in a dataset.
\paragraph{Symbols and meaning.}
Let $d\ge1$ be an integer and $\ell,u\in\R^d$ be finite lower and
upper endpoints satisfying $\ell_i\le u_i$ for indices
$i=1,\ldots,d$. The box is
\[
 [\ell,u]=\prod_{i=1}^{d}[\ell_i,u_i]
 =\{x\in\R^d:\ell_i\le x_i\le u_i\text{ for every }i\}.
\]
An equivalent anchor representation is $x=x_0+\delta$ with
$|\delta_i|\le b_i$, where $x_0=(\ell+u)/2$ and
$b=(u-\ell)/2\in\R^d$ has nonnegative entries.
\paragraph{Derivation or example.}
The intervals $[-1,1]$ and $[2,4]$ form a rectangular box in $\R^2$.
For $f(x)=3x_1-2x_2$, minimizing each independent coordinate term
gives $\min f=3(-1)-2(4)=-11$, attained at $(-1,4)$.
The negative coefficient chooses the upper endpoint, illustrating why
order reversal matters in an affine lower bound.
\paragraph{Why ECDR needs it.}
Boxes support simple coordinatewise membership and efficient interval
propagation. They can overapproximate correlated perturbations by
admitting combinations the data seldom produces; verification must
still cover all admitted combinations.
\paragraph{Conditions and connections.}
Zero-width intervals are legitimate and fix coordinates. Positive widths
are needed when a later argument assumes independent variability in
every coordinate. A box is nonempty under the endpoint inequalities;
reversing an endpoint pair makes its product empty.
\par\label{endconcept:F24}
\typeout{ECDR-CONCEPT-END F24 PAGE=\thepage}

\concept{F25}{Infimum and supremum versus minimum and maximum}{\hyperref[concept:F06]{F06}, \hyperref[concept:F07]{F07}}
\paragraph{Foundation and definition.}
The infimum is the greatest lower bound of a real-valued set, while the
supremum is its least upper bound. A minimum or maximum additionally
must be attained by a member. These definitions separate a best bound
from an actual optimizing input.
\paragraph{Symbols and meaning.}
Let $f:\mathcal S\to\R$ be a scalar function on a nonempty set
$\mathcal S$. Write
$\inf_{x\in\mathcal S}f(x)$ for the infimum of its output set.
A point $x_*\in\mathcal S$ is a minimizer when
$f(x_*)\le f(x)$ for all $x\in\mathcal S$; its value is
$\min_{x\in\mathcal S}f(x)$. The star identifies a candidate
optimizer and does not denote multiplication.
\paragraph{Derivation or example.}
For $f(x)=x$ on $\mathcal S=(0,1)$, the infimum is $0$ but no
minimum exists. Every permitted input has a positive output, yet no
single positive lower bound works throughout the interval. On
$[0,1]$, the minimum exists and equals zero at $x_*=0$.
For a valid enclosure $L\le f(x)$ for every permitted $x$, one has
\[
 L\le\inf_{x\in\mathcal S}f(x).
\]
Here $L\in\R$ is the verifier's lower bound, not necessarily the
exact infimum and not necessarily an attained value.
\paragraph{Why ECDR needs it.}
A positive sound lower bound proves all margins are positive. A loose
nonpositive bound does not identify a violating input. Separating the
bound, true infimum, and optimizer prevents that logical mistake.
\paragraph{Conditions and connections.}
For unbounded-below output sets the infimum may be $-\infty$ in the
extended reals. Empty domains need a separately declared convention and
are excluded here. \hyperref[concept:F27]{F27}--\hyperref[concept:F28]{F28} give sufficient conditions for finite,
attained extrema.
\par\label{endconcept:F25}
\typeout{ECDR-CONCEPT-END F25 PAGE=\thepage}

\concept{F26}{Continuity and the epsilon--delta definition}{\hyperref[concept:F23]{F23}}
\paragraph{Foundation and definition.}
Continuity means that a sufficiently small input change produces any
prescribed small output change. It is a local property of a function,
defined through distances. It does not mean that the function is
constant nearby or that one global rate of change has been proved.
\paragraph{Symbols and meaning.}
Let $f:\mathcal S\subseteq\R^d\to\R^m$ be a function, and let
$x_0\in\mathcal S$. Positive integers $d,m$ count input and output
coordinates. Continuity at $x_0$ means
\[
 \forall\varepsilon>0\;\exists\delta>0\;\forall x\in\mathcal S:
 \norm{x-x_0}_2<\delta
 \Rightarrow\norm{f(x)-f(x_0)}_2<\varepsilon.
\]
The positive scalar $\varepsilon$ is an output tolerance and $\delta$
is an input tolerance, each in its respective coordinate units.
\paragraph{Derivation or example.}
For the scalar function $f(x)=2x+1$, choose
$\delta=\varepsilon/2$. Then
$|f(x)-f(x_0)|=2|x-x_0|<\varepsilon$ whenever
$|x-x_0|<\delta$. If a continuous scalar margin
$g:\mathcal S\to\R$ satisfies $g(x_0)>0$, continuity with output
tolerance $g(x_0)/2$ supplies some neighborhood where
$g(x)>g(x_0)/2>0$. This proves local existence of
a safe neighborhood, without computing a usable radius.
\paragraph{Why ECDR needs it.}
Affine layers and ReLU are continuous; finite compositions therefore
produce continuous logits. A strictly positive winning margin at an
anchor supports a local invariance argument. Quantitative bounds are
needed to turn that existence statement into an executable gate.
\paragraph{Conditions and connections.}
The discrete class decision can be discontinuous at ties even when
logits are continuous. Continuity on a noncompact region does not alone
give a uniform positive margin. \hyperref[concept:F28]{F28} adds the relevant compactness and
strict-positivity conditions.
\par\label{endconcept:F26}
\typeout{ECDR-CONCEPT-END F26 PAGE=\thepage}

\concept{F27}{Compactness in finite-dimensional Euclidean spaces}{\hyperref[concept:F23]{F23}, \hyperref[concept:F24]{F24}, \hyperref[concept:F26]{F26}}
\paragraph{Foundation and definition.}
A set is compact if every collection of open sets covering it has a
finite subcollection that still covers it. In finite-dimensional
Euclidean space, the Heine--Borel theorem says compactness is equivalent
to being closed and bounded. This equivalence is a theorem specific to
the setting, rather than a universal definition for arbitrary spaces.
\paragraph{Symbols and meaning.}
Let $\mathcal S\subseteq\R^d$, where integer $d\ge1$ counts
coordinates. Boundedness means there is a finite scalar $R\ge0$ with
$\norm{x}_2\le R$ for every $x\in\mathcal S$. Closedness means
that limits of convergent sequences of members remain in $\mathcal S$.
An open cover is a family of open sets whose union contains $\mathcal S$.
\paragraph{Derivation or example.}
A finite-endpoint box $[\ell,u]$ is closed and bounded, hence compact.
For $\ell,u\in\R^d$ with $\ell_i\le u_i$, boundedness follows from
$|x_i|\le\max(|\ell_i|,|u_i|)$ for every coordinate $i$.
The interval $(0,1)$ is bounded but not closed: the sequence
$x_n=1/(n+1)$, for positive integers $n$, lies in $(0,1)$ and converges to zero, outside the
interval. The line $\R$ is closed but unbounded. Neither set is compact.
\paragraph{Why ECDR needs it.}
Bounded closed perturbation domains make worst-case margin optimization
well behaved when margins are continuous. Compactness is the bridge
from infinitely many inputs to existence of an attained worst case;
it does not replace the computation of that worst case.
\paragraph{Conditions and connections.}
The empty set is compact, but useful margin statements should separately
require a nonempty domain. Closed and bounded is insufficient in
general infinite-dimensional normed spaces. ECDR's finite tensor inputs
fit the finite-dimensional hypothesis. \hyperref[concept:F28]{F28} states the resulting extrema
and uniform-positivity theorem precisely.
\par\label{endconcept:F27}
\typeout{ECDR-CONCEPT-END F27 PAGE=\thepage}

\concept{F28}{Attained extrema and a uniform positive margin}{\hyperref[concept:F01]{F01}, \hyperref[concept:F25]{F25}, \hyperref[concept:F26]{F26}, \hyperref[concept:F27]{F27}}
\paragraph{Foundation and definition.}
The extreme-value theorem states that a continuous real-valued function
on a nonempty compact set attains a finite minimum and maximum. A useful
consequence is that a continuous function which is strictly positive
everywhere on such a set has one uniform positive lower bound.
\paragraph{Symbols and meaning.}
Let $\mathcal S\subseteq\R^d$ be nonempty and compact. Let
$z_j:\mathcal S\to\R$ be continuous logit functions for class
indices $j=1,\ldots,K$, with integers $d\ge1$ and $K\ge2$.
Fix class $c$ and define each competing margin
$m_{cj}(x)=z_c(x)-z_j(x)$ for $j\ne c$.
\paragraph{Derivation or example.}
Assume $m_{cj}(x)>0$ for every $x\in\mathcal S$ and competitor
$j\ne c$. Each margin attains its minimum at some
$x_j\in\mathcal S$, so $\eta_j=m_{cj}(x_j)>0$. There are only
$K-1$ competitors; therefore
\[
 \eta=\min_{j\ne c}\eta_j>0,
 \qquad m_{cj}(x)\ge\eta
 \quad\text{for every }x\in\mathcal S, j\ne c.
\]
For $m(x)=2-x$ on $[0,1]$, the minimum is $1$. Dropping closedness
can destroy the conclusion: $m(x)=x$ is positive on $(0,1)$ but has
infimum zero.
\paragraph{Why ECDR needs it.}
This conditional theorem explains why strict decision dominance on a
closed bounded region has a positive clearance. A verifier attempts to
compute a sound usable lower bound; the existence theorem alone does
not give that numerical result or an efficient algorithm.
\paragraph{Conditions and connections.}
Tie-consistent decision constancy is weaker than strict dominance: a
chosen class may win ties while its margin equals zero. Neither
empirical agreement nor an unproved continuity assumption can replace
the theorem's hypotheses. A valid finite-precision certificate needs
additional arithmetic-error control.
\par\label{endconcept:F28}
\typeout{ECDR-CONCEPT-END F28 PAGE=\thepage}

\concept{F29}{Probability spaces and probability axioms}{\hyperref[concept:F03]{F03}, \hyperref[concept:F09]{F09}}
\paragraph{Foundation and definition.}
A probability model consists of a sample space, a collection of
measurable events, and a probability measure. The probability axioms
are nonnegativity, normalization of the whole space to one, and
countable additivity for pairwise disjoint events. They describe a
mathematical model; they do not assert that observed windows are
independent draws from that model.
\paragraph{Symbols and meaning.}
Write $(\Omega,\mathcal F,P)$ for the model: $\Omega$ contains
possible outcomes, $\mathcal F$ is a sigma-algebra of events closed
under complements and countable unions, and $P:\mathcal F\to[0,1]$
assigns probabilities. For pairwise disjoint events
$A_1,A_2,\ldots\in\mathcal F$,
\[
 P(\Omega)=1,\qquad
 P\!\left(\bigcup_{i=1}^{\infty}A_i\right)
 =\sum_{i=1}^{\infty}P(A_i).
\]
The integer $i$ indexes events, not time samples unless defined so.
\paragraph{Derivation or example.}
An event $A$ and its complement $A^c=\Omega\setminus A$ partition
the sample space. Additivity gives $P(A^c)=1-P(A)$.
For a uniformly selected member of a finite dataset of $N\ge1$
windows, each dataset index has probability $1/N$. If $M$ indices pass
a gate, its probability under this particular empirical model is $M/N$.
This construction says nothing by itself about future recordings.
\paragraph{Why ECDR needs it.}
Expected cost and reuse rates refer to a distribution of queries.
Different choices, such as uniform windows or uniform subjects, define
different probability models and can produce different reported rates.
\paragraph{Conditions and connections.}
Distributional coverage is not a universal safety certificate.
Probability one may still allow a nonempty zero-probability exception
set. \hyperref[concept:F30]{F30} explains conditioning and independence; \hyperref[concept:F31]{F31} defines expected
cost after a probability model has been specified.
\par\label{endconcept:F29}
\typeout{ECDR-CONCEPT-END F29 PAGE=\thepage}

\concept{F30}{Conditional probability and independence}{\hyperref[concept:F29]{F29}}
\paragraph{Foundation and definition.}
Conditioning restricts a probability calculation to an event that has
occurred. For events $A$ and $B$ with $P(B)>0$, define
$P(A\mid B)=P(A\cap B)/P(B)$. Events are independent when
$P(A\cap B)=P(A)P(B)$. Independence is an additional property of
the model; overlapping windows do not acquire it by being placed in
different rows of an array.
\paragraph{Symbols and meaning.}
Let $A$ be the event ``a query lies in an admitted domain'' and $B$ the
event ``the current anchor has a successful valid certificate.''
The operator $P$ is the chosen query-distribution probability, and
$A\cap B$ means both events occur. These events describe distinct
gates and can be strongly dependent.
\paragraph{Derivation or example.}
If $P(B)=0.4$ and $P(A\mid B)=0.6$, then
\[
 P(A\cap B)=P(A\mid B)P(B)=0.24.
\]
Substituting an unconditional estimate $P(A)=0.7$ would give $0.28$,
which is unjustified unless independence or another valid argument is
available. For a finite sample, a conditional rate uses only the
eligible denominator: $60$ accepted queries among $100$ queries with
certified anchors gives $0.6$, even if $250$ total queries were seen.
\paragraph{Why ECDR needs it.}
Certified reuse requires all applicable gates together. Multiplying
independently reported coverage and certification percentages can give
the wrong total reuse rate when their denominators or conditions differ.
\paragraph{Conditions and connections.}
The elementary conditional formula is undefined when $P(B)=0$.
More advanced conditional distributions require additional machinery
and are not supplied by writing a ratio with zero denominator.
Statistical independence differs from orthogonality and zero
correlation. \hyperref[concept:F32]{F32}--\hyperref[concept:F33]{F33} explain empirical weights and covariance.
\par\label{endconcept:F30}
\typeout{ECDR-CONCEPT-END F30 PAGE=\thepage}

\concept{F31}{Expectation and the linearity of expected cost}{\hyperref[concept:F09]{F09}, \hyperref[concept:F29]{F29}}
\paragraph{Foundation and definition.}
A random variable is a measurable numerical function of an outcome.
Its expectation is a probability-weighted average, formalized as an
integral. For a finite discrete model, this reduces to a weighted sum.
Linearity of expectation does not require independent variables, but
finite ordinary values require suitable integrability.
\paragraph{Symbols and meaning.}
Let $C:\Omega\to\R$ be a random cost measured in microseconds.
For possible costs $c_i$ and probabilities $p_i\ge0$, indexed by
$i=1,\ldots,n$, with integer $n\ge1$ and $\sum_i p_i=1$, define
$\E[C]=\sum_i p_i c_i$. Integrability means
$\E[|C|]<\infty$. The symbol $\E$ refers to the declared probability
model, not an automatic average over an unspecified benchmark.
\paragraph{Derivation or example.}
Let $I$ indicate reuse, so $I=1$ on reuse and $I=0$ otherwise, with
$P(I=1)=r$. Suppose every query pays fixed gate cost $g$ and a miss
pays fixed forward cost $f$, both in microseconds. Then
\[
 C=g+(1-I)f,\qquad \E[C]=g+(1-r)f.
\]
For $g=10$, $f=100$, and $r=0.7$, expected query cost is $40$.
If per-query build cost $B$ is random and integrable, the expectation
becomes $40+\E[B]$. These are hypothetical numbers, not measurements.
\paragraph{Why ECDR needs it.}
Saving forward calls does not establish positive net savings. Gate,
build, failure, and fallback costs must appear in the same expected
cost calculation, with their actual conditional behavior when variable.
\paragraph{Conditions and connections.}
Medians and high quantiles cannot generally replace expectations in
this identity. Unbounded nonintegrable variables require extended-real
care. Correlation between costs does not break finite-sum linearity;
it does affect variance and uncertainty estimates.
\par\label{endconcept:F31}
\typeout{ECDR-CONCEPT-END F31 PAGE=\thepage}

\concept{F32}{Empirical means, variances, and weighting}{\hyperref[concept:F09]{F09}, \hyperref[concept:F31]{F31}}
\paragraph{Foundation and definition.}
An empirical mean summarizes observed numbers using declared weights.
An empirical population variance measures average squared deviation
under those same weights. These are descriptive definitions. Calling a
variance estimator unbiased requires a sampling model and additional
assumptions, commonly independent identically distributed observations.
\paragraph{Symbols and meaning.}
Let $x_1,\ldots,x_n\in\R$ be measurements, with integer $n\ge1$,
and let $w_i\ge0$ satisfy $\sum_{i=1}^{n}w_i=1$. Define
\[
 \bar x=\sum_{i=1}^{n}w_i x_i,
 \qquad v=\sum_{i=1}^{n}w_i(x_i-\bar x)^2.
\]
The mean $\bar x$ has the measurements' units; variance $v$ has squared
units. Equal weighting uses $w_i=1/n$. Standard deviation is $\sqrt v$.
\paragraph{Derivation or example.}
For equal-weight values $(1,3)$, the mean is $2$ and population
variance is $1$. A recording with one value $0$ and another recording
with nine values $10$ have pooled-sample mean $9$, but an equal-recording
mean of their means is $5$. Both are correct for different weights.
For equal-weight data with $n>1$, the statistic
$s^2=\sum_i(x_i-\bar x)^2/(n-1)$ is a different quantity from $v$.
Under the usual independent identical finite-variance sampling model,
$s^2$ estimates population variance without bias.
\paragraph{Why ECDR needs it.}
Pooling windows weights long recordings more heavily, and overlapping
windows count some raw samples repeatedly. Training normalization and
reported averages must state whether they weight samples, windows,
recordings, or subjects.
\paragraph{Conditions and connections.}
Descriptive means remain defined for dependent windows; ordinary
independent-sample uncertainty formulas may fail. \hyperref[concept:F33]{F33} requires a
consistent weighting convention when reporting covariance or correlation.
\par\label{endconcept:F32}
\typeout{ECDR-CONCEPT-END F32 PAGE=\thepage}

\concept{F33}{Covariance and Pearson correlation}{\hyperref[concept:F18]{F18}, \hyperref[concept:F30]{F30}, \hyperref[concept:F32]{F32}}
\paragraph{Foundation and definition.}
Covariance measures whether paired variables deviate in the same
direction from their means. Pearson correlation divides covariance by
both standard deviations to obtain a dimensionless measure of linear
association. It is undefined when either variance is zero; assigning
zero in that case is a software policy, not the mathematical definition.
\paragraph{Symbols and meaning.}
Let $(x_i,y_i)$, for $i=1,\ldots,n$, be paired real measurements with
integer $n\ge1$ and normalized nonnegative weights $w_i$. Let
$\bar x,\bar y$ be their weighted means, and let $v_x,v_y$ be their
weighted population variances. Define
\[
 c_{xy}=\sum_i w_i(x_i-\bar x)(y_i-\bar y),
 \qquad \rho=\frac{c_{xy}}{\sqrt{v_xv_y}}
 \quad\text{when }v_xv_y>0.
\]
Covariance has the product of the two measurement units; correlation
$\rho$ has no units.
\paragraph{Derivation or example.}
For equally weighted pairs $(1,2)$ and $(3,6)$,
$\bar x=2$, $\bar y=4$, $v_x=1$, $v_y=4$, and $c_{xy}=2$,
so $\rho=1$. Nonetheless, the vectors $(1,3)$ and $(2,6)$ are
different. Zero covariance need not imply independence: for a uniformly
distributed scalar $X$ on $[-1,1]$, $Y=X^2$ is determined by $X$,
yet symmetry gives zero covariance.
\paragraph{Why ECDR needs it.}
Lag-one raw correlation, lag-64 correlation, within-window shape
correlation, and correlation of window means concern different paired
objects. Their values answer different questions and cannot substitute
for a perturbation norm or decision certificate.
\paragraph{Conditions and connections.}
Specify pair construction, lag, weighting, and treatment of constant or
invalid segments. High correlation can coexist with large absolute
differences because scaling and offsets affect distance differently.
\par\label{endconcept:F33}
\typeout{ECDR-CONCEPT-END F33 PAGE=\thepage}

\concept{F34}{Quantiles and empirical distribution functions}{\hyperref[concept:F06]{F06}, \hyperref[concept:F29]{F29}, \hyperref[concept:F32]{F32}}
\paragraph{Foundation and definition.}
A quantile locates a probability threshold in an ordered distribution.
For a finite dataset, the empirical cumulative distribution function
counts how many values lie at or below each threshold. Several
interpolation conventions exist; choosing one is part of a reproducible
analysis and should not be confused with a theorem about future data.
\paragraph{Symbols and meaning.}
For real observations $a_1,\ldots,a_n$ with integer $n\ge1$, define
\[
 F_n(t)=\frac1n\sum_{i=1}^{n}\mathbf1\{a_i\le t\},
 \qquad q_p=\inf\{t\in\R:F_n(t)\ge p\},\quad 0<p\le1.
\]
The scalar $t$ is a threshold; $p$ is a probability level; the indicator
$\mathbf1$ equals one when its condition holds and zero otherwise.
The quantile $q_p$ uses the inverse empirical-CDF convention.
\paragraph{Derivation or example.}
For sorted observations $(1,2,4,8)$, this convention gives
$q_{0.75}=4$. A coordinate budget can use
$a_i=|\delta_{ij}|$, where $\delta_{ij}$ is coordinate $j$ of
observed perturbation $i$. Even if every coordinate has $90\%$
marginal coverage, their joint box coverage need not be $90\%$.
For two coordinates, if their $10\%$ failure events are disjoint,
joint success is $80\%$; if independent, it is $81\%$.
\paragraph{Why ECDR needs it.}
Quantiles provide empirically motivated domain budgets. Actual joint
membership must be measured or bounded separately, and decision safety
still requires a whole-domain proof. A large observed coverage value
does not supply that proof.
\paragraph{Conditions and connections.}
Budgets fitted on training data need a stated fitting population and
evaluation split. Dependence, finite-sample uncertainty, distribution
shift, and quantile interpolation affect interpretation. Quantiles of
raw and normalized coordinates are not interchangeable.
\par\label{endconcept:F34}
\typeout{ECDR-CONCEPT-END F34 PAGE=\thepage}

\concept{F35}{Mathematical induction proves a statement at every step}{\hyperref[concept:F01]{F01}, \hyperref[concept:F02]{F02}, \hyperref[concept:F08]{F08}}
\paragraph{Foundation and definition.}
Induction is a proof principle for statements indexed by the
nonnegative integers. It requires a base case and an implication from
each established case to the next. It can be taken as an axiom of the
natural numbers; equivalently, it follows from their well-ordering
principle, which says every nonempty set of nonnegative integers has a
least member. It is not an empirical extrapolation from several examples.
\paragraph{Symbols and meaning.}
Let $P(q)$ be a proposition indexed by the integer $q\ge0$. The base
case is $P(0)$. The induction step is
$\forall q\ge0:\ P(q)\Rightarrow P(q+1)$. Together they imply
$P(q)$ for every nonnegative integer $q$. For a finite horizon
$Q\ge1$, the step need only hold for integers $0\le q<Q$ to conclude
all cases $0\le q\le Q$.
\paragraph{Derivation or example.}
Suppose a counterexample exists despite the base and step. Its least
index $q_*$ cannot be zero because the base holds. Thus $q_*-1$ is a
nonnegative earlier index, where the claim holds by minimality. The step
then proves $P(q_*)$, a contradiction. For a state example, define
$\mathcal R_0=\{x_0\}$ and
$\mathcal R_{q+1}=\{T(x,u):x\in\mathcal R_q,u\in\mathcal U_{q+1}\}$.
Here states lie in a set $\mathcal X$, inputs lie in a set
$\mathcal U$, $x_0\in\mathcal X$ is the initial state,
$T:\mathcal X\times\mathcal U\to\mathcal X$ is a transition function,
and $\mathcal U_{q+1}\subseteq\mathcal U$ permits next inputs. Membership of each
actual input in its permitted set proves actual-state membership by
induction.
\paragraph{Why ECDR needs it.}
A finite-horizon reuse certificate must retain the effects of every
previous accepted wave. Induction connects trusted initial state and
valid transitions to membership in the certified reachable sets.
\paragraph{Conditions and connections.}
The step must cover every permitted previous state and next input.
Checking only one observed path proves no universal step. After a failed
input gate, continuing requires a new valid base or another proved rule.
\par\label{endconcept:F35}
\typeout{ECDR-CONCEPT-END F35 PAGE=\thepage}

\concept{F36}{Convex sets and convex combinations}{\hyperref[concept:F03]{F03}, \hyperref[concept:F11]{F11}, \hyperref[concept:F16]{F16}}
\paragraph{Foundation and definition.}
A set is convex when the full segment joining any two of its members
also belongs to it. This is a geometric definition, without probabilities
or class labels. A convex combination uses nonnegative coefficients
summing to one; unrestricted real coefficients instead define linear
combinations, and unrestricted coefficients summing to one define affine
combinations.
\paragraph{Symbols and meaning.}
Let $\mathcal S\subseteq\R^d$, where integer $d\ge1$ counts
coordinates. For $x,y\in\mathcal S$ and scalar $t\in[0,1]$,
convexity requires $tx+(1-t)y\in\mathcal S$. More generally, a convex
combination of $r\ge1$ points $x_1,\ldots,x_r$ uses coefficients
$\alpha_i\ge0$ satisfying $\sum_{i=1}^{r}\alpha_i=1$, and is
$\sum_i\alpha_i x_i$. The index $i$ selects a point.
\paragraph{Derivation or example.}
For a box $[\ell,u]\subseteq\R^d$ and members $x,y$, multiply each
coordinate inequality by $t$ or $1-t$ and add:
\[
 \ell_j\le tx_j+(1-t)y_j\le u_j,
 \qquad j=1,\ldots,d.
\]
Thus boxes are convex. Norm balls are convex because the triangle
inequality gives
$\norm{tx+(1-t)y}\le t\norm{x}+(1-t)\norm{y}$ for a ball centered
at zero. In contrast, $\{-1,1\}\subset\R$ excludes the midpoint
zero and is not convex. The convex hull of supplied points is the set
of all their finite convex combinations.
\paragraph{Why ECDR needs it.}
A scalar function $f:\mathcal S\to\R$ is convex when
$f(tx+(1-t)y)\le tf(x)+(1-t)f(y)$. ReLU satisfies this inequality,
so its endpoint chord upper-bounds its graph. Convex relaxations use
such enclosures. A convex input domain alone says nothing about a CNN's
decision throughout that domain.
\paragraph{Conditions and connections.}
Affine images and intersections preserve convexity; arbitrary unions
need not. A union of certified regions can be safe even when nonconvex,
provided each accepted query has a valid regional witness. Replacing a
union by its convex hull requires a new certificate on added points.
\par\label{endconcept:F36}
\typeout{ECDR-CONCEPT-END F36 PAGE=\thepage}

\concept{F37}{Eigenvectors and eigenvalues identify preserved directions}{\hyperref[concept:F12]{F12}, \hyperref[concept:F13]{F13}, \hyperref[concept:F15]{F15}}
\paragraph{Foundation and definition.}
An eigenvector of a square matrix is a nonzero vector whose direction
is preserved by the matrix, although its magnitude or sign may change.
The associated scaling factor is its eigenvalue. The zero vector is
excluded because it satisfies the equation for every scalar and would
identify no direction.
\paragraph{Symbols and meaning.}
Let $A\in\R^{d\times d}$ be a square real matrix, with integer
$d\ge1$. A real eigenpair consists of scalar $\lambda\in\R$ and
nonzero column $v\in\R^d$ satisfying
\[
 Av=\lambda v.
\]
The nonzero scalar multiples of $v$ describe the same direction. A real
matrix need not have real eigenvalues or a complete eigenvector basis.
Complex eigenpairs require extending the scalar field and are outside
this page's real-coordinate treatment.
\paragraph{Derivation or example.}
For $A=\begin{pmatrix}2&1\\1&2\end{pmatrix}$,
$v_1=(1,1)^\top$ satisfies $Av_1=3v_1$, while
$v_2=(1,-1)^\top$ satisfies $Av_2=v_2$. Thus the sum direction is
amplified by three and the difference direction by one.
Equivalently, $(A-\lambda I_d)v=0$, where $I_d$ is the identity
matrix. Thus $A-\lambda I_d$ sends a nonzero vector to zero and
cannot have an inverse: an inverse would send zero back to both zero
and $v$. The eigenvalue identifies precisely this degeneracy.
\paragraph{Why ECDR needs it.}
PCA uses eigenvectors of a covariance matrix to describe directions of
variation. The eigenvalues quantify variance along unit eigenvectors.
Large variance along a direction does not establish that a classifier's
margin is insensitive in that direction.
\paragraph{Conditions and connections.}
Generic square matrices can be defective or have complex eigenvalues.
Real symmetric matrices have the stronger orthonormal-basis guarantee
in \hyperref[concept:F38]{F38}. Repeated eigenvalues can make individual eigenvectors
nonunique, even while their eigenspace is well defined.
\par\label{endconcept:F37}
\typeout{ECDR-CONCEPT-END F37 PAGE=\thepage}

\concept{F38}{Symmetry, positive semidefiniteness, and the spectral theorem}{\hyperref[concept:F18]{F18}, \hyperref[concept:F19]{F19}, \hyperref[concept:F27]{F27}, \hyperref[concept:F37]{F37}}
\paragraph{Foundation and definition.}
A real square matrix is symmetric when transposition leaves it
unchanged. A symmetric matrix is positive semidefinite, abbreviated PSD,
when every quadratic form it produces is nonnegative. PSD does not mean
that every matrix entry is nonnegative.
\paragraph{Symbols and meaning.}
Let $A\in\R^{d\times d}$ satisfy $A^\top=A$, with integer
$d\ge1$. The PSD condition is $x^\top Ax\ge0$ for every
$x\in\R^d$. The finite-dimensional real spectral theorem states
\[
 A=V\Lambda V^\top,\qquad V^\top V=I_d,
 \qquad\Lambda=\diag(\lambda_1,\ldots,\lambda_d).
\]
Columns of $V\in\R^{d\times d}$ form an orthonormal eigenbasis;
$\lambda_i\in\R$ are eigenvalues, and $I_d$ is the identity.
This is an established theorem, rather than an additional axiom.
\paragraph{Derivation or example.}
A proof sketch maximizes $x^\top Ax$ on the compact unit sphere.
Its first-order condition produces a unit eigenvector. Symmetry makes
that vector's orthogonal complement invariant under $A$; induction on
dimension constructs the remaining eigenvectors. This sketches the
argument rather than supplying every optimization detail.
For observations $z_1,\ldots,z_N\in\R^d$ centered to have zero
mean, with integer $N\ge2$, the covariance matrix
$C=(N-1)^{-1}\sum_{i=1}^{N}z_i z_i^\top$ is symmetric and
\[
 x^\top Cx=\frac1{N-1}\sum_{i=1}^{N}(z_i^\top x)^2\ge0.
\]
For a unit eigenvector $v$, its eigenvalue is $v^\top Cv\ge0$.
\paragraph{Why ECDR needs it.}
The theorem justifies orthonormal PCA directions and nonnegative
explained-variance contributions. It also distinguishes a covariance
eigenvalue from a signed classifier weight.
\paragraph{Conditions and connections.}
For symmetric $A$, PSD is equivalent to all eigenvalues being
nonnegative. Zero eigenvalues are allowed and identify directions with
zero quadratic form. Finite-precision eigensolvers need numerical
checks; a tiny negative computed covariance eigenvalue may reflect
rounding rather than a negative mathematical variance.
\par\label{endconcept:F38}
\typeout{ECDR-CONCEPT-END F38 PAGE=\thepage}

\concept{F39}{Singular values connect PCA and operator amplification}{\hyperref[concept:F15]{F15}, \hyperref[concept:F21]{F21}, \hyperref[concept:F38]{F38}}
\paragraph{Foundation and definition.}
The singular value decomposition, or SVD, expresses any finite real
matrix through orthogonal input and output coordinates and nonnegative
axis scalings. Unlike eigenvalue decomposition, it applies to rectangular
matrices. Its existence is an established finite-dimensional theorem;
it follows by applying the spectral theorem to $A^\top A$ and
completing orthonormal bases for null spaces.
\paragraph{Symbols and meaning.}
For $A\in\R^{m\times n}$, with integers $m,n\ge1$, write
\[
 A=U\Sigma V^\top,\qquad U^\top U=I_m,\quad V^\top V=I_n.
\]
Here $U\in\R^{m\times m}$ and $V\in\R^{n\times n}$ are
orthogonal, $I_m,I_n$ are identity matrices, and
$\Sigma\in\R^{m\times n}$ is rectangular diagonal. Its singular
values satisfy $\sigma_1\ge\cdots\ge\sigma_p\ge0$, where
$p=\min(m,n)$. Column $v_i$ of $V$ is a right singular direction.
\paragraph{Derivation or example.}
For any input $x\in\R^n$, put $y=V^\top x$. Orthogonality gives
\[
 \norm{Ax}_2^2=\sum_{i=1}^{p}\sigma_i^2y_i^2
 \le\sigma_1^2\norm{x}_2^2.
\]
Taking $x=v_1$ attains the bound, so the induced Euclidean operator
norm is exactly $\norm{A}_{2\to2}=\sigma_1$.
For centered observations arranged as rows of
$X\in\R^{N\times d}$, with integer $N\ge2$, covariance is
$C=X^\top X/(N-1)$. Substituting $X=U_X\Sigma_XV_X^\top$ shows
that columns of $V_X$ are PCA directions; the nonzero covariance
eigenvalues are the squared singular values divided by $N-1$.
\paragraph{Why ECDR needs it.}
One decomposition connects empirical variation geometry and worst-case
linear amplification, while keeping their meanings distinct: variance
describes fitted observations, whereas an operator norm covers every
input direction.
\paragraph{Conditions and connections.}
Repeated singular values permit nonunique directions. A truncated SVD
describes an approximation, not a universal residual bound. Numerical
estimates of the largest singular value need a certified error allowance
before they serve as upper amplification bounds in a proof.
\par\label{endconcept:F39}
\typeout{ECDR-CONCEPT-END F39 PAGE=\thepage}

\concept{F40}{A derivative is a limit of local change ratios}{\hyperref[concept:F01]{F01}--\hyperref[concept:F07]{F07}: logic and the real-number foundations; \hyperref[concept:F23]{F23}: distance; \hyperref[concept:F26]{F26}: limits and continuity.}
\paragraph{Definition and foundation.}
A derivative quantifies the first-order change of a function near one point. For $f:I\to\R$, with $I\subseteq\R$ an open interval and $a\in I$, define
\[
 f'(a)=\lim_{h\to0}\frac{f(a+h)-f(a)}h
\]
when that finite two-sided limit exists. The increment $h$ is nonzero in the quotient. The derivative is not the value obtained by substituting $h=0$, which would leave an undefined division.
\paragraph{Symbols and dimensions.}
The variable $a$ is the reference input, $h$ is an input increment, and $f'(a)\in\R$ is the candidate slope. Its units are output units divided by input units. The limit means: for every output-slope tolerance $\epsilon>0$, some input radius $\delta>0$ makes the quotient differ from $f'(a)$ by less than $\epsilon$ whenever $0<\abs h<\delta$ and $a+h\in I$.
\paragraph{Derivation and example.}
For $f(x)=x^2$, subtraction and factorization give $(f(a+h)-f(a))/h=2a+h$. Choosing $\delta=\epsilon$ proves the limit $2a$ in dimensionless coordinates. Equivalently,
\[
 f(a+h)=f(a)+f'(a)h+r(h),\qquad
 \frac{\abs{r(h)}}{\abs h}\longrightarrow0.
\]
The remainder $r(h)$ is the error of the first-order approximation; for the quadratic it equals $h^2$.
\paragraph{Why ECDR uses this concept.}
Derivatives describe local sensitivity and underlie training gradients. They help explain why network outputs change, but a slope at the anchor alone does not bound every change in a domain.
\paragraph{Conditions and connections.}
Differentiability implies continuity, while continuity alone does not imply differentiability. ReLU is continuous at zero but its left and right slopes differ, so no ordinary derivative exists there. An implementation's chosen gradient at that point is a declared convention. \hyperref[concept:F41]{F41} extends local linear approximation to vector inputs; \hyperref[concept:C01]{C01} requires bounds valid throughout a set rather than at a single point.
\par\label{endconcept:F40}
\typeout{ECDR-CONCEPT-END F40 PAGE=\thepage}

\concept{F41}{Gradient, Jacobian, and full differentiability}{\hyperref[concept:F40]{F40}; vectors, linear maps, and Euclidean norms.}
\paragraph{Definition and foundation.}
A partial derivative changes one coordinate while holding the others fixed. Full differentiability requires a single linear map to approximate changes in every direction simultaneously. These requirements are different: existence of all coordinatewise derivatives at a point does not by itself establish a valid multivariate linear approximation.
\paragraph{Symbols and dimensions.}
Let $f:O\to\R^m$, where $O\subseteq\R^d$ is open, $a\in O$, and $d,m\ge1$ are integer dimensions. An increment is $h\in\R^d$. The Jacobian $J_f(a)\in\R^{m\times d}$ has entry $[J_f(a)]_{ji}=\partial f_j/\partial x_i(a)$ when $f$ is differentiable. For scalar output $m=1$, the gradient $\nabla f(a)\in\R^d$ is the column transpose of this row Jacobian.
\paragraph{Derivation and counterexample.}
Full differentiability means
\[
 f(a+h)=f(a)+J_f(a)h+r_f(h),\qquad
 \frac{\norm{r_f(h)}_2}{\norm h_2}\longrightarrow0.
\]
For $f(x,y)=x^2+3xy$, the gradient is $(2x+3y,3x)^\top$. In contrast, define $q(x,y)=xy/(x^2+y^2)$ off the origin and $q(0,0)=0$. Both partial derivatives at zero are zero, since $q$ vanishes on either axis. Along $(t,t)$ with $t\ne0$, however, $q(t,t)=1/2$. It is not continuous there and hence not differentiable.
\paragraph{Why ECDR uses this concept.}
CNN backpropagation manipulates Jacobians or their products. A gradient is a local linear coefficient; a certified whole-domain sensitivity bound needs additional control of these derivatives or another sound enclosure.
\paragraph{Conditions and connections.}
Continuous partial derivatives in a neighborhood are a sufficient differentiability condition, though not a necessary one. The $(j,i)$ Jacobian entry has units of output coordinate $j$ divided by input coordinate $i$. Finite differences estimate derivatives under numerical conditions; they do not establish a universal Jacobian bound. \hyperref[concept:F42]{F42} explains how valid local linear maps compose.
\par\label{endconcept:F41}
\typeout{ECDR-CONCEPT-END F41 PAGE=\thepage}

\concept{F42}{The chain rule is composition of local linear maps}{\hyperref[concept:F40]{F40}--\hyperref[concept:F41]{F41}; matrix multiplication and small-remainder notation.}
\paragraph{Definition and foundation.}
The chain rule explains why derivatives multiply through successive operations. Its premise is full differentiability at the relevant points. It is not a rule that any collection of independently existing partial derivatives can be multiplied without checking those premises.
\paragraph{Symbols and dimensions.}
Let $f:O\subseteq\R^d\to\R^m$ be differentiable at $a$, and let $g:V\subseteq\R^m\to\R^k$ be differentiable at $b=f(a)$. The sets $O,V$ are open, and $f$ maps a neighborhood of $a$ into $V$. Write $A=J_f(a)\in\R^{m\times d}$, $B=J_g(b)\in\R^{k\times m}$, and input increment $h\in\R^d$.
\paragraph{Derivation and example.}
Differentiability gives $f(a+h)=b+Ah+r_f(h)$, where $r_f(h)=o(\norm h_2)$. Set $v=Ah+r_f(h)$. A fixed finite matrix has bounded operator norm, so $\norm v_2=O(\norm h_2)$. Substitution into the expansion of $g$ gives
\[
 g(f(a+h))=g(b)+BAh+B r_f(h)+r_g(v),
 \qquad J_{g\circ f}(a)=BA.
\]
Both remainder terms are $o(\norm h_2)$, establishing the derivative. The little-$o$ notation means the remainder norm divided by the indicated increment norm tends to zero. For $f(x)=2x-1$ and $g(z)=z^2$, the derivative is $g'(f(x))f'(x)=4(2x-1)$.
\paragraph{Why ECDR uses this concept.}
For a differentiable scalar loss $L:\R^m\to\R$, reverse-mode differentiation uses $\nabla(L\circ f)(a)=J_f(a)^\top\nabla L(f(a))$. The transpose makes dimensions agree and propagates an output sensitivity back to input coordinates.
\paragraph{Conditions and connections.}
Nondifferentiable ReLU boundaries require an explicit training convention or a suitable generalized analysis. An observed product of local derivatives is still local. \hyperref[concept:C01]{C01} instead bounds each operator over the relevant set to obtain a whole-domain sensitivity bound. \hyperref[concept:F44]{F44} defines the big-$O$ notation used in the remainder argument.
\par\label{endconcept:F42}
\typeout{ECDR-CONCEPT-END F42 PAGE=\thepage}

\concept{F43}{Feasibility, optimality, and a constrained witness}{Sets, inequalities, affine functions, and infima; \hyperref[concept:C07]{C07} and \hyperref[concept:C22]{C22}.}
\paragraph{Definition and foundation.}
A constrained problem specifies an admissible set before asking whether an objective is minimized there. Feasibility asks whether any admissible point exists. Optimality asks which admissible point has the smallest objective. Finding one feasible point answers the first question and generally leaves the second unresolved.
\paragraph{Symbols and dimensions.}
Let $u\in\R^q$ be the decision variable, $G\in\R^{r\times q}$, $h\in\R^r$, $E\in\R^{s\times q}$, and $e\in\R^s$ fixed constraints. The feasible set is $D=\{u:Gu\le h,Eu=e\}$, where inequalities are coordinatewise. A linear program minimizes $a^\top u+\beta$ on $D$, with $a\in\R^q$ and scalar $\beta$; all terms of each constraint must have compatible units.
\paragraph{Derivation and example.}
For $D=[0,1]$ and objective $f(u)=u$, witness $u=0.8$ proves nonemptiness but not optimality; the minimum is zero. More generally, any feasible witness gives an upper bound on a minimization value:
\[
 \inf_{u\in D}f(u)\le f(u_{\rm witness}).
\]
Certification needs a lower bound on every feasible objective value, the opposite inequality direction. A lower affine enclosure or a justified dual certificate can supply that evidence; one sampled objective cannot.
\paragraph{Why ECDR uses this concept.}
Generator-domain membership can be written as feasibility of latent coefficients and residuals. A validated witness suffices to accept membership. Proving positive margin everywhere is a different optimization obligation over all admitted inputs.
\paragraph{Conditions and connections.}
An infeasible set has no admissible queries; an unbounded objective need not have a finite minimum. A nonempty compact feasible set with continuous objective attains its minimum. Numerical solver status alone requires an appropriate validation contract. Failure of a heuristic to find a witness does not prove infeasibility, and feasibility of a solver iterate does not establish its optimality.
\par\label{endconcept:F43}
\typeout{ECDR-CONCEPT-END F43 PAGE=\thepage}

\concept{F44}{Big O, lower bounds, and the chosen cost model}{Sequences, inequalities, and algorithmic input representation; \hyperref[concept:C09]{C09} and \hyperref[concept:C26]{C26}.}
\paragraph{Definition and foundation.}
Asymptotic notation compares growth functions under a specified model of cost and input size. For eventually nonnegative functions $A(n)$ and $B(n)$, write $A(n)=O(B(n))$ if constants $c>0$ and $n_0$ exist such that $A(n)\le cB(n)$ for every integer $n\ge n_0$. Write $A(n)=\Omega(B(n))$ for a corresponding lower inequality and $\Theta(B(n))$ when both hold.
\paragraph{Symbols and dimensions.}
The integer $n\ge1$ measures input size; $A(n)$ may count scalar operations, coordinate reads, or bytes moved, but the choice must be fixed. The comparison function $B(n)$ describes growth. Constants $c,n_0$ are independent of $n$ but may depend on the specified algorithm or fixed architecture. These counts are not automatically elapsed milliseconds.
\paragraph{Derivation and example.}
For $A(n)=3n+7$ and $n\ge1$, $3n\le A(n)\le10n$, proving $A(n)=\Theta(n)$. A box scan that can reject after the first failed coordinate still has worst-case $\Theta(n)$ when every coordinate must be read for an interior query. Its average count depends on the query distribution and coordinate order; no such distribution is encoded by big-$O$.
\paragraph{Why ECDR uses this concept.}
Wave membership may scan $Cs$ coordinates while a full-window box scans $CW$, where $C$ is channels, $s$ appended samples, and $W$ window samples. Dense projection onto $q$ directions typically costs $O(nq)$ scalar arithmetic. These statements identify scaling opportunities, not a measured latency advantage.
\paragraph{Conditions and connections.}
A lower bound for one implementation differs from a lower bound for every admissible algorithm under a model. Hardware parallelism, vector instructions, allocation, memory traffic, and small finite sizes affect runtime constants. Input summaries or trusted state can change the information available and invalidate an old access lower bound's premises. For increments approaching zero, the analogous notation bounds a remainder by a constant times the increment norm near zero, as in \hyperref[concept:F42]{F42}. \hyperref[concept:C26]{C26} counts the actual trace costs before claiming amortized saving.
\par\label{endconcept:F44}
\typeout{ECDR-CONCEPT-END F44 PAGE=\thepage}

\concept{F45}{The exponential function and positive score weights}{Real series and limits; \hyperref[concept:F40]{F40} and ordinary multiplication laws.}
\paragraph{Definition and foundation.}
The exponential can be defined by the convergent real series
\[
 \exp(x)=\sum_{n=0}^{\infty}\frac{x^n}{n!},\qquad x\in\R.
\]
The index $n$ is a nonnegative integer; $0!=1$ and $n!=1\cdot2\cdots n$ for $n\ge1$. Absolute convergence follows from the ratio $\abs x/(n+1)\to0$. The convergence facts used here are results of real analysis, rather than extra axioms about neural networks.
\paragraph{Symbols and dimensions.}
The scalar argument $x$ must be dimensionless in a physically meaningful exponential. The value $\exp(x)$ is dimensionless. For softmax, logits $z_j\in\R$ and temperature $\tau>0$ must make $z_j/\tau$ dimensionless; class index $j$ ranges over $1,\ldots,K$, with integer $K\ge2$.
\paragraph{Derivation and example.}
Absolute convergence permits multiplying and regrouping series, yielding $\exp(a+b)=\exp(a)\exp(b)$ for real $a,b$. Since $\exp(0)=1$, $\exp(x)\exp(-x)=1$, so neither factor is zero. Also $\exp(x)=\exp(x/2)^2>0$. For $t>0$, its series gives $\exp(t)>1$, hence $\exp(x+t)>\exp(x)$: strict monotonicity follows. Termwise differentiation, justified on bounded intervals, gives $\exp'(x)=\exp(x)$. The defining series therefore supports both its algebraic identities and its local sensitivity rule.
\paragraph{Why ECDR uses this concept.}
Softmax defines $p_j=\exp(z_j/\tau)/\sum_{i=1}^{K}\exp(z_i/\tau)$. Positivity makes the denominator positive; a common denominator and strict monotonicity preserve the real score ordering. Thus class-only margin proofs can operate on logits directly.
\paragraph{Conditions and connections.}
This argument assumes finite real logits and positive finite temperature. Mathematical positivity does not prevent floating-point overflow or underflow. Subtracting the same maximum logit before exponentiating preserves exact ratios and improves numerical evaluation, but does not prove bitwise equivalence across implementations. \hyperref[concept:F46]{F46} defines the inverse logarithm and connects it to the training objective.
\par\label{endconcept:F45}
\typeout{ECDR-CONCEPT-END F45 PAGE=\thepage}

\concept{F46}{Natural logarithms, inverse maps, and cross-entropy}{\hyperref[concept:F42]{F42} and \hyperref[concept:F45]{F45}; inverse functions and positive probabilities.}
\paragraph{Definition and foundation.}
The natural logarithm $\log:(0,\infty)\to\R$ is the inverse of the real exponential: $\log(\exp(x))=x$ and $\exp(\log y)=y$ for $x\in\R$, $y>0$. The exponential's continuity, strict monotonicity, and range $(0,\infty)$ justify this inverse. Its range follows from unbounded growth for positive inputs and $\exp(-x)=1/\exp(x)\to0$ as $x\to\infty$.
\paragraph{Symbols and dimensions.}
The variable $y$ is a positive dimensionless argument. For classification, let $z\in\R^K$ be logits, $K\ge2$ class count, $c\in\{1,\ldots,K\}$ the supplied target, and $p_j$ each class's softmax probability at temperature one, including $p_c$ for the target. Define scalar loss $L(z,c)=-\log p_c$. The loss and probability are dimensionless.
\paragraph{Derivation and example.}
The inverse-function derivative and \hyperref[concept:F45]{F45} give $(\log y)'=1/y$. Thus for a differentiable positive scalar function $u(x)$ on an open real interval, the chain rule gives $(\log u(x))'=u'(x)/u(x)$. With $p_c=\exp(z_c)/\sum_i\exp(z_i)$,
\[
 L(z,c)=\log\!\left(\sum_{i=1}^{K}\exp(z_i)\right)-z_c,
 \qquad
 \frac{\partial L}{\partial z_j}=p_j-\mathbf1\{j=c\}.
\]
Here $i$ sums over classes, $j$ is one logit coordinate, and the indicator is one exactly when $j=c$. For target probabilities $0.8$ and $0.2$, the losses are approximately $0.223144$ and $1.609438$, respectively.
\paragraph{Why ECDR uses this concept.}
Cross-entropy explains how training rewards the target relative to all competing logits. Backpropagation uses its local gradient, while certification later establishes a different universal property over an input domain.
\paragraph{Conditions and connections.}
The real logarithm is undefined at zero or negative arguments. Finite real logits give strictly positive softmax probabilities, but numerical underflow can still produce zero in a naive program. A stable log-sum-exp evaluation addresses implementation accuracy; it is not a universal robustness certificate. Reuse of a class alone does not reconstruct current probabilities or this loss.
\par\label{endconcept:F46}
\typeout{ECDR-CONCEPT-END F46 PAGE=\thepage}

\concept{B01}{Sensor samples and a sampling clock}{\hyperref[concept:F04]{F04}, \hyperref[concept:F08]{F08}, \hyperref[concept:F14]{F14}}
\paragraph{Foundation and definition.}
A continuous physical signal and its recorded samples are different objects.
A sampling model chooses discrete times and records a finite vector at each
time. For the six-channel inertial input, write
\[
 v[n]=(a_x[n],a_y[n],a_z[n],\omega_x[n],\omega_y[n],\omega_z[n])^\top,
 \qquad t_n=t_0+n/f_s.
\]
\paragraph{Symbols and meaning.}
$n\in\{0,\ldots,T-1\}$ is a sample index, $T$ is the number of samples,
$v[n]\in\R^6$ is one observation, $a_x,a_y,a_z$ are accelerometer axes,
and $\omega_x,\omega_y,\omega_z$ are angular-velocity axes. $t_0$ is the
first sample time in seconds and $f_s>0$ is samples per second. HAPT's
nominal $f_s=50$ corresponds to $1/f_s=0.02$ seconds per sample. Its raw
acceleration unit is g; angular velocity uses rad/s \cite{hapt}.
\paragraph{Derivation and example.}
For $T=128$, the first-to-last elapsed interval is
$t_{127}-t_0=127/50=2.54$ s. The sample-count convention assigns
$T/f_s=2.56$ s. These are different measurements, not contradictory
estimates. Adding a sampling interval to the last timestamp is a convention,
not another observed sample. Raw recordings without timestamps do not
independently establish a perfectly uniform clock.
\paragraph{Why ECDR needs it.}
A query stride measured in samples becomes a physical query rate only after
choosing this clock model. A receptive field of nine samples is not nine
seconds. Equal sample identity is also the foundation of exact overlap.
\paragraph{Conditions and connections.}
Uniform sampling is a modeling assumption here, not a consequence of the
field axioms. Timing jitter, gaps and asynchronous sensors require a richer
input contract. \hyperref[concept:B02]{B02} defines temporal span; \hyperref[concept:B04]{B04} defines the query grid.
\par\label{endconcept:B01}
\typeout{ECDR-CONCEPT-END B01 PAGE=\thepage}

\concept{B02}{Temporal span: what an annotated interval covers}{\hyperref[concept:F08]{F08}, \hyperref[concept:F24]{F24}, \hyperref[concept:B01]{B01}}
\paragraph{Foundation and definition.}
A temporal annotation span is an interval of sample indices assigned one
activity. It is not the linear span of vectors in \hyperref[concept:F12]{F12}. With inclusive
endpoints $a,b\in\mathbb Z$ and $a\le b$, the index set is
\[
 I_{a,b}=\{n\in\mathbb Z:a\le n\le b\},\qquad |I_{a,b}|=b-a+1.
\]
\paragraph{Symbols and meaning.}
$a$ and $b$ denote the first and last annotated sample indices; $n$ ranges
over samples; $I_{a,b}$ is a finite set; $|I_{a,b}|$ is its cardinality,
not a norm. With sampling rate $f_s>0$, sample-count duration is
$(b-a+1)/f_s$ seconds; first-to-last elapsed time is $(b-a)/f_s$ seconds.
The half-open representation $[a,b+1)$ denotes the same integer samples.
\paragraph{Derivation and example.}
The span $[64,191]$ contains $191-64+1=128$ samples. At 50 Hz it has
sample-count duration $2.56$ s and elapsed time $2.54$ s. A span $[1,128]$
has the same length, but a window grid with starts $0,64,128,\ldots$ does
not contain a length-128 window inside it. Duration and grid alignment
are separate constraints, derived in \hyperref[concept:B05]{B05}.
\paragraph{Why ECDR needs it.}
Annotation support after windowing depends on which windows fit inside a
span. Short transitions may disappear from a pure-window evaluation even
when their annotations remain in the raw data. Understanding the unit of
labeling prevents confusing a class transition with a mixed-label window.
\paragraph{Conditions and connections.}
Index origin and endpoint convention must be verified from source metadata.
The inherited HAPT pipeline reproduces one convention; reproduction does
not independently prove its original annotation semantics. \hyperref[concept:B21]{B21} uses the
word span for a receptive-field extent, another distinct meaning.
\par\label{endconcept:B02}
\typeout{ECDR-CONCEPT-END B02 PAGE=\thepage}

\concept{B03}{A tensor, vectorization and a dimension check}{\hyperref[concept:F14]{F14}, \hyperref[concept:F15]{F15}, \hyperref[concept:F19]{F19}}
\paragraph{Foundation and definition.}
A time-series tensor organizes values by named axes. For one window,
$X\in\R^{C\times W}$ has $C$ channels and $W$ sample positions. A
channel-major vectorization is the bijection
\[
 x_{cW+i}=X[c,i],\qquad
 c=0,\ldots,C-1,\quad i=0,\ldots,W-1.
\]
\paragraph{Symbols and meaning.}
$X[c,i]$ is one scalar observation; $x=\operatorname{vec}(X)\in\R^d$ is
the same collection in a vector; $d=CW$ is input dimension. $c$ indexes
channels and $i$ indexes within-window time. A batch adds an axis:
$X^{\mathrm{batch}}\in\R^{B\times C\times W}$ with $B$ windows. Batch
index is not a new physical sensor or a new time sample.
\paragraph{Derivation and example.}
For $C=2,W=3$, rows $(1,2,3)$ and $(4,5,6)$ become
$x=(1,2,3,4,5,6)^\top$. Another fixed layout is a permutation $Px$,
where $P\in\R^{d\times d}$ has one 1 per row and column. Since
$P^\top P=I_d$, $\norm{Px}_2=\norm{x}_2$. Here $I_d$ is the identity
matrix. However, unchanged CNN weights generally do not compensate for an
arbitrary input permutation.
\paragraph{Why ECDR needs it.}
The classifier, domain generator, normalization and membership checker must
agree on the layout. For HAPT, $C=6,W=128$ gives $d=768$; the final feature
dimension may be 64, a different object. A float32 input has
$768\times4=3072$ payload bytes, not total runtime memory.
\paragraph{Conditions and connections.}
Vectorization is lossless only when its inverse and axis convention are
fixed. Unit consistency is addressed by \hyperref[concept:B07]{B07}--\hyperref[concept:B08]{B08}; actual CNN operators
consume the structured tensor rather than an unlabeled coordinate list.
\par\label{endconcept:B03}
\typeout{ECDR-CONCEPT-END B03 PAGE=\thepage}

\concept{B04}{The sliding-window grid and its exact count}{\hyperref[concept:F08]{F08}, \hyperref[concept:F09]{F09}, \hyperref[concept:B01]{B01}}
\paragraph{Foundation and definition.}
A no-padding query grid takes complete windows from one recording. For
$T$ samples, window length $W\ge1$ and stride $s\ge1$, starts are $ks$,
where $k$ is a nonnegative integer. A valid window reads
\[
 X_k[:,i]=v[ks+i],\qquad i=0,\ldots,W-1,
 \qquad ks+W\le T.
\]
\paragraph{Symbols and meaning.}
$v[n]\in\R^C$ is sample $n$, $C$ is channel count, $X_k\in\R^{C\times W}$
is query $k$, and ``:'' means every channel. $T,W,s,k,i$ count samples
or windows, not seconds. No window crosses a recording boundary.
\paragraph{Derivation and example.}
The final inequality is equivalent to $k\le(T-W)/s$. When $T\ge W$,
\[
 N_T=1+\left\lfloor\frac{T-W}{s}\right\rfloor;
 \qquad N_T=0\text{ if }T<W.
\]
There are $N_T$ integers because index 0 is included. At
$T=1000,W=128,s=64$, $N_T=14$. The last start is $13\cdot64=832$,
ending at exclusive index 960; the next would end at 1024 and is invalid.
The nominal query rate is $f_s/s=50/64=0.78125$ queries/s within a long
uniformly sampled recording, ignoring startup and termination.
\paragraph{Why ECDR needs it.}
Query count sets evaluation denominators and amortization frequency.
Overlapping windows do not multiply the independent recording duration.
For multiple recordings, sum their individual counts; there is no valid
pair from the end of one recording to the beginning of another.
\paragraph{Conditions and connections.}
This count assumes fixed length, fixed stride, no padding and a recording
grid starting at index 0. Different alignment policies require a different
count. \hyperref[concept:B05]{B05} applies the same integer reasoning to an annotation span.
\par\label{endconcept:B04}
\typeout{ECDR-CONCEPT-END B04 PAGE=\thepage}

\concept{B05}{How many grid windows fit inside a label span?}{\hyperref[concept:F08]{F08}, \hyperref[concept:B02]{B02}, \hyperref[concept:B04]{B04}}
\paragraph{Foundation and definition.}
A pure window must fit wholly inside one consistently labeled interval.
Use a half-open span $[a,b)$ inside a recording with indices $[0,T)$.
The grid starts remain $ks$; they are not reset at each annotation.
\paragraph{Symbols and meaning.}
$a,b,T\in\mathbb Z$ are span and recording bounds, $W\ge1$ is window
length, $s\ge1$ is stride, and $k\ge0$ is a window index. Define integer
bounds $k_{\min}$ and $k_{\max}$ by
\[
 k_{\min}=\left\lceil\frac{\max(a,0)}s\right\rceil,\qquad
 k_{\max}=\left\lfloor\frac{\min(b-W,T-W)}s\right\rfloor.
\]
\paragraph{Derivation and example.}
The required inequalities are $ks\ge a$, $ks+W\le b$, $ks\ge0$ and
$ks+W\le T$. Dividing by positive $s$, then rounding lower bounds upward
and upper bounds downward, yields
\[
 N_{[a,b)}=\max(0,k_{\max}-k_{\min}+1).
\]
At $a=100,b=400,W=128,s=64,T\ge400$, the bounds are 2 and 4, so starts
128,192,256 are valid. Start 320 ends at 448 and is outside the span.
For inclusive input $[a,b_{\rm inc}]$, substitute $b=b_{\rm inc}+1$.
\paragraph{Why ECDR needs it.}
A span with at least $W$ samples can still produce zero pure windows if
the fixed grid misses it. Thus pure class support is a property of both
the annotated signal and the query policy. It does not directly estimate
the prevalence of that activity in the raw recording or in deployment.
\paragraph{Conditions and connections.}
The formula treats one span and one recording. Adjacent spans with the
same activity require an explicit merge policy before adding counts.
Changing the index convention without validation changes the population
being analyzed; an algebraically correct formula cannot validate that choice.
\par\label{endconcept:B05}
\typeout{ECDR-CONCEPT-END B05 PAGE=\thepage}

\concept{B06}{A label policy is a definition of the evaluation population}{\hyperref[concept:F03]{F03}, \hyperref[concept:F08]{F08}, \hyperref[concept:B05]{B05}}
\paragraph{Foundation and definition.}
Window labels summarize many sample-level labels. For $K$ activity classes,
let $n_j$ count samples assigned class $j\in\{1,\ldots,K\}$ and let $n_0$
count unannotated samples. Then
\[
 \sum_{j=0}^{K}n_j=W,\qquad
 p_{\max}=\max_{1\le j\le K}n_j/W.
\]
\paragraph{Symbols and meaning.}
$W$ is the number of samples in the window; each $n_j$ is an integer
between 0 and $W$. $p_{\max}$ is a whole-window proportion, not a CNN
probability. ID0 denotes absent annotation here, not an additional
recognized activity. An inherited policy is pure if some $n_j=W$;
unlabeled if $n_0=W$; otherwise majority when $p_{\max}\ge0.6$, and
ambiguous when $p_{\max}<0.6$.
\paragraph{Derivation and example.}
Because counts are integer, a majority threshold of 0.6 at $W=128$ needs
$n_j\ge\lceil0.6\cdot128\rceil=77$. If 70 samples have one activity and
58 are unannotated, the observed-label proportion is $70/70=1$, but the
policy proportion is $70/128=0.546875$: the window is ambiguous. The
choice of denominator changes the definition, not just its display.
\paragraph{Why ECDR needs it.}
Filtering to pure windows changes which transitions and mixed boundaries
remain in a study. A reference can be trained on a subset while a runtime
stream still contains every query. Reuse coverage needs its own explicitly
stated population. True activity labels are available for retrospective
analysis, not automatically to a runtime event selector.
\paragraph{Conditions and connections.}
Tie handling for a majority ID and annotation conventions must be fixed.
The policy is a modeling choice; probability axioms do not prescribe it.
\hyperref[concept:B30]{B30} and \hyperref[concept:C28]{C28} explain why metrics with different denominators cannot be renamed.
\par\label{endconcept:B06}
\typeout{ECDR-CONCEPT-END B06 PAGE=\thepage}

\concept{B07}{Frozen normalization as an invertible affine map}{\hyperref[concept:F15]{F15}--\hyperref[concept:F17]{F17}, \hyperref[concept:B03]{B03}}
\paragraph{Foundation and definition.}
Normalization changes coordinates rather than changing the underlying
recorded observation. In vector form, fixed per-coordinate centering and
positive scaling give
\[
 z=\mathcal P(x)=D_\sigma^{-1}(x-\mu),\qquad
 D_\sigma=\operatorname{diag}(\sigma_1,\ldots,\sigma_d).
\]
\paragraph{Symbols and meaning.}
$x\in\R^d$ is raw input, $z\in\R^d$ is normalized input,
$\mu\in\R^d$ contains fitted means, and every $\sigma_i>0$ is a fixed
scale. $D_\sigma\in\R^{d\times d}$ is diagonal. For per-channel fitting,
the same mean and scale repeat across positions of that channel.
\paragraph{Derivation and example.}
Invertibility follows directly: $x=D_\sigma z+\mu$. If one coordinate has
$\mu_i=2,\sigma_i=0.5$ and raw value $x_i=2.1$, then $z_i=0.2$.
For anchor $x^r$ and its image $z^r=\mathcal P(x^r)$,
\[
 z-z^r=D_\sigma^{-1}(x-x^r).
\]
Centering cancels in a difference, but scaling remains. Training-window
means may weight overlapping raw samples repeatedly; they need not equal
means computed on unique raw samples. The fit population must be recorded.
\paragraph{Why ECDR needs it.}
The reference is the composition of preprocessing and CNN, not the CNN
weights alone. A certificate in normalized coordinates must be checked in
those coordinates or transported correctly to raw space. Replacing fitted
scales changes the function being certified.
\paragraph{Conditions and connections.}
Positive finite scales, finite inputs and a frozen fit are assumptions.
Per-window centering defines a different map and can destroy exact
normalized overlap. This operation does not remove gravity or independently
calibrate a sensor. \hyperref[concept:B08]{B08} derives the corresponding transport of bounds.
\par\label{endconcept:B07}
\typeout{ECDR-CONCEPT-END B07 PAGE=\thepage}

\concept{B08}{Transporting a metric, a region and an affine bound}{\hyperref[concept:F16]{F16}, \hyperref[concept:F20]{F20}, \hyperref[concept:B07]{B07}}
\paragraph{Foundation and definition.}
Coordinate transport applies the same invertible map to every object used
in a claim. With $z=D^{-1}(x-\mu)$ and fixed positive diagonal $D$,
\[
 \norm{z-z^r}_2^2=\sum_{i=1}^{d}\frac{(x_i-x_i^r)^2}{\sigma_i^2}.
\]
\paragraph{Symbols and meaning.}
$d$ is input dimension; $\mu\in\R^d$ is the fixed fitted mean;
$x,x^r\in\R^d$ are raw query and anchor;
$z,z^r$ are their normalized images; $\sigma_i=D_{ii}>0$ is a scale.
The normalized norm is dimensionless if scales carry the corresponding
raw units. Raw channels in g and rad/s have no shared physical unit merely
because their numeric values are added in a norm.
\paragraph{Derivation and example.}
A raw box with budgets $b_i^{\rm raw}\ge0$ and $|x_i-x_i^r|\le b_i^{\rm raw}$ maps to
$|z_i-z_i^r|\le b_i^{\rm raw}/\sigma_i$. If normalized margin
$\gamma$ has a valid bound $\gamma(z)\ge a_z^\top z+\beta_z$, substitution gives
\[
 a_x=D^{-\top}a_z,\qquad
 \beta_x=\beta_z-a_z^\top D^{-1}\mu,
 \qquad g(x)=\gamma(\mathcal P(x))\ge a_x^\top x+\beta_x.
\]
$a_z,a_x\in\R^d$ are bound coefficients and $\beta_z,\beta_x\in\R$ are
intercepts. For scalars $a_z=3,\beta_z=1,\mu=2,\sigma=0.5$, obtain
$a_x=6,\beta_x=-11$; both expressions equal 1.6 at raw $x=2.1$.
\paragraph{Why ECDR needs it.}
Using a normalized radius against raw distances accepts the wrong set.
Transport is the link between a mathematical domain and the physical
observations admitted by a checker. Changing units can change a raw
numeric distance while representing the same physical variation.
\paragraph{Conditions and connections.}
The lower bound must hold on the image of the raw domain. Invertibility
alone does not make a bound globally valid. \hyperref[concept:C14]{C14} applies the same reasoning
to projected and generator-based regions.
\par\label{endconcept:B08}
\typeout{ECDR-CONCEPT-END B08 PAGE=\thepage}

\concept{B09}{The shift--append operation is a linear state update}{\hyperref[concept:F15]{F15}, \hyperref[concept:F17]{F17}, \hyperref[concept:B03]{B03}}
\paragraph{Foundation and definition.}
A window advances by dropping its first $s$ columns and appending a new
wave. Define
\[
 \mathcal T_s(X,U)=S_sX+J_sU=[X[:,s:W],U].
\]
\paragraph{Symbols and meaning.}
$X\in\R^{C\times W}$ is the old window, $U\in\R^{C\times s}$ is new
data, and $1\le s\le W$. $S_s$ means drop-and-zero-fill:
$S_sX=[X[:,s:W],0_{C\times s}]$.
$J_sU=[0_{C\times(W-s)},U]$ inserts the new tail. They are linear maps
between the stated tensor spaces; $0_{a\times b}$ is an all-zero matrix.
The brackets concatenate columns, not add them.
\paragraph{Derivation and example.}
With $C=1,W=4,s=2$, vectorized matrices are
\[
 S_s=\begin{pmatrix}0&0&1&0\\0&0&0&1\\0&0&0&0\\0&0&0&0\end{pmatrix},
 \qquad J_s=\begin{pmatrix}0&0\\0&0\\1&0\\0&1\end{pmatrix}.
\]
For $X=(a,b,c,d)^\top$ and $U=(u,v)^\top$, their sum is
$(c,d,u,v)^\top$. $J_s$ is injective: reading its last $s$ positions
recovers $U$, so two different tails cannot create the same full update
when the old window is fixed.
\paragraph{Why ECDR needs it.}
This operation exposes exactly which variables are fixed and which are
new. A one-step wave region can vary $Cs$ coordinates instead of all
$CW$ coordinates. It also defines a state transition for multi-step proofs.
\paragraph{Conditions and connections.}
The update assumes the same recording, channel order and sample grid.
Dropped samples are lost from this state; $S_s$ is generally not invertible.
Reducing free input dimension does not by itself prove lower verifier cost.
\hyperref[concept:B10]{B10} identifies the equality inherited by shared samples.
\par\label{endconcept:B09}
\typeout{ECDR-CONCEPT-END B09 PAGE=\thepage}

\concept{B10}{Aligned overlap is an identity of sample indices}{\hyperref[concept:F09]{F09}, \hyperref[concept:B04]{B04}, \hyperref[concept:B09]{B09}}
\paragraph{Foundation and definition.}
Two consecutive windows from the same fixed grid share samples. For
recording $v[n]\in\R^C$ and $X_k[:,i]=v[ks+i]$,
\[
 X_{k+1}[:,i]=X_k[:,i+s],\qquad 0\le i<W-s.
\]
\paragraph{Symbols and meaning.}
$k$ is window index, $i$ is a within-window sample index, $s$ is stride,
$W$ is window length and $C$ is channel count, with $1\le s\le W$ here. The same vector appears at
different tensor positions: new position $i$ and old position $i+s$.
There are $C(W-s)$ shared scalar coordinates after alignment.
\paragraph{Derivation and example.}
Substitute the window definition:
$X_{k+1}[:,i]=v[(k+1)s+i]=v[ks+(i+s)]$.
The final expression is the definition of $X_k[:,i+s]$. This proof uses
index equality, not smoothness, a common activity label or high correlation.
For $C=6,W=128,s=64$, 384 shared values remain and 384 new values arrive.
Time-homogeneous per-channel normalization preserves equality: the same
channel map is applied at every temporal position. Separately fitted window means need
not preserve it.
\paragraph{Why ECDR needs it.}
Exact feature sharing for a valid CNN starts from this identity. It can
also justify checking only a new tail if the old overlap is guaranteed by
the acquisition contract. A checker without that guarantee must verify the
overlap too. Data similarity alone does not supply sample identity.
\paragraph{Conditions and connections.}
There is no such guarantee across recording boundaries, dropped samples,
changed channel layouts or asynchronous resampling. \hyperref[concept:B11]{B11} shows why this
identity does not imply a small ordinary difference; \hyperref[concept:B24]{B24} lifts it through
appropriate convolutional layers.
\par\label{endconcept:B10}
\typeout{ECDR-CONCEPT-END B10 PAGE=\thepage}

\concept{B11}{Same-coordinate perturbation is not aligned overlap}{\hyperref[concept:F20]{F20}, \hyperref[concept:B09]{B09}--\hyperref[concept:B10]{B10}}
\paragraph{Foundation and definition.}
An ordinary difference compares matching coordinate labels in two input
tensors, not the matching raw samples. For old $X$ and new
$X'=\mathcal T_s(X,U)$,
\[
 \Delta=X'-X=(S_s-I)X+J_sU.
\]
\paragraph{Symbols and meaning.}
$X,X',\Delta\in\R^{C\times W}$; $U\in\R^{C\times s}$ is new data;
$S_s,J_s$ are \hyperref[concept:B09]{B09}'s linear maps; $I$ is the identity on full windows.
The Frobenius norm $\norm{\Delta}_F$ is Euclidean norm after vectorization.
It aggregates differences at the same channel and window position.
\paragraph{Derivation and counterexample.}
For one channel, $W=4,s=2$, take $X=(0,0,M,M)$ and $U=(0,0)$, where
$M\in\R$ is arbitrary amplitude. Then $X'=(M,M,0,0)$ still satisfies exact
50\% aligned overlap, but
\[
 \norm{X'-X}_2=\sqrt{M^2+M^2+M^2+M^2}=2|M|.
\]
No small upper bound independent of amplitude follows from overlap. If an
additional assumption imposes $|X[c,i]|,|U[c,j]|\le A$ for finite $A\ge0$,
each difference is at most $2A$, giving
$\norm{\Delta}_F\le2A\sqrt{CW}$. This bound comes from amplitude control.
\paragraph{Why ECDR needs it.}
An ordinary norm ball may need to include changes due to moving old samples
as well as new data. Reusing the last class merely because half the samples
are shared is unjustified. Wave geometry can encode this structure directly.
\paragraph{Conditions and connections.}
The example disproves a logical implication, not an empirical claim about
human movement. Bounds on input distance still require margin sensitivity
to obtain a decision certificate, as in \hyperref[concept:C02]{C02}.
\par\label{endconcept:B11}
\typeout{ECDR-CONCEPT-END B11 PAGE=\thepage}

\concept{B12}{Marginal coverage, joint coverage and the union bound}{\hyperref[concept:F29]{F29}--\hyperref[concept:F34]{F34}, \hyperref[concept:F24]{F24}}
\paragraph{Foundation and definition.}
A box accepts a query only when every coordinate constraint holds. Let
$\Delta=(\Delta_1,\ldots,\Delta_d)$ be a random perturbation and define
$E_i=\{|\Delta_i|\le b_i\}$ for fixed budgets $b_i\ge0$.
Box membership is the joint event $\bigcap_{i=1}^dE_i$.
\paragraph{Symbols and meaning.}
$d\ge1$ is dimension, $i$ indexes coordinates, $\Pr$ is probability under
a specified data-generating distribution, and $\alpha_i\in[0,1]$ is a
valid upper bound on $\Pr(E_i^c)$. The complement $E_i^c$ is coordinate
failure. A train-set quantile does not automatically provide this
population probability guarantee.
\paragraph{Derivation and example.}
Countable subadditivity from the probability axioms gives
\[
 \Pr\!\left(\bigcap_iE_i\right)
 =1-\Pr\!\left(\bigcup_iE_i^c\right)
 \ge\max\left(0,1-\sum_i\alpha_i\right).
\]
No independence is used. If, additionally, the events are independent and
each has probability $q$, their joint probability is $q^d$.
For $q=0.9,d=768$, it is approximately $7.21518\times10^{-36}$ under
that added assumption. This is a constructed example, not a HAPT estimate.
With two identical events of probability 0.9, joint coverage is instead 0.9.
\paragraph{Why ECDR needs it.}
Per-coordinate q90 budgets do not establish q90 joint membership, much
less q90 certified reuse. Geometry, empirical joint coverage and a
whole-domain model property are three separate questions.
\paragraph{Conditions and connections.}
The distribution and fit/evaluation split must be declared before making
statistical claims. A certificate can be distribution-free conditional on
membership; its usefulness still depends on how often the stream enters
the domain. \hyperref[concept:C18]{C18} and \hyperref[concept:C28]{C28} separate those claims.
\par\label{endconcept:B12}
\typeout{ECDR-CONCEPT-END B12 PAGE=\thepage}

\concept{B13}{Temporal dependence and the unit of a split}{\hyperref[concept:F29]{F29}--\hyperref[concept:F33]{F33}, \hyperref[concept:B10]{B10}}
\paragraph{Foundation and definition.}
Independence means a joint probability factors appropriately; it is not a
consequence of storing observations in different tensor rows. Overlapping
windows contain common raw samples. Multiple recordings from one person
can also share person-specific behavior and device placement.
\paragraph{Symbols and meaning.}
Let $Y_1,\ldots,Y_N$ be scalar evaluation outcomes with finite second
moments, $N\ge1$ their count, and $\bar Y=N^{-1}\sum_iY_i$ their mean.
$\operatorname{Cov}(Y_i,Y_j)$ is covariance. Then
\[
 \operatorname{Var}(\bar Y)=\frac1{N^2}
 \sum_{i=1}^{N}\sum_{j=1}^{N}\operatorname{Cov}(Y_i,Y_j).
\]
\paragraph{Derivation and example.}
Expand the centered square of the sum and use expectation linearity.
The familiar $\sigma^2/N$ requires zero off-diagonal covariances and
appropriate common variance $\sigma^2$. If $Y_i=Y_1$ for every $i$,
averaging duplicate rows gives variance $\operatorname{Var}(Y_1)$, not
$\operatorname{Var}(Y_1)/N$. This extreme example explains why row count
alone does not measure independent information.
\paragraph{Why ECDR needs it.}
A random window split may put the same raw sample into training and test.
Splitting by subject avoids that particular overlap but does not make all
within-subject queries independent. Reuse should be replayed in recording
order because its anchor state depends on preceding queries.
\paragraph{Conditions and connections.}
The statistical unit, clustering and dependence assumptions must match the
uncertainty claim. A subject-level bootstrap is a possible design, not an
already-performed experiment here. Describing all splits for exploration
must not be called an untouched confirmatory test afterward.
\par\label{endconcept:B13}
\typeout{ECDR-CONCEPT-END B13 PAGE=\thepage}

\concept{B14}{The reference classifier is a composition, not just weights}{\hyperref[concept:F04]{F04}, \hyperref[concept:F17]{F17}, \hyperref[concept:B07]{B07}}
\paragraph{Foundation and definition.}
The reference function maps a raw query into class scores through a fixed
preprocessing map followed by a fixed network:
\[
 F=f_\theta\circ\mathcal P:\R^d\to\R^K,
 \qquad F(x)=f_\theta(\mathcal P(x)).
\]
\paragraph{Symbols and meaning.}
$d=CW$ is raw input dimension, $K\ge2$ is class count,
$x\in\R^d$ is one vectorized raw window, $\mathcal P$ is preprocessing,
and $f_\theta$ is the CNN with parameters and frozen operator state
denoted collectively by $\theta$. Its input is normalized vector $z$;
$F$'s input is raw vector $x$. Class ordering and tie handling belong to
the decision contract even though neither appears as a trainable weight.
\paragraph{Derivation and example.}
Let scalar preprocessing be $\mathcal P(x)=2x-4$ and an illustrative
two-score network be $f_\theta(z)=(z,-z)$. Then
$F(x)=(2x-4,4-2x)$. The winner changes at raw value $x=2$, while the
network boundary is normalized value $z=0$. Substituting raw $x$ directly
into the network uses a different classifier.
\paragraph{Why ECDR needs it.}
A reusable certificate must match the graph, weights, preprocessing,
class schema, deterministic mode and domain coordinates used at runtime.
A matching filename is not evidence of matching function identity.
Certification relative to a reference means preserving its decision,
including occasions when that reference disagrees with ground truth.
\paragraph{Conditions and connections.}
Real-valued functional identity and floating-point program identity are
distinct obligations. Updating weights or fitted normalization requires
new evidence unless an explicit transfer result applies. \hyperref[concept:C20]{C20} handles
numerical semantics; \hyperref[concept:C23]{C23} handles anchor identity through time.
\par\label{endconcept:B14}
\typeout{ECDR-CONCEPT-END B14 PAGE=\thepage}

\concept{B15}{Logits, maximum, argmax and deterministic ties}{\hyperref[concept:F06]{F06}, \hyperref[concept:F45]{F45}, \hyperref[concept:B14]{B14}}
\paragraph{Foundation and definition.}
A logit vector is a list of real scores. For $F(x)\in\R^K$, the classifier
returns an index with largest score:
\[
 h(x)=\operatorname{argmax}_{1\le j\le K}F_j(x).
\]
\paragraph{Symbols and meaning.}
$x\in\R^d$ is a query, $F_j(x)\in\R$ is the score for class index $j$,
$K\ge2$ is class count and $h(x)\in\{1,\ldots,K\}$ is the selected
class. A deterministic tie rule makes $h$ a function rather than a
set-valued relation. In contrast, $\max_jF_j(x)$ returns a score value,
not the class index. Logits do not have sensor units or automatic
probabilistic calibration.
\paragraph{Derivation and example.}
For scores $(2.4,1.1,2.0)$, the maximum is 2.4 and the unique argmax is
class~1. For $(2,2,0)$, the set of maximizers is $\{1,2\}$; a
smallest-index rule returns1. Softmax at temperature $\tau>0$ gives
$p_j=\exp(F_j/\tau)/\sum_k\exp(F_k/\tau)$, where $k$ is the denominator
index. Its positive common denominator and the strict monotonicity of
exponentiation preserve score ordering in exact real arithmetic.
\paragraph{Why ECDR needs it.}
Class-only reuse need not recreate current logits or probabilities. A
downstream consumer requiring an embedding or calibrated confidence asks
for a different output contract. Margin certification can therefore work
directly on logits without inserting softmax.
\paragraph{Conditions and connections.}
Strict positive pairwise margins avoid relying on ties, but are sufficient
rather than necessary for a particular deterministic tie rule. NaN scores
are outside the finite-real model. \hyperref[concept:B16]{B16} formalizes the pairwise test;
\hyperref[concept:C20]{C20} accounts for numerical errors in actual programs.
\par\label{endconcept:B15}
\typeout{ECDR-CONCEPT-END B15 PAGE=\thepage}

\concept{B16}{A pairwise decision margin converts argmax into inequalities}{\hyperref[concept:F06]{F06}, \hyperref[concept:B15]{B15}}
\paragraph{Foundation and definition.}
For a fixed target class $c$, compare its score with each competing class.
The pairwise margin is
\[
 g_{cj}(x)=F_c(x)-F_j(x),\qquad j\in\{1,\ldots,K\}\setminus\{c\}.
\]
\paragraph{Symbols and meaning.}
$x\in\R^d$ is raw input; $F:\R^d\to\R^K$ is the frozen reference;
$c$ is a chosen target index and $j$ is a different competitor. Each
$g_{cj}(x)\in\R$ is a scalar score difference. For $K=12$, there are
11 separate comparisons. Define the smallest target margin at an anchor
$x^r$ as $m_r=\min_{j\ne c}g_{cj}(x^r)$.
\paragraph{Derivation and example.}
If every $g_{cj}(x)>0$, then $F_c(x)>F_j(x)$ for every competitor, so
$h(x)=c$ uniquely. For $(F_1,F_2,F_3)=(4.5,-0.5,1)$ and $c=1$,
the two margins are 5 and 3.5. At another input with scores
$(6.5,-2.5,5)$, target score increases but the smaller margin falls to
1.5. Target confidence cannot be inferred from its own logit alone.
\paragraph{Why ECDR needs it.}
Reuse requires preserving ordering, not preserving every score numerically.
Bounding a difference directly can retain cancellations between two logits
that are lost when each logit is bounded independently. A verifier must
check all competitors, not just the class that is second-largest at the
anchor: a different competitor can become dangerous elsewhere in a region.
\paragraph{Conditions and connections.}
Strict positivity is sufficient; equality depends on the tie rule.
A positive margin at one input is a pointwise fact, not a whole-domain
certificate. \hyperref[concept:B17]{B17} supplies a local existence theorem from continuity;
\hyperref[concept:C18]{C18} states the universal lower-bound condition used by ECDR.
\par\label{endconcept:B16}
\typeout{ECDR-CONCEPT-END B16 PAGE=\thepage}

\concept{B17}{Local invariance from continuity and a positive anchor margin}{\hyperref[concept:F26]{F26}, \hyperref[concept:B16]{B16}}
\paragraph{Foundation and definition.}
A classifier can be locally constant in its decision even though its
logits vary continuously. Suppose $F:\R^d\to\R^K$ is continuous at
anchor $x^r$, $K$ is finite, and the chosen class $c$ has
$g_{cj}(x^r)>0$ for every $j\ne c$.
\paragraph{Symbols and meaning.}
$g_{cj}=F_c-F_j$ is a pairwise margin; $d$ is input dimension;
$\norm{\cdot}$ is a fixed norm; and $\epsilon_j>0$ will denote an
input radius for competitor $j$. A ball of radius $\epsilon$ contains
all $x$ satisfying $\norm{x-x^r}<\epsilon$.
\paragraph{Derivation.}
Each margin is continuous because it is a difference of continuous
functions. For each $j$, continuity provides $\epsilon_j>0$ such that
\[
 \norm{x-x^r}<\epsilon_j
 \ \Longrightarrow\
 |g_{cj}(x)-g_{cj}(x^r)|<g_{cj}(x^r)/2.
\]
Consequently $g_{cj}(x)>g_{cj}(x^r)/2>0$. Take
$\epsilon=\min_{j\ne c}\epsilon_j$. A minimum of finitely many positive
numbers is positive. Thus every input in that ball has unique winner $c$.
\paragraph{Why ECDR needs it.}
The theorem explains why a certified invariant neighborhood is possible
around a strictly classified input. It does not compute its useful size.
Operational bounds such as Lipschitz or affine relaxations turn this
existence statement into checkable sufficient conditions.
\paragraph{Conditions and connections.}
Finite $K$, continuity at the anchor and strict margins are necessary
assumptions for this proof. It says neither that the radius is independent
of the model nor that a future query enters the ball. \hyperref[concept:B18]{B18} shows the former
failure; \hyperref[concept:C02]{C02} constructs a quantitative radius under sensitivity bounds.
\par\label{endconcept:B17}
\typeout{ECDR-CONCEPT-END B17 PAGE=\thepage}

\concept{B18}{Why no positive radius can be universal over all classifiers}{\hyperref[concept:F02]{F02}, \hyperref[concept:B16]{B16}--\hyperref[concept:B17]{B17}}
\paragraph{Foundation and definition.}
To disprove a universal claim, construct one permitted counterexample for
each proposed bound. Consider the claim that some fixed radius $r_0>0$
preserves the decision of every continuous classifier at every strictly
classified anchor. The claim reverses the order of quantifiers in \hyperref[concept:B17]{B17}.
\paragraph{Symbols and meaning.}
$r_0\in\R$ is a proposed positive universal radius. Define
$\varepsilon=r_0/2$, scalar input $x\in\R$, anchor $x^r=0$, and a
two-class reference
\[
 F^{(\varepsilon)}(x)=(\varepsilon-x,0).
\]
The superscript labels the chosen model, not exponentiation. Its target
class is $c=1$ at the anchor.
\paragraph{Derivation and counterexample.}
The margin is $g_{12}(x)=\varepsilon-x$, a continuous affine function
with Lipschitz constant~1. At the anchor its margin is $\varepsilon>0$.
Take $x'=3r_0/4$. Then $|x'-x^r|=3r_0/4<r_0$, but
\[
 g_{12}(x')=r_0/2-3r_0/4=-r_0/4<0.
\]
Class2 wins within the proposed radius. This construction works for every
$r_0>0$. It does not contradict a model-dependent radius smaller than
$\varepsilon$.
\paragraph{Why ECDR needs it.}
Input overlap, network continuity or a domain's descriptive name cannot
justify a fixed reusable radius for arbitrary weights and margins.
The region must be linked to a particular reference, anchor and bound.
\paragraph{Conditions and connections.}
This is a distribution-free logical counterexample, not a claim that HAPT
models behave this way. Even a known Lipschitz constant is insufficient
without an anchor margin. \hyperref[concept:C02]{C02} makes that dependence explicit.
\par\label{endconcept:B18}
\typeout{ECDR-CONCEPT-END B18 PAGE=\thepage}

\concept{B19}{Convolution as a sparse affine map with shared coefficients}{\hyperref[concept:F09]{F09}, \hyperref[concept:F15]{F15}--\hyperref[concept:F16]{F16}, \hyperref[concept:B03]{B03}}
\paragraph{Foundation and definition.}
The Conv1D convention used here is cross-correlation: kernel coefficients
are not reversed. For valid padding and stride1,
\[
 V[o,i]=b_o+\sum_{c=0}^{C_{\rm in}-1}\sum_{r=0}^{k-1}
 W[o,c,r]Z[c,i+r].
\]
\paragraph{Symbols and meaning.}
$Z\in\R^{C_{\rm in}\times T}$ is input, $V\in\R^{C_{\rm out}\times(T-k+1)}$
is preactivation, $W\in\R^{C_{\rm out}\times C_{\rm in}\times k}$ is
the kernel and $b\in\R^{C_{\rm out}}$ is bias. $o,c,r,i$ index output
channel, input channel, kernel tap and output position. $1\le k\le T$.
The same $W[o,c,r]$ is used at all output positions.
\paragraph{Derivation and example.}
For one channel with $Z=(1,2,3,4)$, kernel $(1,-1)$ and bias~0.5,
each output is $Z[i]-Z[i+1]+0.5=-0.5$, giving three outputs. Stacking
all output coordinates gives $\operatorname{vec}(V)=\mathsf W\operatorname{vec}(Z)+\mathsf b$.
Writing $T_o=T-k+1$, matrix $\mathsf W$ has shape
$(C_{\rm out}T_o)\times(C_{\rm in}T)$, and $\mathsf b\in\R^{C_{\rm out}T_o}$
repeats channel biases. The sparse matrix's repeated entries encode weight
sharing. This is affine, not generally linear when bias is nonzero.
\paragraph{Why ECDR needs it.}
The affine representation allows interval sign rules and linear
relaxation methods to treat convolution rigorously. Local shared weights
also support exact shifted feature equality, but only at correctly aligned
positions whose receptive fields contain equal inputs.
\paragraph{Conditions and connections.}
Stride, dilation, grouping and padding change connectivity and output
length; they must be stated rather than inferred from the word CNN.
The historical specialized example uses valid stride~1/dilation~1 convolutions; the current reference has padding, BN and pooling.
The separate English CNN document gives an individual page to each
operator parameter. \hyperref[concept:B24]{B24} proves the corresponding feature identity.
\par\label{endconcept:B19}
\typeout{ECDR-CONCEPT-END B19 PAGE=\thepage}

\concept{B20}{ReLU creates affine pieces without preserving every margin}{\hyperref[concept:F16]{F16}--\hyperref[concept:F17]{F17}, \hyperref[concept:B19]{B19}}
\paragraph{Foundation and definition.}
The rectified linear unit is the pointwise scalar function
$\relu(v)=\max(0,v)$. Applied to tensor $V$, it gives $H[o,i]=\relu(V[o,i])$
without changing shape or introducing trainable parameters.
\paragraph{Symbols and meaning.}
$v\in\R$ is one preactivation, $V$ is a preactivation tensor, $H$ is
its activation tensor, and $o,i$ index a feature channel and position.
For $m\ge1$ neurons, let $w\in\R^m$ be their preactivation vector.
On a fixed activation pattern, define $D\in\R^{m\times m}$ diagonal
with entries~0 or~1. Then $\relu(w)=Dw$. Neurons identically zero may use a zero gate.
\paragraph{Derivation and example.}
If every preactivation sign is fixed on an input region, replacing each
ReLU by its fixed gate makes the full affine/ReLU network affine there:
composition of affine maps is affine. For a scalar affine preactivation
$v(x)=2x-1$, ReLU is~0 when $x\le1/2$ and $2x-1$ when $x\ge1/2$.
The pieces agree at the boundary but have different slopes.
For any scalar $a,b$, case analysis gives
$|\relu(a)-\relu(b)|\le|a-b|$; coordinatewise aggregation yields the
same nonexpansion for $L_1,L_2,L_\infty$ norms.
\paragraph{Why ECDR needs it.}
Stable activation patterns allow exact affine reasoning; unstable neurons
require sound bounds. A pattern observed only at the anchor is not
established over an entire candidate region. Network layers around ReLU
may amplify differences even though ReLU alone is nonexpansive.
\paragraph{Conditions and connections.}
Nonnegative activations do not imply increasing class margins: head
coefficients can be negative. \hyperref[concept:C04]{C04} propagates intervals; \hyperref[concept:C05]{C05} derives valid
affine bounds for a neuron whose sign can change.
\par\label{endconcept:B20}
\typeout{ECDR-CONCEPT-END B20 PAGE=\thepage}

\concept{B21}{Receptive-field span and jump on the input clock}{\hyperref[concept:B01]{B01}--\hyperref[concept:B02]{B02}, \hyperref[concept:B19]{B19}}
\paragraph{Foundation and definition.}
An architectural receptive field describes input positions that can affect
one feature. Its span is the distance between extreme potentially accessed
indices plus1. The jump is the input-grid displacement between consecutive
feature positions. This temporal extent is not a vector-space linear span.
\paragraph{Symbols and meaning.}
For local layer $\ell$, let $k_\ell\ge1$ be kernel taps, $d_\ell\ge1$
dilation and $a_\ell\ge1$ stride. Let $R_\ell$ be receptive-field span
in samples and $J_\ell$ input-grid jump in samples. With
$R_0=J_0=1$,
\[
 R_\ell=R_{\ell-1}+(k_\ell-1)d_\ell J_{\ell-1},\qquad
 J_\ell=a_\ell J_{\ell-1}.
\]
\paragraph{Derivation and example.}
The extreme input features in the next kernel are
$(k_\ell-1)d_\ell$ positions apart. Multiply by their raw-grid spacing
$J_{\ell-1}$ and add one preceding field's extent. Advancing the output
by one moves $a_\ell$ preceding positions, giving the jump recurrence.
With kernels~5,3,3 and all strides/dilations~1, spans are~5,7,9 and jumps
remain1. At 50 Hz, RF~9 has first-to-last interval8/50=0.16 s, while
sample-count duration is9/50=0.18 s.
\paragraph{Why ECDR needs it.}
Feature equality requires the whole relevant receptive field to lie in
the shared portion. Input overlap~50\% therefore does not mean feature
overlap~50\%. Padding and striding require extra boundary and phase reasoning.
\paragraph{Conditions and connections.}
With dilation, the span may contain holes; zero learned weights may reduce
actual dependence. The recurrence describes potential connectivity, not
empirical feature importance. \hyperref[concept:B24]{B24} specializes it to a contiguous valid CNN.
\par\label{endconcept:B21}
\typeout{ECDR-CONCEPT-END B21 PAGE=\thepage}

\concept{B22}{Global average pooling and the affine classifier head}{\hyperref[concept:F09]{F09}, \hyperref[concept:F16]{F16}, \hyperref[concept:B20]{B20}}
\paragraph{Foundation and definition.}
Global average pooling (GAP) reduces a final feature map along time while
retaining each channel. For normalized input $Z$, define
\[
 m(Z)=\frac1n\sum_{i=0}^{n-1}H_i(Z),\qquad
 f_\theta(Z)=A\,m(Z)+\beta.
\]
\paragraph{Symbols and meaning.}
$H_i(Z)\in\R^{d_h}$ is the vector of final features at position $i$,
$n\ge1$ is the number of final positions, $d_h$ is feature-channel count,
$m(Z)\in\R^{d_h}$ is the pooled vector, $A\in\R^{K\times d_h}$ is
the head matrix and $\beta\in\R^K$ is its bias. Raw logits are
$F(x)=f_\theta(\mathcal P(x))$. Rows $A_{j,:}$ correspond to classes.
\paragraph{Derivation and example.}
For feature positions $(2,1),(3,1),(4,1),(5,1)$,
$m=(14/4,4/4)=(3.5,1)$. Averaging is linear in the feature map, so
equal channel sums give equal pooled vectors. It discards the explicit
ordering of feature positions at the head. It does not make the full CNN
invariant to arbitrary permutations of raw samples: the local features
themselves can change under such permutations.
\paragraph{Why ECDR needs it.}
Feature sums allow shared contributions to cancel exactly when a window
shifts. An affine head then turns a pooled change into a direct pairwise
margin change. A position-weighted Flatten--MLP head does not generally
have the same cancellation.
\paragraph{Conditions and connections.}
Equal position counts and the stated pooling/head operators are essential.
A nonlinearity after pooling needs separate analysis. \hyperref[concept:B25]{B25} derives the GAP
update, and \hyperref[concept:B26]{B26} connects it to decision certification rather than just
activation reuse.
\par\label{endconcept:B22}
\typeout{ECDR-CONCEPT-END B22 PAGE=\thepage}

\concept{B23}{Parameter count and arithmetic work are different quantities}{\hyperref[concept:B19]{B19}, \hyperref[concept:B22]{B22}}
\paragraph{Foundation and definition.}
Parameters count stored learnable scalars. Arithmetic work counts uses of
those scalars while evaluating one input. For a bias-enabled, groups~1
convolution with input channels $C_{\rm in}$, output channels $C_{\rm out}$,
kernel length $k$ and output length $T_{\rm out}$,
\[
 P=C_{\rm out}(C_{\rm in}k+1),\qquad
 M=T_{\rm out}C_{\rm out}C_{\rm in}k.
\]
\paragraph{Symbols and meaning.}
$P$ is parameter count, $M$ is multiply--accumulate (MAC) count under one
product-and-accumulation convention. The added1 in $P$ is one bias per
output channel; it is not another kernel tap. Output length multiplies
work because weights are reused across positions, not because new weights
are created. A $d_h$-to-$K$ affine head has $K(d_h+1)$ parameters and
$Kd_h$ MAC under the same convention.
\paragraph{Derivation and example.}
For valid kernels~5,3,3, channels $6\to32\to64\to64$ and input length~128,
output lengths are124,122,120. The three parameter counts are~992,6208,
12352; a 12-class head adds~780, totaling~20332. MAC counts are~119040,
749568,1474560 and768, totaling~2343936. Reporting two FLOP/MAC gives
4687872 FLOP for these operators only.
\paragraph{Why ECDR needs it.}
MAC reduction suggests computational opportunity but is not a latency
theorem. ReLU, GAP, bias additions, normalization, memory traffic, dispatch
and certificate work still cost time. \hyperref[concept:C26]{C26}--\hyperref[concept:C27]{C27} account for that overhead.
\paragraph{Conditions and connections.}
Precision, kernel algorithms, buffer schedules and hardware affect actual
runtime and peak memory. Float32 weight payload~20332$\times$4=81328 bytes
is not checkpoint file size or process RAM. These are theoretical counts,
not measured quality or speed of a new HAPT model.
\par\label{endconcept:B23}
\typeout{ECDR-CONCEPT-END B23 PAGE=\thepage}

\concept{B24}{Exact aligned feature equality through a valid CNN}{\hyperref[concept:B10]{B10}, \hyperref[concept:B19]{B19}--\hyperref[concept:B21]{B21}}
\paragraph{Foundation and definition.}
Shared input samples induce shared features only under an operator contract.
Assume frozen time-homogeneous per-channel preprocessing, valid stride~1/dilation~1 convolutions,
position-independent pointwise activations, fixed weights and no pooling
or hidden state. Let the final receptive-field span be $R$ and final
length be $n=W-R+1$.
\paragraph{Symbols and meaning.}
$Z,Z'\in\R^{C\times W}$ are consecutive normalized windows shifted by
$s$ samples, with $1\le s\le n$;
$Z'[:,i]=Z[:,i+s]$ for $0\le i<W-s$.
$H_i(Z)\in\R^{d_h}$ is the final feature at position $i$.
The claim is
\[
 H_i(Z')=H_{i+s}(Z),\qquad 0\le i<n-s.
\]
\paragraph{Derivation.}
The new feature at $i$ reads input indices $i,\ldots,i+R-1$.
For $i\le n-s-1$, the last index is at most $W-s-1$, inside overlap.
Each sample there equals the old sample at offset $s$. Apply equal
kernels, biases and pointwise activations layer by layer to obtain equal
outputs. The proof establishes shifted equality, not equality at the same
tensor position.
\paragraph{Why ECDR needs it.}
For $W=128,R=9,s=64$, final length~120 leaves~56 shared positions.
The~64 new final features depend on~72 raw columns:~64 new samples plus
an~8-sample halo. This explains why forwarding only the~64 newly arrived
raw columns through the entire valid backbone misses needed features.
\paragraph{Conditions and connections.}
Padding, striding, per-window preprocessing or position-dependent operators
require a revised proof. Equality of features alone does not fix the
new pooled sum or class. \hyperref[concept:B25]{B25}--\hyperref[concept:B26]{B26} derive the remaining changes.
\par\label{endconcept:B24}
\typeout{ECDR-CONCEPT-END B24 PAGE=\thepage}

\concept{B25}{The GAP update is a cancellation identity}{\hyperref[concept:F09]{F09}, \hyperref[concept:B22]{B22}, \hyperref[concept:B24]{B24}}
\paragraph{Foundation and definition.}
When final feature maps share a shifted interior, their sums differ only
in dropped and added positions. Define
\[
 D_{\rm drop}=\sum_{i=0}^{s-1}H_i(Z),\qquad
 D_{\rm new}(U)=\sum_{i=n-s}^{n-1}H_i(Z'),
 \qquad Z'=\mathcal T_s(Z,U).
\]
\paragraph{Symbols and meaning.}
$Z,Z'\in\R^{C\times W}$ are old/new normalized windows; $U\in\R^{C\times s}$
is the new wave; $n$ is final feature length; $1\le s\le n$ is shift;
$H_i\in\R^{d_h}$ is a feature vector. Both $D_{\rm drop}$ and
$D_{\rm new}$ are $d_h$-dimensional sums. GAP is
$m(Z)=n^{-1}\sum_iH_i(Z)$.
\paragraph{Derivation and example.}
By \hyperref[concept:B24]{B24}, the shared new sum equals $\sum_{i=s}^{n-1}H_i(Z)$. Therefore
\[
 \sum_iH_i(Z')-\sum_iH_i(Z)=D_{\rm new}(U)-D_{\rm drop},
 \quad m(Z')-m(Z)=\frac{D_{\rm new}(U)-D_{\rm drop}}n.
\]
For old features $(0,1),(1,1),(2,1),(3,1)$ and new features
$(2,1),(3,1),(4,1),(5,1)$, take $n=4,s=2$.
The two sums are $(1,2)$ and $(9,2)$, giving pooled delta $(2,0)$.
\paragraph{Why ECDR needs it.}
Cancellation isolates the unknown future contribution. If that contribution
can be bounded over a wave region, the classifier may be certified before
computing exact features for an observed query. Computing its actual sum
instead gives an incremental evaluation, a different form of reuse.
\paragraph{Conditions and connections.}
This is exact real algebra under \hyperref[concept:B24]{B24}'s assumptions and equal feature
lengths. Floating-point sum updates may not match recomputation bitwise.
An identity provides no bound on all possible new waves by itself;
\hyperref[concept:B26]{B26} and \hyperref[concept:C19]{C19} provide the margin connection.
\par\label{endconcept:B25}
\typeout{ECDR-CONCEPT-END B25 PAGE=\thepage}

\concept{B26}{A pooled-feature change induces a pairwise-margin change}{\hyperref[concept:B16]{B16}, \hyperref[concept:B22]{B22}, \hyperref[concept:B25]{B25}}
\paragraph{Foundation and definition.}
An affine head turns a GAP change into an exact scalar margin change.
Define the normalized-network margin
\[
 \gamma_{cj}(Z)=f_{\theta,c}(Z)-f_{\theta,j}(Z),\qquad
 q_j=(A_{c,:}-A_{j,:})^\top.
\]
\paragraph{Symbols and meaning.}
$A\in\R^{K\times d_h}$ and $\beta\in\R^K$ are head weights and bias;
$c$ is target class, $j\ne c$ is a competitor, and $q_j\in\R^{d_h}$
is a column vector. $Z$ and $Z'$ are old/new normalized inputs;
$n$ is feature length. For raw query $x$, $g_{cj}(x)=\gamma_{cj}(\mathcal P(x))$.
The new wave is $U\in\R^{C\times s}$, where $C$ is input-channel count and
$1\le s\le n$ is the shift. Vectors $D_{\rm new}(U),D_{\rm drop}\in\R^{d_h}$
sum newly added and dropped feature positions, respectively (\hyperref[concept:B25]{B25}).
\paragraph{Derivation and example.}
Using the head formula and subtracting old from new cancels its bias:
\[
 \gamma_{cj}(Z')=\gamma_{cj}(Z)+
 \frac1nq_j^\top\big(D_{\rm new}(U)-D_{\rm drop}\big).
\]
The bias still contributes to the old margin; it must not be removed
from that value. In \hyperref[concept:B25]{B25}'s example, pooled delta is $(2,0)$.
For $q_2=(2,-1)^\top$ the margin increases~4; for
$q_3=(-1,1)^\top$ it decreases~2. The same feature change helps one
comparison and harms another.
\paragraph{Why ECDR needs it.}
This formula expresses exactly what needs a bound for class preservation.
Bounding every feature or logit independently can lose cancellation;
bounding the pairwise scalar objective can be more informative. The
whole-domain lower bound must remain valid for all admitted waves.
\paragraph{Conditions and connections.}
The identity needs the specified affine head and \hyperref[concept:B25]{B25}'s feature update.
It is neither an empirical stability estimate nor a certificate until
the unknown term is enclosed. \hyperref[concept:C19]{C19} chooses the correct feature-sum bounds
according to the signs of the coefficients in $q_j$.
\par\label{endconcept:B26}
\typeout{ECDR-CONCEPT-END B26 PAGE=\thepage}

\concept{B27}{A complete miniature CNN: forward arithmetic}{\hyperref[concept:B19]{B19}--\hyperref[concept:B22]{B22}, \hyperref[concept:B26]{B26}}
\paragraph{Foundation and definition.}
A small fully specified network makes each stage independently checkable.
Take identity preprocessing, one channel, window length~6, shift~2, a
valid length~3 convolution, two filters, ReLU, GAP and a three-class head.
\paragraph{Symbols and meaning.}
The old input is $X=(1,2,3,4,5,6)$, the appended wave is $U=(7,8)$,
and $X'=(3,4,5,6,7,8)$. Filter weights and biases are
\[
 w_1=(1,-1,1),\ b_1=-2;\qquad w_2=(-1,0,1),\ b_2=-1.
\]
The head matrix $A\in\R^{3\times2}$ and bias $\beta\in\R^3$ are
\[
 A=\begin{pmatrix}1&-1\\-1&0\\2&-2\end{pmatrix},\qquad
 \beta=(4,1,0)^\top.
\]
\paragraph{Derivation.}
Filter~1 on the first old patch gives $1-2+3-2=0$; filter~2 gives
$-1+3-1=1$. The position-wise feature vectors after ReLU are
\[
 H(X)=((0,1),(1,1),(2,1),(3,1)),\quad
 H(X')=((2,1),(3,1),(4,1),(5,1)).
\]
Thus $m(X)=(3/2,1)$, $m(X')=(7/2,1)$ and the raw/normalized logits,
identical here because preprocessing is identity, are
\[
 F(X)=(9/2,-1/2,1),\qquad F(X')=(13/2,-5/2,5).
\]
Both choose class~1; its margins change from $(5,7/2)$ to $(9,3/2)$.
\paragraph{Why ECDR needs it.}
The example verifies score ordering, feature alignment, GAP cancellation
and opposite margin responses in a model whose every operation is visible.
It illustrates a mechanism rather than demonstrating HAPT accuracy.
\paragraph{Conditions and connections.}
All values here are exact rational arithmetic. No training or data fit is
claimed. One observed wave is not a region certificate; \hyperref[concept:B28]{B28} varies it over
a continuous box and proves the necessary signs.
\par\label{endconcept:B27}
\typeout{ECDR-CONCEPT-END B27 PAGE=\thepage}

\concept{B28}{Certifying a continuous wave box in the miniature CNN}{\hyperref[concept:B27]{B27}, \hyperref[concept:C03]{C03}, \hyperref[concept:C18]{C18}--\hyperref[concept:C19]{C19}}
\paragraph{Foundation and definition.}
Replace \hyperref[concept:B27]{B27}'s single wave by $U=(u,v)$ in the continuous region
$\mathcal U=[6.5,7.5]\times[7.5,8.5]$. The old anchor and network remain
fixed. Every new input is $X'=(3,4,5,6,u,v)$.
\paragraph{Symbols and meaning.}
$u,v\in\R$ are the two new samples; $\mathcal U$ is a set, not an
assumed probability distribution. There are $n=4$ final positions;
$D_{\rm new}\in\R^2$ sums the two new feature vectors.
The two class~1 margins are $g_{12}$ and $g_{13}$, with competitors~2 and~3.
\paragraph{Derivation.}
The new tail preactivations are
\[
 V_2=(u-3,u-6),\qquad V_3=(4-u+v,v-7).
\]
Their coordinate lower bounds are $(3.5,0.5)$ and $(4,0.5)$, all positive.
Thus post-ReLU features satisfy $H_2=V_2,H_3=V_3$ throughout the box,
not just at its center.
Summing symbolically gives
\[
 D_1=1+v,\quad D_2=u+v-13,\quad
 g_{12}=5+\frac{15+v-u}{4},\quad g_{13}=\frac{u-1}{4}.
\]
$D_1,D_2$ are feature-sum coordinates. Optimizing these affine margins
over the stated intervals yields $\min g_{12}=35/4$ and
$\min g_{13}=11/8$, both positive. Every wave in the box preserves class~1.
\paragraph{Why ECDR needs it.}
This is an actual whole-domain argument without assuming independent or
uniform samples. Keeping shared variables in the objective retains
dependencies. Bounding tail positions independently gives the looser but
still sound lower margins~8 and~7/8.
\paragraph{Conditions and connections.}
An affine objective's corner check can verify arithmetic only after stable
ReLU branches and affine dependence are established. The proof is tied to
the anchor's overlap; \hyperref[concept:C24]{C24} shows why repeating the wave box at a new anchor
without new evidence can fail.
\par\label{endconcept:B28}
\typeout{ECDR-CONCEPT-END B28 PAGE=\thepage}

\concept{B29}{A Cartesian product region does not assert independence}{\hyperref[concept:F24]{F24}, \hyperref[concept:F30]{F30}, \hyperref[concept:B12]{B12}}
\paragraph{Foundation and definition.}
Geometry specifies permitted values; probability specifies how values are
distributed. A finite-horizon wave region may be the Cartesian product
\[
 \mathcal V=\mathcal U_1\times\cdots\times\mathcal U_Q.
\]
\paragraph{Symbols and meaning.}
$Q\ge1$ is the number of future steps, $\mathcal U_q\subseteq\R^{Cs}$
is the admissible wave set at step $q$, $C$ is channels and $s$ is
samples per wave. A member of $\mathcal V$ is a sequence
$(U_1,\ldots,U_Q)$ with $U_q\in\mathcal U_q$. No probability measure
is specified by this set notation.
\paragraph{Derivation and example.}
Let $Q=2$, $\mathcal U_1=\mathcal U_2=[-1,1]$, and random variables
satisfy $U_2=U_1$ with $U_1$ uniform on $[-1,1]$. Their pair always
belongs to the product square but lies only on its diagonal. The variables
are dependent, while a proof that holds throughout the square also holds
on that diagonal. Independence would be a factorization statement about
joint probabilities, not a consequence of the square's shape.
\paragraph{Why ECDR needs it.}
Finite-horizon verification can use a product region without assuming that
sensor waves are statistically independent. A dependent reachable set may
be smaller and easier to certify, but its constraints need a justified
membership or transition contract. A geometry that admits impossible
sequences can still be sound as an overapproximation, though less useful.
\paragraph{Conditions and connections.}
A product domain by itself supplies no coverage probability and no
guarantee that future waves remain inside it. Conversely, observed temporal
correlation is not a deterministic constraint on every future wave.
\hyperref[concept:C25]{C25} combines transition geometry with induction; \hyperref[concept:B12]{B12} handles joint events.
\par\label{endconcept:B29}
\typeout{ECDR-CONCEPT-END B29 PAGE=\thepage}

\concept{B30}{Classification metrics describe a specified labeled population}{\hyperref[concept:F09]{F09}, \hyperref[concept:F31]{F31}--\hyperref[concept:F32]{F32}, \hyperref[concept:B06]{B06}, \hyperref[concept:B15]{B15}}
\paragraph{Foundation and definition.}
Prediction quality against true labels differs from reuse agreement with
a reference. For a confusion matrix with true rows and predicted columns,
consider
\[
 M=\begin{pmatrix}8&2\\1&9\end{pmatrix}.
\]
\paragraph{Symbols and meaning.}
$M_{ij}$ counts examples truly in class $i$ and predicted as class $j$;
$N=\sum_{i,j}M_{ij}=20$ is population size. For a class, TP means true
positives, FP false positives and FN false negatives. Precision is
$\mathrm{TP}/(\mathrm{TP}+\mathrm{FP})$, recall is
$\mathrm{TP}/(\mathrm{TP}+\mathrm{FN})$, and
\[
 \mathrm{F1}=\frac{2\mathrm{TP}}{2\mathrm{TP}+\mathrm{FP}+\mathrm{FN}},
 \qquad \mathrm{MacroF1}=K^{-1}\sum_{j=1}^{K}\mathrm{F1}_j.
\]
$K$ is the number of classes included in that average, not an adjustable
denominator chosen after seeing results.
\paragraph{Derivation and example.}
Accuracy is $(8+9)/20=0.85$. Class1 has TP~8, FP~1, FN~2, so F1 is~16/19.
Class2 has F1 $18/21=6/7$. Macro-F1 is
$(16/19+6/7)/2=113/133\approx0.849624$, not exactly accuracy.
Zero denominators need an explicit convention. Across seeds, mean and
standard deviation are descriptive summaries; standard deviation alone
is not a confidence interval or an independent test split.
\paragraph{Why ECDR needs it.}
A reference may misclassify a true activity while a reuse mechanism
faithfully returns the same mistaken class. Accuracy, reference agreement,
certified coverage and net runtime benefit therefore answer different
questions. Filtering to pure K6 changes the labeled population relative
to full K12 or the all-window stream.
\paragraph{Conditions and connections.}
Declare labels, class schema, policy, split, support and aggregation before
interpreting a metric. Temporal dependence affects uncertainty estimates
even when the arithmetic is correct. \hyperref[concept:C28]{C28} gives separate denominators for
certification and reuse.
\par\label{endconcept:B30}
\typeout{ECDR-CONCEPT-END B30 PAGE=\thepage}

\concept{C01}{Lipschitz continuity and compositional network bounds}{Norms, induced operator norms, and function composition; B-series classifier notation.}
\paragraph{Definition and foundation.}
A Lipschitz constant limits how much a function can amplify a change in its input. For a map $f:\R^d\to\R^m$ on a set $S\subseteq\R^d$, a valid constant $L\ge0$ satisfies
\[
 \norm{f(x)-f(y)}_q\le L\norm{x-y}_p\qquad\text{for every }x,y\in S.
\]
The quantifier is essential: a ratio observed on finitely many pairs is evidence about those pairs, rather than an upper bound on every possible pair. Lipschitz continuity implies continuity, but ordinary continuity need not supply a usable linear rate.
\paragraph{Symbols and dimensions.}
Here $d,m$ are positive integer dimensions; $x,y$ are input columns; $p,q$ identify norms; $L$ has output-unit/input-unit scale. Normalized sensor coordinates and logits are usually dimensionless. A bound must still name the normalization because changing coordinates changes its numerical value.
\paragraph{Derivation and example.}
For $f(x)=Ax+b$ with $A\in\R^{m\times d}$, subtraction cancels the bias $b\in\R^m$. The induced norm gives $\norm{A(x-y)}_q\le\norm{A}_{p\to q}\norm{x-y}_p$. If $f=f_r\circ\cdots\circ f_1$, integer $r\ge1$ counts layers, and layer $i$ has compatible constant $L_i$, repeated substitution yields
\[
 \operatorname{Lip}(f)\le\prod_{i=1}^{r}L_i.
\]
ReLU is 1-Lipschitz for the usual $p$-norms because its coordinatewise differences do not exceed the original differences. For $A=\left(\begin{smallmatrix}1&-2\\0&3\end{smallmatrix}\right)$, the $\infty$-norm constant is the largest absolute row sum, namely 3.
\paragraph{Why ECDR uses this concept.}
A valid rate converts an input-distance budget into a worst-case logit or margin change. It therefore connects data geometry to decision invariance.
\paragraph{Conditions and connections.}
The layer bounds must hold on every intermediate set reached from $S$. Their product may be loose because different layers' worst directions need not align. A gradient at one point is insufficient. \hyperref[concept:C02]{C02} combines a valid rate with the anchor margin.
\par\label{endconcept:C01}
\typeout{ECDR-CONCEPT-END C01 PAGE=\thepage}

\concept{C02}{A margin certificate and its admissible radius}{\hyperref[concept:C01]{C01}; pairwise margin, unique winner, and norm balls.}
\paragraph{Definition and foundation.}
A distance certificate combines a positive margin at one evaluated input with a bound on how quickly that margin can fall. Let $F:\R^d\to\R^K$ be the fixed logit map, $x_r\in\R^d$ an evaluated anchor, and $c\in\{1,\ldots,K\}$ its target class. For each competitor $j\ne c$, define $g_{cj}(x)=F_c(x)-F_j(x)$.
\paragraph{Symbols and dimensions.}
The integer $K\ge2$ counts classes; $r$ labels the anchor, not a radius. Write $\gamma_j=g_{cj}(x_r)>0$ and assume a valid constant $L_j\ge0$ for this scalar margin throughout the considered domain. A radius $\epsilon\ge0$ uses the same norm units as $x$. The index $j$ ranges over all $K-1$ competitors.
\paragraph{Derivation and example.}
From $\abs{g_{cj}(x)-\gamma_j}\le L_j\norm{x-x_r}_p$, the lower half of the absolute-value inequality gives
\[
 g_{cj}(x)\ge\gamma_j-L_j\norm{x-x_r}_p
 \ge\gamma_j-L_j\epsilon.
\]
Consequently, a closed ball of radius $\epsilon$ is certified by the strict gate $\gamma_j-L_j\epsilon>0$ for every $j\ne c$. With margins $(0.8,0.3)$ and constants $(2,1)$, radius 0.2 gives bounds $(0.4,0.1)$. Radius 0.3 gives a zero second bound and does not pass this strict gate.
\paragraph{Why ECDR uses this concept.}
Once these bounds are valid, runtime can test membership in the certified ball and return $c$. Confidence or small distance alone supplies neither half of this argument.
\paragraph{Conditions and connections.}
For $L_j>0$, use $\epsilon<\gamma_j/L_j$; a competitor with $L_j=0$ and $\gamma_j>0$ imposes no further radius restriction within the domain where its bound is valid. These are sufficient conditions. A failed gate does not prove a changed class. A smaller certified ball may exist, and a non-ball domain can succeed where this isotropic construction fails. Numerical executable guarantees additionally require \hyperref[concept:C20]{C20}--\hyperref[concept:C21]{C21}.
\par\label{endconcept:C02}
\typeout{ECDR-CONCEPT-END C02 PAGE=\thepage}

\concept{C03}{Interval bounds for an affine operation}{Order on real numbers, coordinate intervals, and matrix multiplication.}
\paragraph{Definition and foundation.}
An interval enclosure records a lower and upper bound for every possible coordinate. Suppose an unknown vector $v\in\R^n$ satisfies $l\le v\le u$ coordinatewise, and an affine operation computes $w=Av+b\in\R^m$. The problem is to enclose every resulting $w$, rather than to estimate its average.
\paragraph{Symbols and dimensions.}
The columns $l,u\in\R^n$ obey $l_i\le u_i$. The matrix $A\in\R^{m\times n}$ and bias $b\in\R^m$ are fixed. Define $A^+_{ki}=\max(A_{ki},0)$ and $A^-_{ki}=\min(A_{ki},0)$. Thus $A^-$ contains nonpositive entries; it is not the absolute magnitude of negative entries. All coordinates use the input or activation units of their layer.
\paragraph{Derivation and example.}
For a positive coefficient, the smallest product uses the smaller input. Multiplication by a negative coefficient reverses order, so the smallest product uses the larger input. Adding these scalar inequalities row by row yields
\[
 \underline w=A^+l+A^-u+b,\qquad
 \overline w=A^+u+A^-l+b.
\]
For $v_1\in[1,2]$, $v_2\in[-1,3]$, and $w=2v_1-3v_2+1$, the resulting interval is $[-6,8]$. The lower endpoint is $2(1)-3(3)+1=-6$. Using both lower inputs would produce 6, which fails to enclose valid outputs.
\paragraph{Why ECDR uses this concept.}
Convolutions, dense heads, fixed affine normalization, and channel sums all admit this rule. Shared convolution weights do not require an explicitly materialized dense matrix.
\paragraph{Conditions and connections.}
The rule is exact for each output coordinate over the full input box, but it can lose relationships between outputs. If the true input set occupies only part of the box, the result remains sound and may be loose. \hyperref[concept:C04]{C04} propagates the enclosure through a network; \hyperref[concept:C19]{C19} applies the same sign logic to tail sums.
\par\label{endconcept:C03}
\typeout{ECDR-CONCEPT-END C03 PAGE=\thepage}

\concept{C04}{Interval propagation and inductive soundness}{\hyperref[concept:F35]{F35}, \hyperref[concept:C03]{C03}; monotone functions and induction.}
\paragraph{Definition and foundation.}
Interval bound propagation, or IBP, transports an input enclosure through each operation of a fixed computation graph. Its central invariant is set inclusion: every true intermediate activation lies inside the interval assigned to that activation. Soundness means preserving this invariant, rather than obtaining narrow intervals.
\paragraph{Symbols and dimensions.}
Let $S\subseteq\R^{d_0}$ be the input domain, $f_i:\R^{d_{i-1}}\to\R^{d_i}$ the $i$th layer, and $v_i=f_i(v_{i-1})$ its true activation, with $v_0=x\in S$. Bounds $l_i,u_i\in\R^{d_i}$ are coordinatewise enclosures. The layer index is $i=1,\ldots,r$.
\paragraph{Derivation and example.}
Begin with $S\subseteq[l_0,u_0]$. For an affine layer use \hyperref[concept:C03]{C03}. Because $\relu(a)=\max(0,a)$ is nondecreasing,
\[
 l_i=\relu(l_{i-1}),\qquad u_i=\relu(u_{i-1})
\]
for a ReLU layer. Thus $[-6,8]$ becomes $[0,8]$. A maximum over coordinates is enclosed by the maximum of their lowers and the maximum of their uppers. Assume the invariant at layer $i-1$; a sound enclosure rule establishes it at layer $i$. Induction proves the final enclosure for every $x\in S$.
\paragraph{Why ECDR uses this concept.}
If final logit intervals satisfy $l_{r,c}>u_{r,j}$ for every competitor $j\ne c$, class $c$ wins throughout $S$. This is one route to a reusable certificate.
\paragraph{Conditions and connections.}
Every operation must have a correct enclosure rule, including preprocessing. Training-time batch normalization is not automatically a fixed pointwise affine map. Independent intervals can lose dependencies: for $F_c(t)=t+1$ and $F_j(t)=t$ on $[-1,1]$, the true margin is 1, while separate logit intervals yield a lower margin of $-1$. \hyperref[concept:C06]{C06} explains a tighter relational alternative. Floating-point soundness requires more than this real-arithmetic induction.
\par\label{endconcept:C04}
\typeout{ECDR-CONCEPT-END C04 PAGE=\thepage}

\concept{C05}{Affine relaxations of an unstable ReLU}{\hyperref[concept:C03]{C03}--\hyperref[concept:C04]{C04}; affine lines, convexity, and inequality direction.}
\paragraph{Definition and foundation.}
A relaxation replaces an operation with inequalities that contain its entire graph on a specified interval. A ReLU is unstable when its input interval crosses zero: some admissible inputs activate it and others do not. Its exact behavior is then not one affine map over the entire interval.
\paragraph{Symbols and dimensions.}
Let $v\in[l,u]\subset\R$, with $l<0<u$, and let $r=\relu(v)$. The endpoints $l,u$ have activation units. Define the dimensionless slope $a_U=u/(u-l)$ and intercept $b_U=-ul/(u-l)$, with activation units. A lower slope $\alpha$ is any real number in $[0,1]$.
\paragraph{Derivation and example.}
The line through the graph endpoints $(l,0)$ and $(u,u)$ has slope $(u-0)/(u-l)$, giving
\[
 \alpha v\le\relu(v)\le a_Uv+b_U
 =\frac{u}{u-l}(v-l).
\]
For $v\le0$, the lower expression is nonpositive; for $v\ge0$, it is at most $v$. For the upper inequality, when $v\le0$ the chord is nonnegative. When $v\ge0$, subtracting $v$ from the chord gives $l(v-u)/(u-l)\ge0$. Both inequalities therefore hold on the entire interval.
For $[l,u]=[-1,2]$, the upper line is $2(v+1)/3$, including an upper value $2/3$ at $v=0$, where ReLU is exactly zero.
\paragraph{Why ECDR uses this concept.}
These lines let a verifier carry affine information through nonlinear CNN layers and retain dependencies that independent intervals discard.
\paragraph{Conditions and connections.}
If $l\ge0$, use the exact identity; if $u\le0$, use exact zero. The interval itself must be sound. A relaxation valid only on $[l,u]$ cannot be reused on a wider interval without justification. \hyperref[concept:C06]{C06} chooses the correct line according to the sign of the downstream coefficient.
\par\label{endconcept:C05}
\typeout{ECDR-CONCEPT-END C05 PAGE=\thepage}

\concept{C06}{Backward linear relaxation and sign-aware substitution}{\hyperref[concept:C05]{C05}; linear objectives, lower bounds, and cancellation.}
\paragraph{Definition and foundation.}
Linear relaxation based perturbation analysis replaces nonlinear terms while preserving a bound on the particular output objective of interest. To lower-bound a scalar term $a r$, use a lower bound on $r$ when $a\ge0$, and an upper bound when $a<0$. A negative multiplier reverses inequality direction.
\paragraph{Symbols and dimensions.}
Here $r$ is a scalar activation and $a$ its downstream coefficient. For the example let $x\in[-1,1]$, $r_1=\relu(x)$, $r_2=\relu(x+1)$, and margin $g(x)=3/2+2r_1-r_2$. All example quantities are dimensionless. The lower ReLU slope is $\alpha\in[0,1]$.
\paragraph{Derivation and example.}
Because $x+1\in[0,2]$, $r_2=x+1$ exactly. Because the coefficient of $r_1$ is positive, substitute $r_1\ge\alpha x$:
\[
 g(x)\ge\tfrac32+2\alpha x-(x+1)
 =\tfrac12+(2\alpha-1)x.
\]
Minimizing the affine expression on $[-1,1]$ gives $1/2-\abs{2\alpha-1}$. Choosing $\alpha=1/2$ certifies the exact minimum $1/2$, since the true expression is $g(x)=1/2+\abs{x}$. Independent intervals $r_1\in[0,1]$ and $r_2\in[0,2]$ instead give $-1/2$.
\paragraph{Why ECDR uses this concept.}
Backward substitution can produce $g_{cj}(x)\ge a_x^\top x+\beta_x$, keeping shared input dependencies until the final domain optimization. Here $a_x\in\R^d$ and $\beta_x\in\R$ define that affine lower expression.
\paragraph{Conditions and connections.}
Every substitution must be valid over the reachable activation range. The example explains a CROWN/LiRPA principle, not every algorithmic variant. Relaxation quality and computation cost remain separate questions. A sign error creates an invalid certificate; a loose but sound relaxation merely fails to certify. \hyperref[concept:C07]{C07} converts the final affine expression into a constant domain bound. Optimizing the relaxation slope can improve that constant without changing the domain. It changes the proof approximation, rather than changing the underlying CNN or adding information to the observed input.
\par\label{endconcept:C06}
\typeout{ECDR-CONCEPT-END C06 PAGE=\thepage}

\concept{C07}{Support functions and minimization of affine bounds}{Linear functionals, supremum and infimum; \hyperref[concept:C06]{C06}.}
\paragraph{Definition and foundation.}
For a nonempty set $D\subseteq\R^d$, its support function in direction $v\in\R^d$ is $\sigma_D(v)=\sup_{\delta\in D}v^\top\delta$. It measures the farthest extent of the set along a linear direction. If $D$ is unbounded, this value can be $+\infty$; nonempty compact domains avoid that difficulty and attain their extrema.
\paragraph{Symbols and dimensions.}
Write the input domain as $S=c_0+D$, where $c_0\in\R^d$ is its center and $\delta\in D$ is a displacement. Suppose the scalar margin obeys $g(x)\ge a^\top x+\beta$ throughout $S$, with $a\in\R^d$ and scalar $\beta$. The subscript 0 distinguishes the center from a class label. Products $a^\top x$ and $\beta$ use margin units.
\paragraph{Derivation and example.}
Substitute $x=c_0+\delta$ and use $\inf z=-\sup(-z)$:
\[
 \inf_{x\in S}(a^\top x+\beta)
 =a^\top c_0+\beta-\sigma_D(-a).
\]
For the asymmetric interval $D=[0,2]$ and $a=-1$, the correct value $\sigma_D(-a)=\sigma_D(1)=2$ gives minimum $-2$ when $c_0=\beta=0$. Incorrectly using $\sigma_D(a)=\sigma_D(-1)=0$ would claim minimum zero and fail at displacement 2. Thus a missing minus sign can invalidate the bound, rather than merely reduce its precision.
\paragraph{Why ECDR uses this concept.}
Network relaxation and domain geometry perform different jobs. The former supplies an affine lower expression; the support function evaluates its worst case over the chosen input domain.
\paragraph{Conditions and connections.}
Only a centrally symmetric displacement set, $D=-D$, permits replacing $\sigma_D(-a)$ with $\sigma_D(a)$. The lower expression must hold for every domain point. \hyperref[concept:C08]{C08}--\hyperref[concept:C10]{C10} and \hyperref[concept:C13]{C13} derive supports for balls, boxes, generators, and orthogonal residuals. The construction asks a directional worst-case question about a set; it neither averages data points nor assigns probabilities to them.
\par\label{endconcept:C07}
\typeout{ECDR-CONCEPT-END C07 PAGE=\thepage}

\concept{C08}{Norm-ball support and duality through Hölder's inequality}{\hyperref[concept:C07]{C07}; $p$-norms, dual norms, and Hölder's inequality.}
\paragraph{Definition and foundation.}
The dual norm of a vector $a\in\R^d$ is the largest linear response $a^\top v$ over the unit ball of the chosen input norm. For the standard $p$-norm, $1\le p\le\infty$, the dual is the $p^*$-norm, where $1/p+1/p^*=1$, interpreting the endpoints as $1^*=\infty$ and $\infty^*=1$.
\paragraph{Symbols and dimensions.}
Let $D=\{\delta\in\R^d:\norm{\delta}_p\le\epsilon\}$ with radius $\epsilon\ge0$, and let $S=c_0+D$. An affine margin lower bound has coefficient $a\in\R^d$, center $c_0\in\R^d$, and intercept $\beta\in\R$. The radius has input units; $\epsilon\norm{a}_{p^*}$ has margin units.
\paragraph{Derivation and example.}
Hölder gives $a^\top\delta\le\norm{a}_{p^*}\norm{\delta}_p\le\epsilon\norm{a}_{p^*}$. In finite dimension, a maximizing direction on the closed unit ball attains the dual norm. Hence this upper bound is the exact support:
\[
 \inf_{x\in S}(a^\top x+\beta)
 =a^\top c_0+\beta-\epsilon\norm{a}_{p^*}.
\]
For $a=(3,-4)^\top$, $c_0=0$, $\beta=2$, and $\epsilon=0.1$, the minima are $1.5$ for a Euclidean ball, $1.3$ for an $\infty$-ball, and $1.6$ for a 1-ball. Their penalties are respectively $0.1(5)$, $0.1(7)$, and $0.1(4)$.
\paragraph{Why ECDR uses this concept.}
It turns an affine verifier output into a scalar lower margin using a closed-form geometric calculation.
\paragraph{Conditions and connections.}
The three examples use different sets: an equal numeric radius does not make coverage or volume comparable. This exact support does not make the preceding neural relaxation exact. It evaluates that relaxation exactly. \hyperref[concept:C09]{C09} replaces a uniform radius with coordinate-specific budgets. A Euclidean ball pays according to the coefficient vector's Euclidean length; a coordinate box can choose adverse signs separately. This difference explains the distinct penalties without appealing to a probabilistic interpretation of either set.
\par\label{endconcept:C08}
\typeout{ECDR-CONCEPT-END C08 PAGE=\thepage}

\concept{C09}{Diagonal boxes: exact support and coordinate-access limits}{\hyperref[concept:C07]{C07}; Cartesian products and adversarial indistinguishability.}
\paragraph{Definition and foundation.}
A diagonal box is $S=\{c_0+\delta:\abs{\delta_i}\le b_i,\ i=1,\ldots,d\}$, with center $c_0\in\R^d$ and budgets $b_i\ge0$. Coordinate constraints are independent geometrically; this says nothing about statistical independence of the observed sensor channels.
\paragraph{Symbols and dimensions.}
The affine lower expression is $a^\top x+\beta$, where $a\in\R^d$, $\beta\in\R$, and each $b_i$ has the units of coordinate $x_i$. The integer $d\ge1$ is the number of coordinates tested. Define a coordinate-access model in which reading one previously unknown query coordinate has unit cost.
\paragraph{Derivation and example.}
Each term $a_i\delta_i$ is minimized independently at $-b_i\abs{a_i}$. Summing yields
\[
 \inf_{x\in S}(a^\top x+\beta)
 =a^\top c_0+\beta-\sum_{i=1}^d b_i\abs{a_i}.
\]
For $a=(2,-1)^\top$, $b=(0.1,0.3)^\top$, $c_0=0$, and $\beta=1$, the lower value is $1-0.2-0.3=0.5$.
An exact membership algorithm has worst-case $\Omega(d)$ coordinate accesses for a finite nonempty box with every width positive. Run it on an interior point. If a coordinate remains unread, change only that coordinate beyond its finite bound. The algorithm observes the same values but must now return a different answer, a contradiction. A complete scan achieves $O(d)$, so the bound is $\Theta(d)$.
\paragraph{Why ECDR uses this concept.}
Both certification and membership have transparent formulas, making boxes a useful reference geometry.
\paragraph{Conditions and connections.}
The lower bound assumes exact membership for arbitrary inputs, no prior query summary, and all nondegenerate finite coordinates. Early rejection can improve average cost. Parallel hardware changes elapsed time, and an always-reject conservative checker avoids the exactness requirement while providing no reuse. Big-O alone establishes no speedup. The access argument concerns the information needed to distinguish inputs, not the number of machine instructions generated by one particular loop. It therefore survives routine code refactoring under the same input model.
\par\label{endconcept:C09}
\typeout{ECDR-CONCEPT-END C09 PAGE=\thepage}

\concept{C10}{Generator domains and additive residual support}{\hyperref[concept:C07]{C07}; matrix images, Minkowski sums, and coordinate boxes.}
\paragraph{Definition and foundation.}
A generator domain represents coordinated movement through a smaller latent vector. Let $S=\{c_0+Bu+v:\norm{u}_\infty\le1,\ \abs{v_i}\le\eta_i\}$. Its first component follows columns of $B$; its residual permits coordinatewise departures. These are set constraints, not a probability model.
\paragraph{Symbols and dimensions.}
The center is $c_0\in\R^d$, generator matrix $B\in\R^{d\times q}$, latent column $u\in\R^q$, residual $v\in\R^d$, and residual budgets $\eta_i\ge0$. The dimension $q\ge1$ counts generators. For affine lower bound $a^\top x+\beta$, $a\in\R^d$ and $\beta\in\R$. Columns of $B$ have input units when $u$ is dimensionless; they need not be orthogonal.
\paragraph{Derivation and example.}
The two constraint sets form a Cartesian product, so their infima separate. Since $a^\top Bu=(B^\top a)^\top u$,
\[
 \inf_{x\in S}(a^\top x+\beta)
 =a^\top c_0+\beta-\norm{B^\top a}_1
   -\sum_{i=1}^d\eta_i\abs{a_i}.
\]
For margin $1+x_1-x_2$, choose $c_0=0$, $B=(0.8,0.8)^\top$, and $\eta_1=\eta_2=0.1$. The generator contribution vanishes because $(1,-1)B=0$; the residual penalty is 0.2, leaving bound 0.8. The enclosing box $\abs{x_i}\le0.9$ instead gives $-0.8$.
\paragraph{Why ECDR uses this concept.}
A correlated temporal displacement can be broad in directions harmless to the decision margin, while keeping a small support in harmful directions.
\paragraph{Conditions and connections.}
Low dimension alone gives no guarantee: margin $1+x_1+x_2$ changes along the same generator. Coordinate residuals here are not orthogonal Euclidean residuals. \hyperref[concept:C11]{C11}--\hyperref[concept:C14]{C14} derive that distinct geometry and explain how normalization changes it. The support is especially cheap once $B^\top a$ is computed, but checking a query's generator representation can be substantially harder. Proof cost and membership cost must therefore be evaluated separately, even for this simple closed-form lower bound.
\par\label{endconcept:C10}
\typeout{ECDR-CONCEPT-END C10 PAGE=\thepage}

\concept{C11}{Linear span, orthogonal projection, and residual decomposition}{Vector spaces, inner products, and linear independence; \hyperref[concept:C10]{C10}.}
\paragraph{Definition and foundation.}
The linear span of columns $v_1,\ldots,v_q\in\R^d$ is the set of every real linear combination of those columns. It is a subspace because it contains zero and is closed under addition and scalar multiplication. A bounded generator domain occupies only a constrained portion of a span; ``span'' itself does not impose coefficient budgets.
\paragraph{Symbols and dimensions.}
Let $V=[v_1\ \cdots\ v_q]\in\R^{d\times q}$ have orthonormal columns, so $V^\top V=I_q$, with $1\le q\le d$. The symbols $I_q$ and $I_d$ denote identity matrices of their indicated sizes. For displacement $w\in\R^d$, define score $s=V^\top w\in\R^q$, projector $P=VV^\top\in\R^{d\times d}$, and complementary projector $P_\perp=I_d-P$.
\paragraph{Derivation and example.}
The identity $P^2=V(V^\top V)V^\top=P$ proves idempotence; $P^\top=P$ proves symmetry. Also $V^\top P_\perp=0$. Therefore
\[
 w=Vs+P_\perp w,\qquad
 \norm{w}_2^2=\norm{s}_2^2+\norm{P_\perp w}_2^2.
\]
The second equation follows because the two components are orthogonal. With $V=(1,1)^\top/\sqrt2$ and $w=(0.2,0)^\top$, the score is $0.1\sqrt2$, the projection is $(0.1,0.1)^\top$, and the residual is $(0.1,-0.1)^\top$.
\paragraph{Why ECDR uses this concept.}
A query can be tested against separate budgets along retained temporal directions and outside their span. This gives a precise meaning to ``structured motion plus residual uncertainty.''
\paragraph{Conditions and connections.}
The formula $VV^\top$ is an orthogonal projector only under the stated orthonormality condition. A general full-column-rank matrix $B$ requires $B(B^\top B)^{-1}B^\top$. \hyperref[concept:C12]{C12} obtains directions from training data, while \hyperref[concept:C14]{C14} explains why rescaling them can invalidate the original projector. The residual is the component not represented by the chosen span; it is not automatically measurement noise, annotation error, or a random variable independent of the scores.
\par\label{endconcept:C11}
\typeout{ECDR-CONCEPT-END C11 PAGE=\thepage}

\concept{C12}{PCA, explained variance, and what a fitted subspace does not prove}{\hyperref[concept:F32]{F32}, \hyperref[concept:F37]{F37}--\hyperref[concept:F39]{F39}, \hyperref[concept:C11]{C11}; training/test separation.}
\paragraph{Definition and foundation.}
Principal component analysis, or PCA, selects orthogonal directions maximizing variation in a finite collection of centered observations. It describes the data used to fit it. A high explained-variance fraction is neither a universal support constraint nor a statement that every future displacement lies close to the selected span.
\paragraph{Symbols and dimensions.}
Let $\Delta^{(i)}\in\R^d$, $i=1,\ldots,N$, be $N\ge2$ training displacements. Their mean is $\bar\Delta=N^{-1}\sum_i\Delta^{(i)}$. The sample covariance is
\[
 C=\frac1{N-1}\sum_{i=1}^{N}(\Delta^{(i)}-\bar\Delta)
 (\Delta^{(i)}-\bar\Delta)^\top\in\R^{d\times d}.
\]
Its eigenvalues satisfy $\lambda_1\ge\cdots\ge\lambda_d\ge0$. The retained matrix $V\in\R^{d\times q}$ contains $q$ orthonormal eigenvectors, with $1\le q\le d$.
\paragraph{Derivation and example.}
For a unit direction $v$, its sample projected variance is $v^\top Cv$. Maximizing this expression selects a leading eigenvector; successive orthogonal directions select subsequent eigenvectors. If $\sum_i\lambda_i>0$, the retained variance fraction is $\sum_{i=1}^{q}\lambda_i/\sum_{i=1}^{d}\lambda_i$. Eigenvalues $(9,1)$ give 90\% for $q=1$, yet an unseen displacement can still lie arbitrarily far along the discarded direction.
\paragraph{Why ECDR uses this concept.}
PCA proposes a geometry worth testing: score budgets along $V$ and a residual budget outside its span. Actual query membership and a sound whole-domain certificate perform the subsequent correctness work.
\paragraph{Conditions and connections.}
If PCA fits centered deltas around $\bar\Delta$, the corresponding query-domain center is generally $x_r+\bar\Delta$, not automatically the anchor $x_r$. Training-only fitting prevents leakage; it does not prove future coverage. A quantile budget remains empirical. With zero total variance the displayed fraction is undefined. \hyperref[concept:C13]{C13} supplies the support formula for a specified PCA-shaped domain.
\par\label{endconcept:C12}
\typeout{ECDR-CONCEPT-END C12 PAGE=\thepage}

\concept{C13}{Support of an orthogonal residual domain}{\hyperref[concept:C07]{C07}--\hyperref[concept:C08]{C08} and \hyperref[concept:C11]{C11}; score boxes and Euclidean orthogonality.}
\paragraph{Definition and foundation.}
An orthogonal residual domain combines bounded retained scores with a Euclidean residual restricted to their orthogonal complement:
\[
 S=\{c_0+V\diag(b)u+v:\norm{u}_\infty\le1,
          V^\top v=0,\ \norm{v}_2\le\eta\}.
\]
Here $V\in\R^{d\times q}$ has orthonormal columns, $c_0\in\R^d$ is a center, $b\in\R^q$ has nonnegative score budgets, and $\eta\ge0$ is a residual radius. Columns $u\in\R^q$ and $v\in\R^d$ are constrained latent and residual variables.
\paragraph{Symbols and dimensions.}
Let $a^\top x+\beta$ be an affine margin lower expression, with $a\in\R^d$, $\beta\in\R$, and $P_\perp=I_d-VV^\top$. Budgets $b$ and $\eta$ use the displacement coordinates' units. The diagonal matrix $\diag(b)$ is $q\times q$.
\paragraph{Derivation and example.}
The score contribution minimizes to $-\norm{\diag(b)V^\top a}_1$. Since $v$ is orthogonal to the retained span, $a^\top v=(P_\perp a)^\top v$. Cauchy--Schwarz gives its exact minimum as $-\eta\norm{P_\perp a}_2$, attained opposite $P_\perp a$ when that vector is nonzero. Thus
\[
 \ell=a^\top c_0+\beta-\norm{\diag(b)V^\top a}_1
                  -\eta\norm{P_\perp a}_2.
\]
For $V=(1,1)^\top/\sqrt2$, $a=(1,-1)^\top$, base value $a^\top c_0+\beta=1$, and $\eta=0.1$, the score term vanishes and $\ell=1-0.1\sqrt2\approx0.858579$.
\paragraph{Why ECDR uses this concept.}
It penalizes the margin only for residual directions that can actually occur in the defined domain, instead of charging the entire coefficient norm.
\paragraph{Conditions and connections.}
The residual must satisfy both orthogonality and the Euclidean bound. A coordinate residual box is a different set and has a different support. This exact geometry calculation still relies on the preceding affine neural bound being sound. \hyperref[concept:C14]{C14} handles coordinate changes. The construction does not assert that a learned direction has physical causal meaning. Its immediate mathematical benefit is an explicitly defined feasible set whose linear worst case can be derived and checked.
\par\label{endconcept:C13}
\typeout{ECDR-CONCEPT-END C13 PAGE=\thepage}

\concept{C14}{Normalization changes generators and orthogonality}{\hyperref[concept:B07]{B07}--\hyperref[concept:B08]{B08}; \hyperref[concept:C11]{C11}--\hyperref[concept:C13]{C13}; invertible diagonal coordinate transformations.}
\paragraph{Definition and foundation.}
Normalization is a change of coordinates. A raw-coordinate domain can be transformed exactly, but its old geometric interpretation must also be transported. In particular, independent channel scaling does not generally preserve Euclidean angles, orthonormality, or a spherical residual constraint.
\paragraph{Symbols and dimensions.}
Let $x\in\R^d$ be raw input, $z=D^{-1}(x-\mu)$ normalized input, $\mu\in\R^d$ a fixed mean, and $D=\diag(\sigma_1,\ldots,\sigma_d)$ with every $\sigma_i>0$. For raw domain $x=c_0+Vu+v$, $c_0,v\in\R^d$, $V\in\R^{d\times q}$, and $u\in\R^q$, the transformed generator is $B=D^{-1}V$.
\paragraph{Derivation and example.}
Substitution gives $z=D^{-1}(c_0-\mu)+Bu+D^{-1}v$. Even when $V^\top V=I_q$,
\[
 B^\top B=V^\top D^{-2}V
\]
need not equal $I_q$. For $V=(1,1)^\top/\sqrt2$ and $D=\diag(1,2)$, $\norm{B}_2^2=5/8$, so $BB^\top$ is not idempotent. Consequently $I_d-BB^\top$ is not the required complementary orthogonal projector.
A raw orthogonality constraint $V^\top v=0$ becomes $V^\top Dv_z=0$ for $v_z=D^{-1}v$; its raw Euclidean budget becomes $\norm{Dv_z}_2\le\eta$.
\paragraph{Why ECDR uses this concept.}
A certificate and membership checker must describe the same coordinates. Silently placing scaled directions in an old PCA projector can accept inputs outside the intended domain.
\paragraph{Conditions and connections.}
One safe option keeps the raw domain and includes normalization inside the certified network. Another transports every constraint and recomputes support and membership. Reorthonormalizing $B$ alone does not preserve the original coefficient box: its coefficient constraints must also be transformed. Positive scales and fixed fitted statistics are necessary for this invertible transformation. If a scale were zero, the inverse map would not exist; if statistics changed between certificate construction and use, runtime would be evaluating a different coordinate map.
\par\label{endconcept:C14}
\typeout{ECDR-CONCEPT-END C14 PAGE=\thepage}

\concept{C15}{Containment and the direction of certificate transfer}{Set inclusion, universal statements, and infima; \hyperref[concept:C18]{C18}'s margin criterion.}
\paragraph{Definition and foundation.}
A certificate valid on a larger set also applies to every subset. This follows directly from a universal quantifier: if a statement is true for every member of $\widehat S$, then it is true for every member of $S\subseteq\widehat S$. Transferring a bound in the opposite direction has no such justification.
\paragraph{Symbols and dimensions.}
Let $S,\widehat S\subseteq\R^d$ be nonempty domains, with $S\subseteq\widehat S$. For target class $c$ and competitor $j\ne c$, the scalar margin is $g_{cj}:\R^d\to\R$. A sound lower bound $\ell_j$ on $\widehat S$ means $\ell_j\le\inf_{x\in\widehat S}g_{cj}(x)$.
\paragraph{Derivation and example.}
Enlarging a minimization set cannot increase the minimum or infimum:
\[
 \inf_{x\in S}g_{cj}(x)\ge
 \inf_{x\in\widehat S}g_{cj}(x)\ge\ell_j.
\]
If $\ell_j>0$ for every competitor, the target wins on both sets. For $g(x)=1-x$, the smaller interval $[-0.5,0.5]$ has minimum $0.5$, but the larger interval $[-2,2]$ has minimum $-1$. Certification of the smaller interval therefore says nothing sufficient about the larger one.
\paragraph{Why ECDR uses this concept.}
A difficult structured domain can be enclosed by a simpler box or ball. Proving positivity over that enclosure is sufficient for the original domain, provided containment has actually been established.
\paragraph{Conditions and connections.}
A verifier's numerical output need not improve monotonically when the domain shrinks because its optimization procedure can behave differently. The mathematical property and any already-valid bound do transfer. Lower-bound failure on an enclosure does not prove failure on the enclosed set. Changing a center, radius, normalization, or generator may destroy containment and must be checked explicitly before reusing a certificate.
\par\label{endconcept:C15}
\typeout{ECDR-CONCEPT-END C15 PAGE=\thepage}

\concept{C16}{Selectors and unions of certified domains}{\hyperref[concept:C15]{C15} and \hyperref[concept:C18]{C18}; disjunction, witnesses, and finite unions.}
\paragraph{Definition and foundation.}
Selection decides which certificate to try; membership determines whether that certificate applies. Their responsibilities can be separated without weakening soundness. A selector may be heuristic, but an accepted query must have a witness: a specific valid certificate whose domain contains it.
\paragraph{Symbols and dimensions.}
Let $\mathcal E$ be a finite set of record identifiers. Record $e\in\mathcal E$ contains domain $S_e\subseteq\R^d$, target $c_e\in\{1,\ldots,K\}$, reference identity, and sound positive lower margins. A selector $s(x)$ proposes an identifier for query $x\in\R^d$. The conservative checker $M_e(x)\in\{0,1\}$ must satisfy $M_e(x)=1\Rightarrow x\in S_e$.
\paragraph{Derivation and example.}
If all chosen records share target $c$, then
\[
 \bigl(\exists e\in\mathcal E:M_e(x)=1\bigr)
 \Longrightarrow x\in\bigcup_{e\in\mathcal E}S_e
 \Longrightarrow h(x)=c,
\]
where $h$ is the certified reference decision. The first implication uses checker soundness; the second uses the certificate of whichever domain contains $x$. A selector mistake can cause a miss and extra cost. It cannot cause a false accept if the selected certificate and membership check remain sound.
\paragraph{Why ECDR uses this concept.}
Temporal event names, nearest-anchor rules, and data-driven routing can propose useful domains without themselves proving the class. The proof travels with each record, rather than with the selector's confidence.
\paragraph{Conditions and connections.}
Domains with different targets can also coexist: return the target of a valid witnessed record. For a single deterministic reference, sound certified domains of distinct strict winners cannot overlap. If two such records both accept, some assumption or implementation is inconsistent. Do not merge centers or budgets while keeping old bounds; a union is represented by a disjunction of memberships, not automatically by its enclosing box.
\par\label{endconcept:C16}
\typeout{ECDR-CONCEPT-END C16 PAGE=\thepage}

\concept{C17}{A wave domain is an embedded full-window domain}{\hyperref[concept:B09]{B09}--\hyperref[concept:B10]{B10} and \hyperref[concept:B24]{B24}--\hyperref[concept:B26]{B26}; \hyperref[concept:C15]{C15}--\hyperref[concept:C16]{C16}.}
\paragraph{Definition and foundation.}
A wave-domain certificate fixes the retained part of an anchor window and varies only the newly appended samples. It therefore describes a particular subset of full windows. Checking the new tail alone is valid only when the retained part is guaranteed to be the one used in that subset's definition.
\paragraph{Symbols and dimensions.}
Let $X_r\in\R^{C\times W}$ be a fixed anchor with $C$ channels and $W$ samples, and let integer shift $s$ satisfy $1\le s\le W$. A new wave is $U\in\R^{C\times s}$ with admissible set $\mathcal U$. Define $\mathcal T_s(X_r,U)$ by dropping the first $s$ columns of $X_r$ and appending $U$. Vectorization maps the window to $\R^{CW}$.
\paragraph{Derivation and example.}
The full certified domain corresponding to a wave set is
\[
 S_r=\{\operatorname{vec}(\mathcal T_s(X_r,U)):U\in\mathcal U\}.
\]
This is an affine embedding of $Cs$ variable coordinates into $CW$ coordinates. For $C=1,W=4,s=2$, anchor $(1,2,3,4)$ and wave box $[5,6]^2$ produce exactly windows $(3,4,u_1,u_2)$ with $u_1,u_2\in[5,6]$. Window $(30,40,5,6)$ has an admissible tail but is outside $S_r$.
\paragraph{Why ECDR uses this concept.}
The geometry explains why a checker may inspect only $Cs$ new values when trusted stream state already guarantees the overlap. It also identifies precisely what can invalidate that shortcut.
\paragraph{Conditions and connections.}
The domain record must bind the anchor, overlap, layout, normalization, and stream position. A full-window perturbation box varies different coordinates and cannot be silently substituted. After one shift, the overlap changes; \hyperref[concept:C23]{C23}--\hyperref[concept:C25]{C25} address lifecycle rules rather than assuming the wave certificate follows the predicted class. The reduced number of free variables is useful only while that affine embedding remains fixed.
\par\label{endconcept:C17}
\typeout{ECDR-CONCEPT-END C17 PAGE=\thepage}

\concept{C18}{The whole-domain strict-positive margin theorem}{\hyperref[concept:B15]{B15}--\hyperref[concept:B16]{B16}; universal quantification and sound lower bounds.}
\paragraph{Definition and foundation.}
The central decision theorem needs every competitor and every input in the certified domain. Let $F:\R^d\to\R^K$ be fixed, and let $h(x)$ choose an argmax logit under a fixed deterministic tie rule. For target $c$, each margin is $g_{cj}(x)=F_c(x)-F_j(x)$.
\paragraph{Symbols and dimensions.}
The domain $S\subseteq\R^d$ is nonempty, $K\ge2$, and $j\in\{1,\ldots,K\}\setminus\{c\}$. A scalar $\ell_j$ is sound when $\ell_j\le g_{cj}(x)$ for every $x\in S$. Logits, margins, and lower bounds are dimensionless scores unless a different scale is declared.
\paragraph{Derivation and proof.}
Assume $\ell_j>0$ for every competitor. Pick an arbitrary $x\in S$. For each $j\ne c$, soundness gives $F_c(x)-F_j(x)\ge\ell_j>0$, so $F_c(x)>F_j(x)$. Therefore $c$ is the unique maximizer and $h(x)=c$. Since the chosen $x$ was arbitrary,
\[
 \left[\forall j\ne c:\ \ell_j\le\inf_{x\in S}g_{cj}(x),\ \ell_j>0\right]
 \Longrightarrow\left[\forall x\in S:\ h(x)=c\right].
\]
No distributional assumption is used in this implication.
\paragraph{Why ECDR uses this concept.}
The theorem separates offline proof from online membership. Runtime need not evaluate the CNN again once it establishes that the current input belongs to this unchanged, valid domain.
\paragraph{Conditions and connections.}
The implication is universal but conditional on valid premises. It does not claim that every domain is certifiable. A zero bound fails the strict gate, but may merely reflect looseness or a tie handled in favor of $c$. Even a true positive margin at every point need not have a positive uniform infimum on a noncompact set. Numerical code, model identity, preprocessing, and membership add \hyperref[concept:C20]{C20}--\hyperref[concept:C23]{C23} obligations.
\par\label{endconcept:C18}
\typeout{ECDR-CONCEPT-END C18 PAGE=\thepage}

\concept{C19}{Sign-aware certification from new and dropped feature sums}{\hyperref[concept:B24]{B24}--\hyperref[concept:B26]{B26}; \hyperref[concept:C03]{C03} and \hyperref[concept:C18]{C18}.}
\paragraph{Definition and foundation.}
For a valid, stride-one, local CNN with ReLU and global average pooling, aligned overlapping feature positions cancel in the difference between consecutive windows. Only newly arriving and dropped final-layer features contribute to the pooled-feature change. The resulting margin identity permits a specialized certificate.
\paragraph{Symbols and dimensions.}
Let $Z_r\in\R^{C\times W}$ be the old normalized window, $U\in\R^{C\times s}$ a wave, and $Z^\prime=\mathcal T_s(Z_r,U)$. Integers $C,W$ count channels/samples, $d_h$ counts final channels, $n\ge1$ counts final positions, and $1\le s\le n$ is the shift. The head is $A\in\R^{K\times d_h}$ for $K\ge2$ classes; $c$ is target and $j\ne c$ is competitor. Define $q_j=(A_{c,:}-A_{j,:})^\top$, $\gamma_{cj}=f_{\theta,c}-f_{\theta,j}$ and $m_j=\gamma_{cj}(Z_r)$. Dropped/new sums $D_{\rm drop},D_{\rm new}(U)\in\R^{d_h}$ are as in \hyperref[concept:B25]{B25}.
\paragraph{Derivation and example.}
The cancellation identity gives $\gamma_{cj}(Z^\prime)=m_j+q_j^\top(D_{\rm new}(U)-D_{\rm drop})/n$. Suppose $l\le D_{\rm new}(U)\le u$ for every admissible wave, with $l,u\in\R^{d_h}$. Writing $q_j^+=\max(q_j,0)$ and $q_j^-=\min(q_j,0)$ coordinatewise gives
\[
 \ell_j=m_j+\frac{(q_j^+)^\top l+(q_j^-)^\top u
                   -q_j^\top D_{\rm drop}}{n}.
\]
For $n=10$, $m_j=0.4$, $q_j=(2,-1)^\top$, $D_{\rm drop}=(1,2)^\top$, $l=(0,1)^\top$, and $u=(2,3)^\top$, the bound is $0.4-3/10=0.1$.
\paragraph{Why ECDR uses this concept.}
It uses known anchor features and localized uncertainty rather than rebuilding an unrelated full-window bound.
\paragraph{Conditions and connections.}
The feature-shift identity must actually hold; padding, pooling alignment, or stateful operators require separate treatment. New-sum bounds must cover the whole wave domain. Dropped sums must be known with justified numerical semantics. Every competitor must pass \hyperref[concept:C18]{C18}, and a reused class alone does not supply new exact feature sums. The numerical example covers one competitor only. In a three-class problem, the second competitor needs its own coefficient column and positive bound before any query can be accepted as class-invariant.
\par\label{endconcept:C19}
\typeout{ECDR-CONCEPT-END C19 PAGE=\thepage}

\concept{C20}{Numerical soundness requires two separate error budgets}{\hyperref[concept:C18]{C18}--\hyperref[concept:C19]{C19}; real arithmetic versus executable arithmetic.}
\paragraph{Definition and foundation.}
A real-number theorem describes an ideal mathematical function. A deployed program evaluates rounded operations, while a verifier also computes rounded bounds. These are three distinct objects: the ideal margin, the executable margin, and the computed certificate. A positive computed number proves executable positivity only after both discrepancies are bounded.
\paragraph{Symbols and dimensions.}
For one competitor, let $g:S\to\R$ be the ideal margin, $\widetilde g:S\to\R$ the actual executable-reference margin, and $\widehat\ell\in\R$ the verifier's computed lower value. Suppose proved budgets $\delta\ge0$ and $\eta\ge0$ satisfy
\[
 \widehat\ell-\delta\le\inf_{x\in S}g(x),\qquad
 \abs{\widetilde g(x)-g(x)}\le\eta\quad\text{for every }x\in S.
\]
All four scalar quantities have margin units. The budgets may differ by competitor and domain.
\paragraph{Derivation and example.}
For any accepted query whose exact represented input belongs to $S$,
\[
 \widetilde g(x)\ge g(x)-\eta
 \ge\widehat\ell-\delta-\eta.
\]
Thus $\widehat\ell>\delta+\eta$ is a sufficient strict gate. With $\widehat\ell=0.020$, $\delta=0.003$, and $\eta=0.004$, the executable lower margin is 0.013. With $\widehat\ell=0.005$, the lower value becomes $-0.002$ and this gate cannot certify.
\paragraph{Why ECDR uses this concept.}
The operational objective is agreement with the frozen reference program, not merely agreement with an idealized formula that may round differently near a decision boundary.
\paragraph{Conditions and connections.}
These example budgets are assumed valid for illustration; they are not arbitrary tolerances. Their derivation must cover the actual operations, weights, precisions, reduction order, preprocessing, and supported inputs. NaNs or unsupported exceptional values require explicit rejection or treatment. Zero mismatches in sampled tests cannot replace a universal error bound. Membership arithmetic needs its own enclosure, addressed in \hyperref[concept:C21]{C21}.
\par\label{endconcept:C20}
\typeout{ECDR-CONCEPT-END C20 PAGE=\thepage}

\concept{C21}{Conservative numerical membership}{\hyperref[concept:C09]{C09} and \hyperref[concept:C20]{C20}; interval containment and rounding error.}
\paragraph{Definition and foundation.}
A conservative membership checker may reject an admissible query when uncertain, but must never accept an inadmissible one. In numerical code, the exact mathematical coordinate represented after preprocessing can differ from a rounded point estimate. The checker therefore needs a justified enclosure of that coordinate, rather than an informal tolerance added to the certified set.
\paragraph{Symbols and dimensions.}
Let the certified box be $S=\{z\in\R^d:l_i\le z_i\le u_i\}$, where $l,u\in\R^d$ and $l\le u$. For the particular query, let $[\underline z_i,\overline z_i]$ be a sound interval for its exact normalized coordinate $z_i$. These endpoints must include preprocessing and comparison-related numerical uncertainty. All coordinates are in the normalization used when building the certificate.
\paragraph{Derivation and example.}
Accept the box test only if
\[
 l_i\le\underline z_i\le\overline z_i\le u_i
 \qquad\text{for every }i=1,\ldots,d.
\]
Then every possible exact query consistent with the enclosures lies in $S$, so the actual query does too. For a one-dimensional domain $[0,0.1]$, interval $[0.099,0.101]$ cannot be accepted even if the rounded estimate is 0.1. Interval $[0.098,0.099]$ lies entirely inside and can pass.
\paragraph{Why ECDR uses this concept.}
Membership is the logical bridge from a whole-domain theorem to the concrete runtime query. An unsound bridge invalidates even a perfectly computed margin bound.
\paragraph{Conditions and connections.}
The domain endpoints themselves must have fixed, justified meaning. Directed rounding or interval arithmetic can implement enclosures, but their correctness depends on the actual operation sequence. For PCA membership, score errors, residual errors, and orthogonality errors must also be enclosed. Enlarging a comparison threshold without certifying the enlarged domain is unsafe. Rejecting uncertainty and falling back can preserve soundness, at the cost of lower coverage.
\par\label{endconcept:C21}
\typeout{ECDR-CONCEPT-END C21 PAGE=\thepage}

\concept{C22}{Membership algorithms and their computational obligations}{\hyperref[concept:C09]{C09}--\hyperref[concept:C14]{C14} and \hyperref[concept:C21]{C21}; arithmetic complexity and existential constraints.}
\paragraph{Definition and foundation.}
Domain membership is an algorithmic task with a geometry-specific proof obligation. An explicit box requires comparisons; a projection domain requires linear algebra; a general generator domain asks whether a feasible latent representation exists. Giving these all the same complexity hides materially different work.
\paragraph{Symbols and dimensions.}
Let $x,c_0\in\R^d$ and $w=x-c_0$. A PCA matrix $V\in\R^{d\times q}$ has orthonormal columns, score budgets $b\in\R^q_{\ge0}$, and residual radius $\eta\ge0$. For a general generator domain use $B\in\R^{d\times q}$, latent $u\in\R^q$, and residual $v\in\R^d$. Here complexity counts scalar arithmetic or coordinate accesses, not measured elapsed time.
\paragraph{Derivation and example.}
A diagonal box uses $d$ coordinate comparisons in a full scan. The PCA domain of \hyperref[concept:C13]{C13} has equivalent tests
\[
 \abs{V^\top w}\le b,\qquad
 \norm{w-VV^\top w}_2\le\eta.
\]
Dense score and reconstruction operations cost $O(dq+d)$, with the $d$ term accounting for residual processing. This page assumes $1\le q\le d$. If some $b_i=0$, the corresponding exact score must be zero. Under exact orthonormality, the residual squared norm also equals $\norm{w}_2^2-\norm{V^\top w}_2^2$, but numerical cancellation requires care.
\paragraph{Why ECDR uses this concept.}
The checker can dominate a small CNN's cost. A domain with better geometric coverage can still produce worse total latency if membership is expensive.
\paragraph{Conditions and connections.}
For $x=c_0+Bu+v$ with box constraints, membership is feasibility of a linear system of equalities and inequalities. A verified feasible witness suffices to accept; a failed heuristic witness does not prove nonmembership. General residual norms can require other constrained solvers. Replacing feasibility by a projection shortcut needs an equivalence proof. Wave-only checks reduce the dimension to $Cs$ only with \hyperref[concept:C17]{C17}'s trusted overlap identity.
\par\label{endconcept:C22}
\typeout{ECDR-CONCEPT-END C22 PAGE=\thepage}

\concept{C23}{The fixed-anchor reuse theorem}{\hyperref[concept:C18]{C18} and \hyperref[concept:C20]{C20}--\hyperref[concept:C21]{C21}; frozen reference identity and two-branch reasoning.}
\paragraph{Definition and foundation.}
Fixed-anchor reuse keeps the meaning of a certificate unchanged while testing successive queries against it. A reused output is a decision result; it does not automatically become a newly evaluated anchor. The proof can apply independently to each accepted query, without assuming successive queries are close.
\paragraph{Symbols and dimensions.}
Let $R$ denote the frozen reference identity: graph, weights, evaluation state, preprocessing, class order, precision semantics, and tie rule. Its deterministic decision is $h_R:\mathcal X\to\{1,\ldots,K\}$ on valid input space $\mathcal X\subseteq\R^d$. A record $(x_r,S_r,c_r,R)$ stores an anchor $x_r$, a certified domain $S_r$, and target $c_r$. Queries are $x_t\in\mathcal X$, indexed by integer $t\ge1$.
\paragraph{Derivation and proof.}
Assume the record has sound strictly positive margins over $S_r$, including any necessary numerical justification. Assume a conservative checker $M_r$ satisfies $M_r(x)=1\Rightarrow x\in S_r$. The policy is
\[
 \widehat h(x_t)=
 \begin{cases}c_r,&\text{if identity matches and }M_r(x_t)=1,\\
 h_R(x_t),&\text{otherwise.}\end{cases}
\]
On acceptance, \hyperref[concept:C18]{C18} gives $c_r=h_R(x_t)$. On rejection or uncertainty, the fallback returns $h_R(x_t)$ directly. These exhaustive branches establish equality for every valid query under the assumptions.
\paragraph{Why ECDR uses this concept.}
It states precisely what certified decision reuse promises: agreement with the chosen reference while avoiding some reference evaluations.
\paragraph{Conditions and connections.}
The anchor, domain, and reference identity must remain the same. An updated model or recentered domain requires a new certificate or a proved transfer. Reused predictions do not provide uncomputed feature tensors. The theorem guarantees neither ground-truth accuracy, useful hit rate, nor lower cost. \hyperref[concept:C24]{C24} shows why moving the anchor without transferring the proof is unsafe.
\par\label{endconcept:C23}
\typeout{ECDR-CONCEPT-END C23 PAGE=\thepage}

\concept{C24}{A stale certificate and the danger of automatic recentering}{\hyperref[concept:C17]{C17} and \hyperref[concept:C23]{C23}; \hyperref[concept:B27]{B27}--\hyperref[concept:B28]{B28}'s computed toy CNN.}
\paragraph{Definition and foundation.}
A certificate is stale when runtime applies it to inputs outside the domain for which its bounds were proved. Class agreement at one accepted point is insufficient to transfer a neighborhood or a wave-domain certificate to that point. This is a logical failure even in exact arithmetic.
\paragraph{Symbols and dimensions.}
For a scalar counterexample let $g(x)=1-x$, where positive margin selects class 1. The original anchor is $x_r=0$ and the certified set is $S_r=[-0.6,0.6]$. Let $x_1=0.6$ be an accepted query and $x_2=1.2$ a later query. All quantities in this counterexample are dimensionless.
\paragraph{Derivation and example.}
The original minimum is $\inf_{x\in S_r}g(x)=0.4>0$. Although $x_1$ is certified, replacing the center by $x_1$ while keeping radius 0.6 admits $x_2$. Yet $g(x_2)=-0.2$, so class 1 loses. The triangle inequality gives $\abs{x_2-x_r}\le\abs{x_2-x_1}+\abs{x_1-x_r}=1.2$, not 0.6.
The \hyperref[concept:B27]{B27}--\hyperref[concept:B28]{B28} toy CNN gives the same lifecycle failure with $W=6$ and shift $s=2$. Repeating admissible wave $(7,8)$ produces second-step window $(5,6,7,8,7,8)$. Its pooled features are $(5,1/2)$ and logits $(17/2,-4,9)$, so class 3 beats the formerly certified class 1.
\paragraph{Why ECDR uses this concept.}
A tail checker can appear valid while silently ignoring a changed overlap. The counterexample identifies what record identity must protect.
\paragraph{Conditions and connections.}
In the CNN example, $W$ counts input samples, $s$ counts dropped/appended samples, and the three logits correspond to the toy's fixed class order. The first-step proof remains correct: the second overlap is a different object. A fresh reference evaluation with a fresh certificate, a proved transfer rule, or \hyperref[concept:C25]{C25}'s reachable-domain construction can justify continuation.
\par\label{endconcept:C24}
\typeout{ECDR-CONCEPT-END C24 PAGE=\thepage}

\concept{C25}{Finite-horizon reachable sets and inductive reuse}{\hyperref[concept:C17]{C17}--\hyperref[concept:C18]{C18} and \hyperref[concept:C24]{C24}; mathematical induction and set-valued transitions.}
\paragraph{Definition and foundation.}
A horizon certificate covers every input reachable through an entire admissible sequence of waves. Its guarantee is stronger than certifying each wave against the original overlap independently. The state uncertainty from all earlier waves must be preserved in each reachable set.
\paragraph{Symbols and dimensions.}
Let integer $C\ge1$ count channels and $W\ge1$ count samples; $X_r\in\R^{C\times W}$ is the evaluated starting window, $1\le s\le W$ the shift, and $Q\ge1$ the finite number of steps. For each step $q=1,\ldots,Q$, let $\mathcal U_q\subseteq\R^{C\times s}$ be a predefined wave domain. Define the shift-and-append map $\mathcal T_s$ as in \hyperref[concept:B09]{B09} and the sets
\[
 \mathcal R_0=\{X_r\},\qquad
 \mathcal R_q=\{\mathcal T_s(X,U):X\in\mathcal R_{q-1},\ U\in\mathcal U_q\}.
\]
The fixed target class is $c$; vectorization is implicit when evaluating the reference.
\paragraph{Derivation and proof.}
Assume sound strictly positive target margins on every $\mathcal R_q$. The actual initial state belongs to $\mathcal R_0$. If the actual state belongs to $\mathcal R_{q-1}$ and the next wave belongs to $\mathcal U_q$, the transition definition puts the next state in $\mathcal R_q$. Induction gives membership at every step, and \hyperref[concept:C18]{C18} gives class $c$ throughout.
\paragraph{Why ECDR uses this concept.}
Runtime may check the current wave and horizon counter when trusted state guarantees all earlier conditions and the certificate covers their complete possible history.
\paragraph{Conditions and connections.}
A failed wave check breaks this induction path; runtime must fall back or use another justified record. For $W=4,s=2$, two steps remove all original samples, yet a valid reachable-set proof can still work. In general only $\max(W-qs,0)$ original sample positions remain per channel. Enlarging reachable sets can loosen bounds and increase build cost. Soundness alone does not establish that a long useful horizon exists.
\par\label{endconcept:C25}
\typeout{ECDR-CONCEPT-END C25 PAGE=\thepage}

\concept{C26}{Trace accounting and amortized inference cost}{Finite sums, means, and conditional expectation; \hyperref[concept:C23]{C23}--\hyperref[concept:C25]{C25}.}
\paragraph{Definition and foundation.}
Amortized cost distributes every operation over the queries that incur it. The first task is complete trace accounting, before replacing variable costs by constants or averages. A failed certificate attempt still consumes time and must be charged even though it produces no reuse.
\paragraph{Symbols and dimensions.}
Let $N\ge1$ be the query count, $\mathcal M\subseteq\{1,\ldots,N\}$ the miss indices, and $B\ge0$ the number of all build attempts. Let $f_t$ be baseline forward time for query $t$, $f_t^{\rm miss}$ its fallback time, $v_t$ lookup/verification time, $b_k$ build time for attempt $k$, and $T_{\rm state}$ uncounted state/initialization time. Every cost uses one time unit, such as milliseconds.
\paragraph{Derivation and example.}
The totals are
\[
 T_{\rm full}=\sum_{t=1}^{N}f_t,\qquad
 T_{\rm ECDR}=\sum_{t=1}^{N}v_t+
 \sum_{t\in\mathcal M}f_t^{\rm miss}+
 \sum_{k=1}^{B}b_k+T_{\rm state}.
\]
Divide by $N$ only after counting all terms. If forward costs vary, the expected miss contribution is $\Pr(\mathrm{miss})\E[f\mid\mathrm{miss}]$, whenever that conditional expectation is defined. It is not automatically miss probability times unconditional mean forward time.
For forward costs $(1,1,9,9)$ ms, hitting the first two queries saves only 2 ms. Verification of 1 ms per query gives total 22 ms versus baseline 20 ms, despite a 50\% hit rate.
\paragraph{Why ECDR uses this concept.}
The arithmetic prevents a high reuse percentage from hiding expensive builds, misses, or membership checks.
\paragraph{Conditions and connections.}
Baseline and policy must share the same inclusion boundary for normalization, input handling, and initialization. Shadow-reference evaluations used for validation need a separately stated timing role. Avoid double-counting state work already inside verification or builds. \hyperref[concept:C27]{C27} derives a break-even threshold only after adding explicit constant-cost assumptions.
\par\label{endconcept:C26}
\typeout{ECDR-CONCEPT-END C26 PAGE=\thepage}

\concept{C27}{Break-even thresholds and useful certificate lifetimes}{\hyperref[concept:C26]{C26}; algebraic inequalities and dimensionally consistent costs.}
\paragraph{Definition and foundation.}
Break-even means total ECDR time is strictly less than the corresponding reference-only total. It is a performance condition, independent of decision soundness. A sound policy can be slower, and an apparently faster unsound policy does not satisfy the certified objective.
\paragraph{Symbols and dimensions.}
Assume constant forward cost $t_f>0$, verification cost $t_v\ge0$, additional state cost $t_m\ge0$ per query, and build cost $t_b\ge0$. Let $N\ge1,H\ge0,B\ge0$ be integer query, hit and build counts, with $H\le N$. Let $r=H/N\in[0,1]$ be hits divided by queries, and $\rho_b=B/N\ge0$ builds per query. All $t$ quantities use milliseconds; the rates are dimensionless.
\paragraph{Derivation and example.}
\hyperref[concept:C26]{C26} reduces under these assumptions to
\[
 \bar t_{\rm ECDR}=t_v+t_m+(1-r)t_f+\rho_b t_b.
\]
Subtracting this from baseline $t_f$ gives the saving condition $rt_f>t_v+t_m+\rho_b t_b$. If one build occurs on every miss, so $\rho_b=1-r$, then
\[
 r>\frac{t_v+t_m+t_b}{t_f+t_b}.
\]
With $t_f=4$ ms, $t_v+t_m=0.08$ ms, and $t_b=1.5$ ms, the threshold is $1.58/5.5\approx28.7273\%$. For 100 queries, 80 hits, and 20 builds, the total is $8+80+30=118$ ms versus 400 ms, an illustrative ratio of approximately 3.38983.
\paragraph{Why ECDR uses this concept.}
It identifies the overhead budget that a proposed domain and lifecycle must meet, rather than using coverage alone as evidence of efficiency.
\paragraph{Conditions and connections.}
For one build followed by $L$ all-hit queries, benefit requires $L>t_b/(t_f-t_v-t_m)$ when the denominator is positive. With $t_b=8$ ms and the same other costs, at least 3 such hits are needed. These are assumed costs, not HAPT measurements. Variable or path-dependent costs require \hyperref[concept:C26]{C26}'s trace equation instead. Likewise, a build scheduled periodically rather than on every miss must retain its actual $\rho_b$; substituting $1-r$ would charge the wrong number of attempts. The useful lifetime is the number of later verified hits, not merely elapsed wall time.
\par\label{endconcept:C27}
\typeout{ECDR-CONCEPT-END C27 PAGE=\thepage}

\concept{C28}{Soundness, agreement, accuracy, coverage, and net benefit}{\hyperref[concept:B30]{B30}; \hyperref[concept:C18]{C18}, \hyperref[concept:C23]{C23}, and \hyperref[concept:C26]{C26}--\hyperref[concept:C27]{C27}; explicit metric denominators.}
\paragraph{Definition and foundation.}
These quantities answer different questions. Soundness is a universal implication under declared assumptions. Agreement measures equality with the reference on evaluated cases. Accuracy measures equality with a ground-truth label. Coverage measures how often a domain or gate accepts. Net benefit measures saved cost after overhead. None is a substitute for the others.
\paragraph{Symbols and dimensions.}
Let $N$ be all queries, $N_b$ build attempts, $N_b^+$ successful certified builds, $N_m$ membership trials, $N_m^+$ membership successes, $N_r$ actual reused queries, and $N_a$ reused queries agreeing with shadow reference. These are nonnegative integer counts. Let $y_t$ be a known ground-truth class and $\widehat h_t$ the policy's prediction on query $t$.
\paragraph{Derivation and example.}
The operational rates are
\[
 p_b=\frac{N_b^+}{N_b},\quad p_m=\frac{N_m^+}{N_m},\quad
 r=\frac{N_r}{N},\quad A_r=\frac{N_a}{N_r}.
\]
A zero denominator makes the associated rate undefined, reported as N/A. Accuracy on a declared labeled index set $\mathcal L$ is $\abs{\mathcal L}^{-1}\sum_{t\in\mathcal L}\mathbf1\{\widehat h_t=y_t\}$ when $\abs{\mathcal L}>0$.
If, on the same 100 queries, 60 satisfy membership and 50 have a valid relevant certificate, their intersection can contain 10--50 queries. The product $0.6(0.5)=0.3$ requires an independence assumption, not just the two marginal rates.
\paragraph{Why ECDR uses this concept.}
It prevents reporting geometric hits as certified reuse, empirical zero mismatch as a universal theorem, or reference agreement as ground-truth correctness.
\paragraph{Conditions and connections.}
Build counts and query counts often have different denominators, making even the intersection example inapplicable without alignment. Adjacent-pair coverage also differs from fixed-anchor lifecycle coverage. A singleton domain can be sound but useless; a broad certified domain can be too costly. An ECDR result therefore needs separate evidence for correctness assumptions, useful reuse, and net time or energy saving.
\par\label{endconcept:C28}
\typeout{ECDR-CONCEPT-END C28 PAGE=\thepage}

\clearpage
\renewcommand{\conceptid}{Coverage crosswalks}
\appendix
\section{Coverage of the previous mathematical explanation}
The previous report grouped its mathematics as M01--M29. These groups are
retained as a coverage inventory, not as new concepts or claims of novelty.
The pages below identify where each group is explained at greater depth.
\small
\begin{longtable}{@{}P{1.2cm}P{7.3cm}P{6.7cm}@{}}\toprule
Group & Mathematical content & Concepts\\\midrule\endhead
M01 & Reference composition, logits and deterministic argmax & F04, F17, B14--B16\\
M02 & Universal implication, regions and sufficient certificates & F01--F03, C18, C23\\
M03 & Sensor observations, tensors and vectorization & B01--B03\\
M04 & Sampling clock, query stride and tensor payload & B01, B04, B23\\
M05 & Window counts and annotated-span grid counts & F08--F09, B02, B04--B05\\
M06 & Label counts and majority thresholds & B06\\
M07 & Normalization, transported budgets and weighted norms & B07--B08, C14\\
M08 & Split boundaries, temporal dependence and uncertainty & F30--F34, B13\\
M09 & Shift--append maps and aligned overlap & F15, B09--B10\\
M10 & Counterexample to overlap implying a small difference & B11\\
M11 & Wave-region dimension and fixed overlap & B09, C17\\
M12 & Marginal/joint coverage and the union bound & F29--F34, B12, B29\\
M13 & CNN parameters, MAC/FLOP, receptive field and jump & B19--B23; expanded CNN companion\\
M14 & Accuracy, precision, recall, F1 and dispersion & F31--F32, B30\\
M15 & Pairwise margins, infima and strict sufficient gates & F25, B16, C18\\
M16 & Continuity, compactness and local invariance & F26--F28, B17\\
M17 & Universal-radius, shift and stale-state counterexamples & B11, B18, C24\\
M18 & Lipschitz constants, operator norms and certified balls & F21--F23, C01--C02\\
M19 & Interval signs, activation propagation and induction & C03--C04\\
M20 & Affine ReLU relaxation and sign-aware substitution & C05--C06\\
M21 & Support functions, duality, boxes and generators & F22, C07--C10\\
M22 & Linear span, PCA, residuals and normalization & F12--F19, C11--C14\\
M23 & Containment, selection and the certified union & C15--C16\\
M24 & Feature shift, GAP and margin-delta identities & B24--B28, C19\\
M25 & Membership, witnesses and access-complexity bounds & C09, C21--C22\\
M26 & Fixed anchor, stale certificates and reachable horizons & C23--C25\\
M27 & Verifier error, reference-program error and containment & C20--C21\\
M28 & Trace accounting, amortization and break-even cost & C26--C27\\
M29 & Research metrics, denominators and hypothesis status & B30, C28, Appendix B\\\bottomrule
\end{longtable}
\normalsize

\clearpage
\section{Interpreting all twenty-four research hypotheses}
This ledger preserves the earlier H01--H24 questions while separating
mathematical implications from claims that require experiments. ``Conditional''
means the concept page states and uses the necessary hypotheses; it is not
a report that an implementation has established those hypotheses.
\small
\begin{longtable}{@{}P{0.9cm}P{7.2cm}P{7.1cm}@{}}\toprule
ID & Question or claim & Logical status and concept connection\\\midrule\endhead
H01 & Identical raw samples remain equal after a correctly aligned shift. & Identity under the window contract; B10.\\
H02 & Large aligned overlap forces small ordinary input distance. & False without additional assumptions; B11 gives an amplitude-scalable counterexample.\\
H03 & The same true activity forces the same CNN decision. & Unsupported in general; B06, B14--B16 distinguish labels and reference decisions.\\
H04 & Continuity and a strict anchor margin yield a local invariant region. & Conditional theorem; B17. Its usefulness is not implied.\\
H05 & Sound positive lower margins over a whole region preserve the target. & Conditional sufficient theorem; C18.\\
H06 & A valid Lipschitz rate supplies a certified distance budget. & Conditional bound; C01--C02, with all competitors and coordinate units specified.\\
H07 & IBP finds every genuinely invariant region. & False in general because interval independence loses dependencies; C04.\\
H08 & Direct margin reasoning can preserve cancellation. & Algebraically possible; C06 and B28 give examples, not a universal speed advantage.\\
H09 & A tighter certificate method always costs less. & Does not follow; C26--C27 require complete cost accounting.\\
H10 & Temporal variation directions have small impact on the relevant margins. & Model/data hypothesis; C10--C14 show what directional coefficients mean.\\
H11 & A structured domain is superior to a box. & Open without matched coverage, cost, bounds and output semantics; C07--C17.\\
H12 & A fixed-overlap wave region can have only $Cs$ free variables. & Conditional geometric result; B09, C17. Additional overlap uncertainty changes it.\\
H13 & Aligned valid-CNN features coincide on fully shared receptive fields. & Conditional identity; B24.\\
H14 & GAP changes by new feature sum minus dropped feature sum. & Conditional exact real-arithmetic identity; B25--B26.\\
H15 & A reused query can automatically replace the old certified anchor. & False without new or transferred evidence; C24 gives a computed counterexample.\\
H16 & A sound positive certificate remains safe under repeated fixed-anchor membership. & Conditional theorem; C23, with matching identities and conservative acceptance.\\
H17 & Positive bounds on reachable sets certify a finite horizon. & Conditional induction; C25. Long horizons need not be useful or cheap.\\
H18 & Exact diagonal-box membership requires linear worst-case coordinate access. & Conditional access-model lower bound, finite positive-width box; C09, C22.\\
H19 & A sound domain necessarily covers future queries. & False; conditional safety supplies no nonzero coverage guarantee; C18, C28.\\
H20 & Many reuse hits necessarily reduce total runtime. & False if overhead dominates; C26--C27.\\
H21 & ECDR outperforms full CNN inference or another cache on HAPT. & Empirical claim not established by these definitions or toy examples; C28.\\
H22 & An ideal real-arithmetic proof directly certifies a float32 program. & False without executable-error and membership guarantees; C20--C21.\\
H23 & An event name alone authorizes returning an old class. & False; selector output, membership and positive certificate have different duties; C16--C18.\\
H24 & There is a regime that is simultaneously sound, useful and cheaper. & Central research hypothesis requiring both proof and correctly scoped measurements; C18, C26--C28.\\\bottomrule
\end{longtable}
\normalsize

\clearpage
\renewcommand{\conceptid}{Sources and evidence}
\section{Sources, authoring provenance and scope}
The conceptual coverage comes from the existing ECDR Week1--3 theory and
the previous mathematical explanation in this workspace. The new exposition
is written in English and starts from the prerequisites needed to understand
their formulas. Standard linear-algebra background can be checked in
Axler's author-hosted open text, especially Sections1B,2A--2C,3A,6A--6C
and7F \cite{axler}. The concept pages supply their own definitions,
examples and project-specific reasoning rather than reproducing that text.

The verification methods are established prior work: Lipschitz-margin
reasoning \cite{lmt}, interval propagation \cite{ibp}, CROWN's output
relaxations \cite{crown}, and computational-graph LiRPA \cite{lirpa}.
These references identify method provenance; using them is not a novelty
claim or proof that a current implementation is compatible with every
operator. Architecture-dependent streaming reuse has separate prior art
\cite{streaminnc}. Metadata about HAPT is distinguished from the local
window policy and from a model actually trained on that data \cite{hapt}.

The theoretical CNN configuration is valid Conv1D--ReLU--GAP with an affine
head. Its worked counts and miniature network are specified explicitly in
the relevant concept pages. They are not new training, device measurements
or certified HAPT coverage results. Mathematical proofs use finite real
inputs unless a numerical-error assumption is explicitly added.

The English atlas is saved as a standalone editable source:
\code{ECDR_Mathematics_Concepts_EN.tex}. The former Vietnamese explanation
remains \code{ECDR_Mathematics_Explained.tex}. The English CNN source uses
the existing editor path \code{ECDR_CNN_Model_Analysis.tex}; its preserved
Vietnamese counterpart is \code{ECDR_CNN_Model_Analysis_VI.tex}.
The data-analysis document is outside this revision's scope.

\begin{thebibliography}{99}
\bibitem{axler} S. Axler. \emph{Linear Algebra Done Right}, fourth edition.
Author-hosted open-access edition, revised13 July2026.
\url{https://linear.axler.net/LADR4e.pdf}.
\bibitem{hapt} J. Reyes-Ortiz et al. \emph{Smartphone-Based Recognition of
Human Activities and Postural Transitions}. UCI Machine Learning Repository,
2015. \url{https://doi.org/10.24432/C54G7M}.
\bibitem{lmt} Y. Tsuzuku, I. Sato and M. Sugiyama. \emph{Lipschitz-Margin
Training: Scalable Certification of Perturbation Invariance for Deep Neural
Networks}. NeurIPS,2018. \url{https://arxiv.org/abs/1802.04034}.
\bibitem{ibp} S. Gowal et al. \emph{On the Effectiveness of Interval Bound
Propagation for Training Verifiably Robust Models}.2018.
\url{https://arxiv.org/abs/1810.12715}.
\bibitem{crown} H. Zhang et al. \emph{Efficient Neural Network Robustness
Certification with General Activation Functions}. NeurIPS,2018.
\url{https://arxiv.org/abs/1811.00866}.
\bibitem{lirpa} K. Xu et al. \emph{Automatic Perturbation Analysis for Scalable
Certified Robustness and Beyond}. NeurIPS,2020.
\url{https://arxiv.org/abs/2002.12920}.
\bibitem{streaminnc} C. Kechris et al. \emph{Don't Think It Twice: Exploit Shift
Invariance for Efficient Online Streaming Inference of CNNs}.2024.
\url{https://arxiv.org/abs/2408.03223}.
\end{thebibliography}
\end{document}
